% Hyperbolic and Inverse hyperbolic Functions — Further Mathematics Upper Sixth — Cours signé OnBuch+
% Official MINESEC syllabus. Compiler avec : tectonic FMA01-hyperbolic-and-inverse-hyperbolic-functions.tex
\newcommand{\DOCMATIERE}{Further Mathematics}
\newcommand{\DOCNIVEAU}{Upper Sixth}
\newcommand{\DOCMODULE}{Upper Sixth — Further mathematics}
\newcommand{\DOCLECON}{1}
\newcommand{\DOCTITRE}{Hyperbolic and Inverse hyperbolic Functions}
% OnBuch+ — préambule « cours signés » ENRICHI (compilé avec tectonic / XeLaTeX)
% Système visuel « Pop » d'OnBuch+ : fond crème (jamais blanc), bordures encre
% épaisses, ombres « dures » décalées, titres Archivo Black en capitales.
% Version MATHÉMATIQUES : graphes (pgfplots/tikz), boîtes pédagogiques
% (définition, propriété, méthode, attention, le savais-tu, exercice résolu,
% exercice, TP, bilan), page de garde et signature OnBuch+ ancrée en bas.
%
% Macros attendues dans main.tex AVANT \input{preamble} :
%   \DOCMATIERE  \DOCNIVEAU  \DOCMODULE  \DOCLECON (numéro)  \DOCTITRE
\documentclass[11pt]{article}

\usepackage{fontspec}
\usepackage[english]{babel}
\usepackage[a4paper,margin=2cm,top=3.4cm,bottom=2.8cm,headheight=1.4cm,footskip=1.3cm]{geometry}
\usepackage{amsmath,amssymb,amsfonts}
\usepackage{mathtools} % \xmapsto, \coloneqq... souvent utilisés par les modèles
\usepackage{cancel} % simplifications barrées (bilans de forces)
% \abs{z} (module, valeur absolue) et \norm{u} (norme), utilisés par les modèles
\providecommand{\abs}{}\let\abs\relax\DeclarePairedDelimiter{\abs}{\lvert}{\rvert}
\providecommand{\norm}{}\let\norm\relax\DeclarePairedDelimiter{\norm}{\lVert}{\rVert}
\usepackage[table]{xcolor} % [table] : fournit aussi \rowcolors (alternance de lignes)
\usepackage{pagecolor}
\usepackage{tikz}
\usetikzlibrary{arrows.meta,positioning,calc,shapes.geometric,shapes.misc,shapes.arrows,decorations.pathmorphing,decorations.pathreplacing,decorations.markings,patterns,shadows,mindmap,trees,fit,backgrounds,matrix,intersections,angles,quotes}
\usepackage{pgfplots}
\usepackage[version=4]{mhchem} % équations nucléaires \ce{^{235}_{92}U}
\usepackage[european,siunitx]{circuitikz} % schémas électriques
\pgfplotsset{compat=1.18}
\usepgfplotslibrary{fillbetween,groupplots} % aires entre courbes (calcul intégral)
\usepackage{enumitem}
\usepackage{fancyhdr}
% [most] charge toutes les bibliothèques tcolorbox (listings, minted, raster,
% documentation...) et gonfle inutilement la mémoire TeX sur un document long
% (~300 boîtes) : on ne charge que ce qui est réellement utilisé.
\usepackage[skins,breakable]{tcolorbox}
\usepackage{graphicx}
\usepackage{colortbl}
\usepackage{booktabs}
\usepackage{tabularx}
\usepackage{diagbox} % case d'en-tête diagonale des tableaux à double entrée
% Colonnes extensibles centrée / gauche / droite (C, L, R), souvent utilisées par les modèles
\newcolumntype{C}{>{\centering\arraybackslash}X}
\newcolumntype{L}{>{\raggedright\arraybackslash}X}
\newcolumntype{R}{>{\raggedleft\arraybackslash}X}
\usepackage{array}
\usepackage{multicol}
\usepackage{siunitx}
% parse-numbers=false : accepte aussi \SI{2,5 \times 10^{-3}}{...}
\sisetup{locale=UK,output-decimal-marker={.},group-minimum-digits=4,parse-numbers=false}
\DeclareSIUnit\ohm{\ensuremath{\symup{\Omega}}} % U+2126 absent de la police
\DeclareSIUnit\division{div} % réglages d'oscilloscope (ms/div, V/div)
\DeclareSIUnit\tr{tr} % tours (vitesse de rotation, ex. \SI{1500}{\tr\per\minute})
\usepackage{qrcode}
% \degree : commande fréquemment utilisée par les modèles pour un angle ou une
% température en dehors de \SI{}{} (non fournie par les paquets chargés).
\providecommand{\degree}{^{\circ}}
% \qed (fin de démonstration), utilisé par les modèles sans amsthm
\providecommand{\qed}{\hfill$\square$}
% \barre : double barre des valeurs interdites dans un tableau de variations (tkz-tab)
\providecommand{\barre}{\ensuremath{\|}}
% environnement proof (amsthm non chargé) utilisé par les modèles
\ifdefined\proof\else\newenvironment{proof}[1][Proof]{\par\noindent\textit{#1.}\ }{\qed\par}\fi
\usepackage{lastpage}
\usepackage{hyperref}
\hypersetup{hidelinks}

% ============================================================
% Palette « Pop » OnBuch+ — mêmes hex que la classe P (plus_ui.dart)
% ============================================================
\definecolor{poporange}{HTML}{F59321}
\definecolor{popdark}{HTML}{E07A0C}
\definecolor{popcream}{HTML}{FFF6E8}
\definecolor{popink}{HTML}{1C1714}
\definecolor{poppurple}{HTML}{7A5AE0}
\definecolor{popgreen}{HTML}{2E9E6A}
\definecolor{poppink}{HTML}{E0457B}
\definecolor{popblue}{HTML}{2F7DE1}
\definecolor{popgold}{HTML}{C98A12}
\definecolor{popmuted}{HTML}{8A7F76}
\definecolor{popline}{HTML}{EADFD2}
\definecolor{popgray}{HTML}{8A7F76}
\colorlet{popgrey}{popgray}
% Alias tolérants (le modèle nomme parfois les couleurs différemment)
\colorlet{popwhite}{white}
\colorlet{popblack}{popink}
\colorlet{popred}{poppink}
\colorlet{popyellow}{popgold}
% Teintes claires pour fonds de tableaux / zones
\colorlet{popblueL}{popblue!10!white}
\colorlet{poporangeL}{poporange!14!white}
\colorlet{popgreenL}{popgreen!12!white}
\colorlet{poppurpleL}{poppurple!12!white}
\colorlet{poppinkL}{poppink!12!white}
\colorlet{popgoldL}{popgold!14!white}

\pagecolor{popcream}

% ============================================================
% Typographie
% ============================================================
\defaultfontfeatures{Ligatures=TeX}
\setmainfont{PlusJakartaSans-Medium.ttf}[Path=../../../fonts/,BoldFont=PlusJakartaSans-Bold.ttf]
% Maths en sans-serif (Fira Math) avec chiffres et lettres droites de Plus
% Jakarta, pour que les formules se fondent dans le texte.
\usepackage[math-style=ISO,bold-style=ISO]{unicode-math}
\setmathfont{FiraMath-Regular.otf}[Path=../../../fonts/]
\setmathfont{PlusJakartaSans-Medium.ttf}[Path=../../../fonts/,range={up/num,up/{Latin,latin},\mathup}]
% (icomma retiré : décimales avec point en anglais)
% Caractères mathématiques Unicode absents des polices de texte (Plus Jakarta,
% Archivo Black) : rendus par leur équivalent mathématique, en texte comme en
% formule — sinon ils s'affichent en carré vide (ex. « Arithmétique dans ℤ »).
\usepackage{newunicodechar}
\newunicodechar{ℝ}{\ensuremath{\mathbb{R}}}
\newunicodechar{ℤ}{\ensuremath{\mathbb{Z}}}
\newunicodechar{ℂ}{\ensuremath{\mathbb{C}}}
\newunicodechar{ℕ}{\ensuremath{\mathbb{N}}}
\newunicodechar{ℚ}{\ensuremath{\mathbb{Q}}}
\newunicodechar{𝔻}{\ensuremath{\mathbb{D}}}
\newunicodechar{α}{\ensuremath{\alpha}}
\newunicodechar{β}{\ensuremath{\beta}}
\newunicodechar{γ}{\ensuremath{\gamma}}
\newunicodechar{δ}{\ensuremath{\delta}}
\newunicodechar{ε}{\ensuremath{\varepsilon}}
\newunicodechar{θ}{\ensuremath{\theta}}
\newunicodechar{λ}{\ensuremath{\lambda}}
\newunicodechar{μ}{\ensuremath{\mu}}
\newunicodechar{σ}{\ensuremath{\sigma}}
\newunicodechar{τ}{\ensuremath{\tau}}
\newunicodechar{φ}{\ensuremath{\varphi}}
\newunicodechar{ψ}{\ensuremath{\psi}}
\newunicodechar{ω}{\ensuremath{\omega}}
\newunicodechar{Σ}{\ensuremath{\Sigma}}
% majuscules produites par \MakeUppercase dans les titres (α-aminés -> Α)
\newunicodechar{Α}{\ensuremath{\alpha}}
\newunicodechar{Β}{\ensuremath{\beta}}
\newunicodechar{Γ}{\ensuremath{\Gamma}}
\newunicodechar{Δ}{\ensuremath{\Delta}}
\newunicodechar{Ε}{\ensuremath{\varepsilon}}
\newunicodechar{Θ}{\ensuremath{\Theta}}
\newunicodechar{Λ}{\ensuremath{\Lambda}}
\newunicodechar{Μ}{\ensuremath{\mu}}
\newunicodechar{Τ}{\ensuremath{\tau}}
\newunicodechar{Φ}{\ensuremath{\Phi}}
\newunicodechar{Ψ}{\ensuremath{\Psi}}
\newunicodechar{Ω}{\ensuremath{\Omega}}
\newunicodechar{↦}{\ensuremath{\mapsto}}
\newunicodechar{∈}{\ensuremath{\in}}
\newunicodechar{∉}{\ensuremath{\notin}}
\newunicodechar{∧}{\ensuremath{\wedge}}
\newunicodechar{⇔}{\ensuremath{\Leftrightarrow}}
\newunicodechar{⟺}{\ensuremath{\Longleftrightarrow}}
\newunicodechar{⇒}{\ensuremath{\Rightarrow}}
\newunicodechar{⟹}{\ensuremath{\Longrightarrow}}
\newunicodechar{∘}{\ensuremath{\circ}}
\newunicodechar{∪}{\ensuremath{\cup}}
\newunicodechar{∩}{\ensuremath{\cap}}
\newunicodechar{≡}{\ensuremath{\equiv}}
\newunicodechar{⊂}{\ensuremath{\subset}}
\newunicodechar{⊥}{\ensuremath{\perp}}
\newunicodechar{∃}{\ensuremath{\exists}}
\newunicodechar{∀}{\ensuremath{\forall}}
\newunicodechar{∛}{\ensuremath{\sqrt[3]{\,}}}
\newunicodechar{∜}{\ensuremath{\sqrt[4]{\,}}}
\newunicodechar{⃗}{\ensuremath{{}^{\to}}}
\newunicodechar{ⁿ}{\ifmmode^{n}\else\textsuperscript{\ensuremath{n}}\fi}
\newunicodechar{ˣ}{\ifmmode^{x}\else\textsuperscript{\ensuremath{x}}\fi}
\newunicodechar{ᵇ}{\ifmmode^{b}\else\textsuperscript{\ensuremath{b}}\fi}
\newunicodechar{ʳ}{\ifmmode^{r}\else\textsuperscript{\ensuremath{r}}\fi}
\newunicodechar{ᵘ}{\ifmmode^{u}\else\textsuperscript{\ensuremath{u}}\fi}
\newunicodechar{ʸ}{\ifmmode^{y}\else\textsuperscript{\ensuremath{y}}\fi}
\newunicodechar{ᵃ}{\ifmmode^{a}\else\textsuperscript{\ensuremath{a}}\fi}
\newunicodechar{ᵉ}{\ifmmode^{e}\else\textsuperscript{\ensuremath{e}}\fi}
\newunicodechar{ᵏ}{\ifmmode^{k}\else\textsuperscript{\ensuremath{k}}\fi}
\newunicodechar{ᵖ}{\ifmmode^{p}\else\textsuperscript{\ensuremath{p}}\fi}
\newunicodechar{ᴺ}{\ifmmode^{N}\else\textsuperscript{\ensuremath{N}}\fi}
\newunicodechar{ᶜ}{\ifmmode^{c}\else\textsuperscript{\ensuremath{c}}\fi}
\newunicodechar{ʰ}{\ifmmode^{h}\else\textsuperscript{\ensuremath{h}}\fi}
\newunicodechar{ᵠ}{\ifmmode^{q}\else\textsuperscript{\ensuremath{q}}\fi}
\newunicodechar{ₙ}{\ifmmode_{n}\else\textsubscript{\ensuremath{n}}\fi}
\newunicodechar{ₐ}{\ifmmode_{a}\else\textsubscript{\ensuremath{a}}\fi}
\newunicodechar{ᵢ}{\ifmmode_{i}\else\textsubscript{\ensuremath{i}}\fi}
\newunicodechar{ᵣ}{\ifmmode_{r}\else\textsubscript{\ensuremath{r}}\fi}
\newunicodechar{ₖ}{\ifmmode_{k}\else\textsubscript{\ensuremath{k}}\fi}
\newunicodechar{ᵦ}{\ifmmode_{\beta}\else\textsubscript{\ensuremath{\beta}}\fi}
\newunicodechar{ₓ}{\ifmmode_{x}\else\textsubscript{\ensuremath{x}}\fi}
\newfontfamily\popdisplay{ArchivoBlack-Regular.ttf}[Path=../../../fonts/]
\newfontfamily\poplogo{SpaceGrotesk-Bold.ttf}[Path=../../../fonts/]
\newfontfamily\popsemi{PlusJakartaSans-SemiBold.ttf}[Path=../../../fonts/]
\newfontfamily\popxbold{PlusJakartaSans-ExtraBold.ttf}[Path=../../../fonts/]
\newfontfamily\popxboldls{PlusJakartaSans-ExtraBold.ttf}[Path=../../../fonts/,LetterSpace=3.0]

\setlist[enumerate]{leftmargin=1.7em,itemsep=5pt,topsep=4pt}
\setlist[itemize]{leftmargin=1.7em,itemsep=4pt,topsep=4pt,label={\color{poporange}\textbullet}}
\setlength{\parindent}{0pt}
\setlength{\parskip}{5pt}
\renewcommand{\arraystretch}{1.25}

% pgfplots : style Pop par défaut
\pgfplotsset{
  every axis/.append style={axis line style={popink,line width=1pt},tick style={popink},
    grid style={popline,line width=0.5pt},label style={font=\small},tick label style={font=\footnotesize}},
  popcurve/.style={poporange,line width=1.8pt,no markers,smooth},
}
% Même style disponible dans un \draw TikZ simple (hors pgfplots)
\tikzset{popcurve/.style={poporange,line width=1.8pt}}
\tikzset{popgrid/.style={popline,very thin}}
\colorlet{popcurve}{poporange} % le nom du style est parfois utilisé comme couleur

% ============================================================
% Pill / squiggle
% ============================================================
\newcommand{\poppill}[3]{%
  \tikz[baseline=-0.55ex]\node[draw=popink,line width=1.2pt,rounded corners=2.6pt,%
    fill=#2,inner xsep=6.5pt,inner ysep=3pt]{%
    \popxboldls\fontsize{7}{7}\selectfont\color{#3}\MakeUppercase{#1}};%
}
\newcommand{\popsquiggle}[1]{%
  \tikz[baseline=-0.4ex]{
    \draw[#1,line width=2.6pt,line cap=round]
      plot [smooth] coordinates {(0,0.09) (0.34,0.01) (0.68,0.21) (1.02,0.01) (1.36,0.10)};
  }%
}
% Mot-clé surligné (marqueur orange sous le texte)
\newcommand{\cle}[1]{{\popxbold\color{popdark}#1}}

% ============================================================
% En-tête / pied
% ============================================================
\pagestyle{fancy}
\fancyhf{}
\renewcommand{\headrulewidth}{0pt}
\renewcommand{\footrulewidth}{0pt}
\lhead{\raisebox{-0.15ex}{\poplogo\fontsize{13}{13}\selectfont\color{popink}OnBuch\color{poporange}+}}
\rhead{\poppill{\DOCMATIERE\ \textbullet\ \DOCNIVEAU\ \textbullet\ Chapter \DOCLECON}{white}{popink}}
\lfoot{\popxbold\fontsize{7.6}{7.6}\selectfont\color{popmuted}\MakeUppercase{Signed course OnBuch+}\ \textbullet\ {\popsemi\DOCTITRE}}
\rfoot{\tikz[baseline=-0.6ex]\node[rounded corners=8.5pt,fill=popink,minimum height=17pt,minimum width=17pt,inner xsep=5pt,inner sep=0]{\popxbold\fontsize{8}{8}\selectfont\color{popcream}\thepage\,/\,\pageref*{LastPage}};}

% ============================================================
% Titres
% ============================================================
\newcommand{\chaptitre}[2]{%
  \begin{tcolorbox}[enhanced,colback=poporange,colframe=popink,boxrule=2.6pt,arc=11pt,
      drop shadow=popink,left=15pt,right=15pt,top=12pt,bottom=11pt,before skip=3pt,after skip=18pt]
    \poppill{#2}{popink}{popcream}\par\vspace{9pt}
    {\raggedright\hyphenpenalty=10000\exhyphenpenalty=10000\popdisplay\fontsize{22}{24}\selectfont\color{popcream}\MakeUppercase{#1}\par}
  \end{tcolorbox}
}
% Section numérotée : pastille orange avec le numéro + titre Archivo
\newcounter{popsec}
\newcommand{\coursec}[1]{%
  \stepcounter{popsec}\setcounter{popsub}{0}%
  \par\vspace{16pt}\noindent
  \begin{minipage}{\linewidth}
  \tikz[baseline=-0.7ex]\node[draw=popink,line width=1.6pt,fill=poporange,rounded corners=4pt,minimum size=22pt,inner sep=0]{\popdisplay\fontsize{12}{12}\selectfont\color{popcream}\thepopsec};%
  \hspace{8pt}{\popdisplay\fontsize{15}{17}\selectfont\color{popink}\MakeUppercase{#1}}\par
  \vspace{4pt}\noindent{\color{poporange}\rule{36pt}{3.6pt}}
  \end{minipage}\par\nopagebreak\vspace{8pt}
}
\newcounter{popsub}
\newcommand{\courssub}[1]{%
  \stepcounter{popsub}%
  \par\vspace{9pt}\noindent{\popxbold\fontsize{11.5}{13}\selectfont\color{popink}{\color{poporange}\thepopsec.\thepopsub}\hspace{6pt}#1}\par\nopagebreak\vspace{4pt}
}

% ============================================================
% Boîtes pédagogiques — même recette : fond blanc, bordure épaisse colorée,
% ombre dure encre, badge plein. Titre optionnel [#1].
% ============================================================
\tcbset{popbase/.style={breakable,enhanced,colback=white,boxrule=2.3pt,arc=9pt,
  shadow={3pt}{-3pt}{0pt}{popink},left=14pt,right=14pt,top=11pt,bottom=11pt,before skip=11pt,after skip=15pt,
  fontupper=\popsemi\fontsize{10.6}{14.5}\selectfont\color{popink},
  fontlower=\popsemi\fontsize{10.6}{14.5}\selectfont\color{popink}}}
% Coche dessinée (le glyphe ✓ n'existe pas dans les polices du document)
\newcommand{\popcheck}[1]{\tikz[baseline=-0.5ex]\draw[#1,line width=1.6pt,line cap=round,line join=round](-0.13,0.0)--(-0.04,-0.09)--(0.13,0.1);}
\providecommand{\coche}{\popcheck{popgreen}}
\providecommand{\resultat}[1]{\par\smallskip\noindent\fcolorbox{popgreen}{popgreenL}{\parbox{\dimexpr\linewidth-2\fboxsep-2\fboxrule\relax}{\textbf{Result.} #1}}\par\smallskip}
\newcommand{\popboxhead}[3]{\poppill{#1}{#2}{white}\ifx\relax#3\relax\else\hspace{6pt}{\popxbold\fontsize{10.6}{12}\selectfont\color{popink}#3}\fi\par\vspace{7pt}}

\newtcolorbox{aretenir}[1][]{popbase,colframe=popgreen,before upper={\popboxhead{Key points}{popgreen}{#1}}}
\newtcolorbox{exemplebox}[1][]{popbase,colframe=poporange,before upper={\popboxhead{Exemple}{poporange}{#1}}}
\newtcolorbox{definition}[1][]{popbase,colframe=popblue,before upper={\popboxhead{Definition}{popblue}{#1}}}
\newtcolorbox{propriete}[1][]{popbase,colframe=poppurple,before upper={\popboxhead{Property}{poppurple}{#1}}}
\newtcolorbox{methode}[1][]{popbase,colframe=popgold,before upper={\popboxhead{Method}{popgold}{#1}}}
\newtcolorbox{attention}[1][]{popbase,colframe=poppink,before upper={\popboxhead{Warning}{poppink}{#1}}}
\newtcolorbox{experience}[1][]{popbase,colframe=popblue,colback=popblueL,before upper={\popboxhead{Experiment}{popblue}{#1}}}
% « Le savais-tu ? » : carte violette pleine (contexte, culture, Cameroun)
\newtcolorbox{savaistu}[1][]{popbase,colback=poppurple,colframe=popink,
  fontupper=\popsemi\fontsize{10.4}{14.2}\selectfont\color{white},
  before upper={\poppill{Did you know?}{white}{poppurple}\ifx\relax#1\relax\else\hspace{6pt}{\popxbold\color{white}#1}\fi\par\vspace{7pt}}}
% Exercice résolu : énoncé en haut, \tcblower puis la solution sur fond crème
\newcounter{popexo}\newcounter{popexr}
\newtcolorbox{exoresolu}[1][]{popbase,colframe=poporange,colbacklower=popcream,
  segmentation style={popink,line width=1.2pt,dashed},
  before upper={\stepcounter{popexr}\popboxhead{Worked example \thepopexr}{poporange}{#1}},
  before lower={\poppill{Solution}{popink}{popcream}\par\vspace{6pt}}}
% Exercice (non résolu en place) — la correction est dans la partie Corrigés
\newtcolorbox{exercice}[1][]{popbase,colframe=popink,
  before upper={\stepcounter{popexo}\popboxhead{Exercise \thepopexo}{popink}{#1}}}
% Corrigé
\newtcolorbox{corrige}[1][]{popbase,colframe=popgreen,colback=popgreenL,
  before upper={\popboxhead{Solution}{popgreen}{#1}}}
% Objectifs / prérequis / bilan
\newtcolorbox{objectifs}{popbase,colframe=popink,colback=poporangeL,
  before upper={\popboxhead{Chapter objectives}{popink}{}}}
\newtcolorbox{prerequis}{popbase,colframe=popink,colback=popblueL,
  before upper={\popboxhead{Prerequisites}{popblue}{}}}
\newtcolorbox{bilan}{popbase,colframe=popink,colback=white,boxrule=2.8pt,
  before upper={\popboxhead{Fiche bilan}{poporange}{L'essentiel à connaître pour l'examen}}}
% Figure encadrée avec légende
\newtcolorbox{popfigure}[1][]{enhanced,breakable=false,colback=white,colframe=popline,boxrule=1.4pt,arc=8pt,
  left=10pt,right=10pt,top=10pt,bottom=8pt,before skip=10pt,after skip=12pt,halign=center,
  lower separated=false,halign lower=center,
  fontlower=\popsemi\fontsize{9}{11.5}\selectfont\color{popmuted},
  #1}
\newcommand{\legende}[1]{\par\vspace{6pt}{\popsemi\fontsize{9}{11.5}\selectfont\color{popmuted}#1}}

% ============================================================
% Signature OnBuch+
% ============================================================
\newcommand{\poplogoinline}[1]{{\poplogo\fontsize{#1}{#1}\selectfont\color{popink}OnBuch\color{poporange}+}}
\newcommand{\signatureonbuch}{%
  \begin{tcolorbox}[enhanced,colback=popink,colframe=popink,boxrule=2.2pt,arc=10pt,
      drop shadow=poporange,left=14pt,right=14pt,top=10pt,bottom=10pt,before skip=0pt,after skip=0pt]
    \tikz[baseline=-0.7ex]\node[circle,fill=popgreen,draw=popcream,line width=1.3pt,minimum size=22pt,inner sep=0]{\popcheck{white}};%
    \hspace{9pt}\begin{minipage}[c]{0.62\linewidth}
      {\popxbold\fontsize{12}{13}\selectfont\color{popcream}Signed\ \textbullet\ {\poplogo OnBuch\color{poporange}+}}\par\vspace{2pt}
      {\raggedright\popsemi\fontsize{8}{10.5}\selectfont\color{popline}\MakeUppercase{OnBuch+ teaching team}\\\MakeUppercase{Follows the official MINESEC syllabus}\par}
    \end{minipage}\hfill
    {\poplogo\fontsize{22}{22}\selectfont\color{popcream}OnBuch\color{poporange}+}
  \end{tcolorbox}%
}

% ============================================================
% Page de garde
% #1 titre · #2 matière · #3 niveau · #4 module · #5 n° leçon · #6 durée ·
% #7 description / objectifs (texte ou itemize)
% ============================================================
\newcommand{\pagegarde}[7]{%
  \thispagestyle{empty}
  \newgeometry{a4paper,margin=1.7cm,top=1.6cm,bottom=1.4cm}%
  \begin{tikzpicture}[remember picture,overlay]
    \fill[poporange!18] ([xshift=-3.2cm,yshift=-3.6cm]current page.north east) circle (4.6cm);
    \fill[poppurple!14] ([xshift=2.4cm,yshift=7.6cm]current page.south west) circle (3.8cm);
  \end{tikzpicture}%
  {\poplogo\fontsize{30}{30}\selectfont\color{popink}OnBuch\color{poporange}+}\hfill
  \raisebox{0.6ex}{\poppill{Signed course \textbullet\ Premium}{popink}{popcream}}\par
  \vspace{1pt}\popsquiggle{poppurple}\par
  \vspace{8pt}
  \begin{tcolorbox}[enhanced,colback=poporange,colframe=popink,boxrule=2.8pt,arc=13pt,
      drop shadow=popink,left=18pt,right=18pt,top=12pt,bottom=13pt,after skip=10pt]
    \poppill{#2 \textbullet\ #3}{popink}{popcream}\hspace{5pt}\poppill{#4}{white}{popink}\par\vspace{8pt}
    {\popdisplay\fontsize{14}{14}\selectfont\color{popink}CHAPTER #5}\par\vspace{4pt}
    {\raggedright\hyphenpenalty=10000\exhyphenpenalty=10000\popdisplay\fontsize{25}{27}\selectfont\color{popcream}\MakeUppercase{#1}\par}
    \vspace{8pt}
    {\popxbold\fontsize{9.5}{9.5}\selectfont\color{popink}\MakeUppercase{Suggested time: #6}}
  \end{tcolorbox}
  \begin{tcolorbox}[enhanced,colback=white,colframe=popink,boxrule=2.2pt,arc=10pt,
      drop shadow=popink,left=16pt,right=16pt,top=10pt,bottom=10pt,after skip=8pt]
    \poppill{In this chapter}{poppurple}{white}\par\vspace{5pt}
    \popsemi\fontsize{9.3}{12.6}\selectfont\color{popink}#7
  \end{tcolorbox}
  \noindent\begin{minipage}[t]{0.487\linewidth}
    \begin{tcolorbox}[enhanced,colback=popgreen,colframe=popink,boxrule=2.2pt,arc=10pt,
        drop shadow=popink,left=11pt,right=11pt,top=10pt,bottom=10pt,height=3.1cm,valign=center]
      \tcbox[colback=white,colframe=popink,boxrule=1.3pt,arc=5pt,boxsep=0pt,nobeforeafter,
        left=3pt,right=3pt,top=3pt,bottom=3pt,valign=center]{\qrcode[height=1.9cm]{https://play.google.com/store/apps/details?id=com.luvvix.onbuch&hl=en}}%
      \hspace{8pt}\begin{minipage}[c]{0.5\linewidth}
        {\popdisplay\fontsize{11}{12.5}\selectfont\color{white}\MakeUppercase{Download OnBuch}}\par\vspace{4pt}
        {\raggedright\popsemi\fontsize{8.4}{11}\selectfont\color{white}Free \textbullet\ Google Play\\AI tutor, past papers, results\par}
      \end{minipage}
    \end{tcolorbox}
  \end{minipage}\hfill
  \begin{minipage}[t]{0.487\linewidth}
    \begin{tcolorbox}[enhanced,colback=poppurple,colframe=popink,boxrule=2.2pt,arc=10pt,
        drop shadow=popink,left=11pt,right=11pt,top=10pt,bottom=10pt,height=3.1cm,valign=center]
      \tikz[baseline=-0.7ex]\node[circle,fill=white,draw=popink,line width=1.3pt,minimum size=1.6cm,inner sep=0]{\popdisplay\fontsize{24}{24}\selectfont\color{poppurple}+};%
      \hspace{8pt}\begin{minipage}[c]{0.6\linewidth}
        {\popdisplay\fontsize{11}{12.5}\selectfont\color{white}\MakeUppercase{Go OnBuch+}}\par\vspace{4pt}
        {\raggedright\popsemi\fontsize{8.4}{11}\selectfont\color{white}All signed courses, unlimited AI tutor, PDF sheets — OnBuch+ tab in the app\par}
      \end{minipage}
    \end{tcolorbox}
  \end{minipage}
  \vspace{8pt}
  \popcontenu
  \vfill
  \signatureonbuch
  \restoregeometry
  \newpage
}

% Rangée « Ce cours contient » (page de garde)
\newcommand{\poptile}[3]{% #1 couleur #2 gros texte #3 légende
  \begin{tcolorbox}[enhanced,colback=#1,colframe=popink,boxrule=2pt,arc=9pt,shadow={2.5pt}{-2.5pt}{0pt}{popink},
      left=6pt,right=6pt,top=9pt,bottom=9pt,halign=center,valign=center,height=1.75cm,width=\linewidth,nobeforeafter]
    {\popdisplay\fontsize{12.5}{13}\selectfont\color{white}\MakeUppercase{#2}}\par\vspace{3pt}
    {\centering\popsemi\fontsize{7.8}{9.5}\selectfont\color{white}#3\par}
  \end{tcolorbox}}
\newcommand{\popcontenu}{%
  \par\noindent{\popxboldls\fontsize{7.5}{7.5}\selectfont\color{popmuted}THIS SIGNED COURSE CONTAINS}\par\vspace{7pt}
  \noindent\begin{minipage}[t]{0.235\linewidth}\poptile{popblue}{Course}{complete, illustrated, step by step}\end{minipage}\hfill
  \begin{minipage}[t]{0.235\linewidth}\poptile{poporange}{Methods}{and worked examples}\end{minipage}\hfill
  \begin{minipage}[t]{0.235\linewidth}\poptile{poppink}{Exercises}{exam style + solutions}\end{minipage}\hfill
  \begin{minipage}[t]{0.235\linewidth}\poptile{popgreen}{Summary sheet}{the essentials to remember}\end{minipage}\par}

% Sommaire
\newtcolorbox{sommaire}{enhanced,colback=white,colframe=popink,boxrule=2.2pt,arc=10pt,
  drop shadow=popink,left=18pt,right=18pt,top=13pt,bottom=13pt,before skip=6pt,after skip=20pt,
  before upper={\poppill{Contents}{poppurple}{white}\par\vspace{9pt}\popsemi\fontsize{10.6}{14.5}\selectfont\color{popink}}}

% Signature de fin de cours — poussée tout en bas de la dernière page
\newcommand{\findecours}{%
  \par\vfill\nopagebreak
  \begin{center}\popsquiggle{poporange}\end{center}\vspace{4pt}
  \signatureonbuch
}
\AtBeginDocument{\def\llbracket{\mathopen{[\mkern-3.5mu[}}\def\rrbracket{\mathclose{]\mkern-3.5mu]}}} % intervalles d'entiers (glyphe ⟦ absent de la police)
% Glyphes absents de Fira Math : redéfinis à partir de glyphes disponibles
\AtBeginDocument{%
  \def\setminus{\mathbin{\backslash}}\let\smallsetminus\setminus
  \def\vdots{\mathord{\vbox{\baselineskip3.2pt\lineskiplimit0pt\kern4pt\hbox{.}\hbox{.}\hbox{.}}}}%
  \def\ddots{\mathinner{\mkern1mu\raise6.5pt\hbox{.}\mkern2mu\raise3.6pt\hbox{.}\mkern2mu\raise0.7pt\hbox{.}\mkern1mu}}%
  \def\triangle{\mathord{\scalebox{0.9}{$\symup{\Delta}$}}}%
  \def\nabla{\mathord{\raisebox{\depth}{\scalebox{1}[-1]{$\symup{\Delta}$}}}}%
  \def\checkmark{\popcheck{popdark}}%
  \def\star{\mathbin{\ast}}\def\bigstar{\mathord{\ast}}%
  \def\diamond{\mathbin{\scalebox{0.6}{$\Diamond$}}}%
  \def\bigcup{\mathop{\mathchoice{\raisebox{-0.3em}{\scalebox{1.7}{$\cup$}}}{\scalebox{1.25}{$\cup$}}{\cup}{\cup}}\displaylimits}%
  \def\bigcap{\mathop{\mathchoice{\raisebox{-0.3em}{\scalebox{1.7}{$\cap$}}}{\scalebox{1.25}{$\cap$}}{\cap}{\cap}}\displaylimits}%
  \def\lesseqqgtr{\mathrel{\lessgtr}}\def\bigoplus{\mathop{\oplus}\displaylimits}%
}
\providecommand{\ssi}{if and only if} % abréviation fréquente
\AtBeginDocument{\providecommand{\wideparen}{\overparen}} % arc de cercle

% Fonctions usuelles que les modèles utilisent sans que amsmath les fournisse
\providecommand{\arsinh}{\operatorname{arsinh}}\providecommand{\arcosh}{\operatorname{arcosh}}
\providecommand{\artanh}{\operatorname{artanh}}\providecommand{\arsech}{\operatorname{arsech}}
\providecommand{\arcsch}{\operatorname{arcsch}}\providecommand{\arcoth}{\operatorname{arcoth}}
\providecommand{\arcsinh}{\operatorname{arsinh}}\providecommand{\arccosh}{\operatorname{arcosh}}
\providecommand{\arctanh}{\operatorname{artanh}}\providecommand{\sech}{\operatorname{sech}}
\providecommand{\csch}{\operatorname{csch}}\providecommand{\arcsec}{\operatorname{arcsec}}
\providecommand{\arccsc}{\operatorname{arccsc}}\providecommand{\arccot}{\operatorname{arccot}}
\providecommand{\sgn}{\operatorname{sgn}}\providecommand{\lcm}{\operatorname{lcm}}
\providecommand{\Var}{\operatorname{Var}}\providecommand{\Cov}{\operatorname{Cov}}

%% FIGFIX-BEGIN v4 — correctifs globaux des figures (docs/onbuch-generation/tools/figfix)
\usepackage{adjustbox}\usepackage{etoolbox}
% toute figure TikZ/pgfplots plus large que la ligne est réduite à la largeur disponible
\BeforeBeginEnvironment{tikzpicture}{\begin{adjustbox}{max width=\linewidth}}
\AfterEndEnvironment{tikzpicture}{\end{adjustbox}}
% légendes pgfplots lisibles par-dessus les courbes
\pgfplotsset{every axis/.append style={legend style={fill=white,fill opacity=0.92,text opacity=1,draw=black!15,inner sep=3pt}}}
% étiquettes d'arêtes (syntaxe edge["..."]) avec fond blanc
\tikzset{every edge quotes/.append style={fill=white,inner sep=1.2pt}}
%% FIGFIX-END
\begin{document}

\pagegarde{Hyperbolic and Inverse hyperbolic Functions}{Further Mathematics}{Upper Sixth}{Upper Sixth — Further mathematics}{1}{13 periods}{This chapter introduces hyperbolic functions, their identities, inverses, calculus, and series expansions. It equips students with the analytical techniques required for the GCE Advanced Level Further Mathematics examination and for modelling natural phenomena.
\vspace{4pt}
\begin{multicols}{2}\small
\begin{itemize}
\item Define the six basic hyperbolic functions and compute their numerical values.
\item Sketch and interpret the graphs of hyperbolic functions, noting asymptotes and symmetry.
\item Derive and apply fundamental hyperbolic identities, including Osborn's rule.
\item Define inverse hyperbolic functions, sketch their graphs, and express them in logarithmic form.
\item Solve equations involving hyperbolic and inverse hyperbolic functions.
\item Differentiate and integrate hyperbolic and inverse hyperbolic functions.
\end{itemize}\end{multicols}}

\begin{sommaire}
\begin{multicols}{2}
\begin{enumerate}
\item 1. Definition and Numerical Evaluation of Hyperbolic Functions
\item 2. Graphs of Hyperbolic Functions
\item 3. Hyperbolic Identities and Osborn's Rule
\item 4. Inverse Hyperbolic Functions: Definitions, Graphs, and Logarithmic Forms
\item 5. Solving Hyperbolic Equations
\item 6. Calculus of Hyperbolic Functions: Differentiation and Integration
\item 7. Advanced Applications: Hyperbolic Substitutions and Maclaurin Series
\item Integration activity
\item Methods and worked examples
\item Exercises
\item Solutions to the exercises
\item Summary sheet
\end{enumerate}
\end{multicols}
\end{sommaire}

\begin{objectifs}
\begin{itemize}
\item Define the six basic hyperbolic functions and compute their numerical values.
\item Sketch and interpret the graphs of hyperbolic functions, noting asymptotes and symmetry.
\item Derive and apply fundamental hyperbolic identities, including Osborn's rule.
\item Define inverse hyperbolic functions, sketch their graphs, and express them in logarithmic form.
\item Solve equations involving hyperbolic and inverse hyperbolic functions.
\item Differentiate and integrate hyperbolic and inverse hyperbolic functions.
\item Use hyperbolic substitutions to evaluate integrals that contain quadratic expressions.
\item Obtain Maclaurin series expansions for sinh x, cosh x, and tanh x and use them for approximations.
\end{itemize}
\end{objectifs}

\begin{prerequis}
\begin{itemize}
\item Exponential and logarithmic functions (Lower Sixth).
\item Trigonometric identities and their derivations.
\item Basic differentiation and integration techniques (chain rule, substitution).
\item Maclaurin series for elementary functions (e\^{}x, sin x, cos x).
\item Graph sketching using derivatives and asymptotes.
\end{itemize}
\end{prerequis}

\newpage

\coursec{Definition and Numerical Evaluation of Hyperbolic Functions}

Imagine a suspension bridge over the Wouri River: the main cables hang in a graceful curve called a \emph{catenary}, which is described by the hyperbolic cosine function. Engineers use these curves to design safe, elegant structures, and the same functions appear in the study of hanging chains, transmission lines, special relativity and the shape of soap films. In this chapter you will discover the mathematics behind such real-world shapes and learn to manipulate them with the same tools used in advanced calculus.

\courssub{From the Unit Circle to the Unit Hyperbola}

You are familiar with the trigonometric functions defined on the unit circle $x^2 + y^2 = 1$. A point on this circle corresponding to an angle $\theta$ has coordinates $(\cos\theta,\;\sin\theta)$, and the identity $\cos^2\theta+\sin^2\theta=1$ is simply the statement that the point lies on the circle.

Hyperbolic functions arise from a completely analogous construction on the \emph{unit hyperbola} $x^2 - y^2 = 1$ (right branch). Let $t$ be a real parameter. The coordinates of a point on this hyperbola are given by
\[
x = \cosh t,\qquad y = \sinh t.
\]
Here $t$ represents \emph{twice the area} of the hyperbolic sector bounded by the hyperbola, the $x$-axis and the line joining the origin to the point $(\cosh t,\sinh t)$. This geometric interpretation is the deep reason why hyperbolic functions satisfy identities that mirror those of trigonometric functions, but with crucial sign differences (Osborn's rule, studied later in this chapter).

\begin{popfigure}
\begin{tikzpicture}[scale=1.2]
% Axes
\draw[->] (-2.5,0) -- (2.5,0) node[right] {$x$};
\draw[->] (0,-2) -- (0,2) node[above] {$y$};
% Hyperbola x^2 - y^2 = 1 (right branch)
\draw[popblue, thick, domain=-1.8:1.8, smooth, variable=\t]
  plot ({cosh(\t)}, {sinh(\t)});
% Asymptotes
\draw[popmuted, dashed] (-2,-2) -- (2,2);
\draw[popmuted, dashed] (-2,2) -- (2,-2);
% Area shading for t=1 (sector)
\fill[popblueL, opacity=0.35] (0,0) -- plot[domain=0:1, variable=\t] ({cosh(\t)}, {sinh(\t)}) -- cycle;
\node[popdark, font=\small] at (1.35,0.28) {Area $=\dfrac{t}{2}$};
% Points for t = -1, -0.5, 0, 0.5, 1
\foreach \tval/\lbl in {-1/{t=-1}, -0.5/{t=-0.5}, 0/{t=0}, 0.5/{t=0.5}, 1/{t=1}}{
  \pgfmathsetmacro{\xval}{cosh(\tval)}
  \pgfmathsetmacro{\yval}{sinh(\tval)}
  \fill[poporange] (\xval,\yval) circle (2pt);
  \node[poporange, anchor=west, font=\tiny] at (\xval,\yval) {$\lbl$};
}
% Labels
\node[popblue, anchor=south west] at (1.6,1.7) {$x^2-y^2=1$};
\node[popmuted, anchor=south east, font=\tiny] at (1.9,1.9) {$y=x$};
\node[popmuted, anchor=north east, font=\tiny] at (1.9,-1.9) {$y=-x$};
\end{tikzpicture}
\legende{The unit hyperbola $x^2-y^2=1$ with points $(\cosh t,\sinh t)$ for several values of $t$. The shaded sector has area $t/2$.}
\end{popfigure}

\courssub{Definitions via Exponential Functions}

While the geometric definition is beautiful, the algebraic definition using exponentials is the one we use for calculations, and it is the definition required by the GCE Advanced Level syllabus.

\begin{definition}[Hyperbolic Sine and Cosine]
For any real number $x$,
\[
\sinh x = \frac{e^x - e^{-x}}{2},\qquad 
\cosh x = \frac{e^x + e^{-x}}{2}.
\]
The names are read ``shine $x$'' (or ``sinch $x$'') and ``cosh $x$''.
\end{definition}

From these two basic functions we define the other four hyperbolic functions, exactly as we define $\tan$, $\cot$, $\sec$, $\csc$ from $\sin$ and $\cos$.

\begin{definition}[The Six Hyperbolic Functions]
For $x\in\mathbb{R}$ (with the usual restrictions where denominators vanish):
\begin{align*}
\tanh x &= \frac{\sinh x}{\cosh x} = \frac{e^x - e^{-x}}{e^x + e^{-x}}, \\
\coth x &= \frac{\cosh x}{\sinh x} = \frac{e^x + e^{-x}}{e^x - e^{-x}},\quad x\neq 0,\\
\operatorname{sech} x &= \frac{1}{\cosh x} = \frac{2}{e^x + e^{-x}}, \\
\operatorname{csch} x &= \frac{1}{\sinh x} = \frac{2}{e^x - e^{-x}},\quad x\neq 0.
\end{align*}
\end{definition}

\begin{attention}[Notation]
There is no standard keyboard symbol for hyperbolic functions: always write them out in full ($\sinh$, $\cosh$, $\tanh$, $\coth$, $\operatorname{sech}$, $\operatorname{csch}$). Note also that $\sinh^2 x$ means $(\sinh x)^2$, exactly as for trigonometric functions.
\end{attention}

\courssub{Domain, Range, Parity and Values at Zero}

Because they are built from exponentials, the domains and ranges are straightforward to determine. For example, since $e^x>0$ and $e^{-x}>0$, we have $\cosh x > 0$ for all $x$; moreover, by the inequality of arithmetic and geometric means, $\cosh x = \dfrac{e^x+e^{-x}}{2} \ge \sqrt{e^x \cdot e^{-x}} = 1$, with equality only at $x=0$.

\begin{propriete}[Domains, Ranges and Parity]
\begin{center}
\begin{tabularx}{\linewidth}{|l|X|X|X|}\hline
\rowcolor{popblueL}
\textbf{Function} & \textbf{Domain} & \textbf{Range} & \textbf{Parity} \\ \hline
$\sinh x$ & $\mathbb{R}$ & $\mathbb{R}$ & Odd: $\sinh(-x)=-\sinh x$ \\ \hline
$\cosh x$ & $\mathbb{R}$ & $[1,+\infty)$ & Even: $\cosh(-x)=\cosh x$ \\ \hline
$\tanh x$ & $\mathbb{R}$ & $(-1,1)$ & Odd \\ \hline
$\coth x$ & $\mathbb{R}\setminus\{0\}$ & $(-\infty,-1)\cup(1,+\infty)$ & Odd \\ \hline
$\operatorname{sech} x$ & $\mathbb{R}$ & $(0,1]$ & Even \\ \hline
$\operatorname{csch} x$ & $\mathbb{R}\setminus\{0\}$ & $\mathbb{R}\setminus\{0\}$ & Odd \\ \hline
\end{tabularx}
\end{center}
The parity claims follow immediately from the definitions, since replacing $x$ by $-x$ simply interchanges $e^x$ and $e^{-x}$.
\end{propriete}

\begin{aretenir}[Values at $x=0$]
\[
\sinh 0 = 0,\quad \cosh 0 = 1,\quad \tanh 0 = 0,\quad \operatorname{sech} 0 = 1.
\]
$\coth 0$ and $\operatorname{csch} 0$ are undefined (the graphs have the $y$-axis as a vertical asymptote).
\end{aretenir}

\courssub{The Fundamental Identity}

Just as $\cos^2\theta+\sin^2\theta=1$ for the unit circle, the unit hyperbola gives the fundamental hyperbolic identity.

\begin{propriete}[Fundamental Hyperbolic Identity]
For all $x\in\mathbb{R}$,
\[
\cosh^2 x - \sinh^2 x = 1.
\]
\emph{Proof (examinable):} Substitute the exponential definitions:
\begin{align*}
\cosh^2 x - \sinh^2 x 
&= \left(\frac{e^x+e^{-x}}{2}\right)^2 - \left(\frac{e^x-e^{-x}}{2}\right)^2 \\
&= \frac{1}{4}\Bigl[\bigl(e^{2x}+2+e^{-2x}\bigr) - \bigl(e^{2x}-2+e^{-2x}\bigr)\Bigr] \\
&= \frac{1}{4}(4) = 1.
\end{align*}
Alternatively, factorise as a difference of two squares:
\[
\cosh^2 x - \sinh^2 x = (\cosh x - \sinh x)(\cosh x + \sinh x) = e^{-x}\cdot e^{x} = 1,
\]
since $\cosh x + \sinh x = e^x$ and $\cosh x - \sinh x = e^{-x}$.
\end{propriete}

\begin{aretenir}[Two Companion Identities]
Dividing the fundamental identity by $\cosh^2 x$ or by $\sinh^2 x$ yields (valid where the denominators are non-zero):
\[
1 - \tanh^2 x = \operatorname{sech}^2 x,\qquad \coth^2 x - 1 = \operatorname{csch}^2 x.
\]
\end{aretenir}

\begin{attention}[Sign Change]
The single most common error with hyperbolic functions is writing $\cosh^2 x + \sinh^2 x = 1$ by analogy with the trigonometric identity. The correct identity is $\cosh^2 x \cle{-} \sinh^2 x = 1$. Note the useful consequence: $\cosh^2 x + \sinh^2 x$ is \emph{not} constant; in fact $\cosh^2 x + \sinh^2 x = \cosh 2x$ (a double-angle identity proved later in this chapter).
\end{attention}

\courssub{Numerical Evaluation: Exact and Approximate Values}

The exponential form makes exact evaluation possible when $x$ is the logarithm of a rational number. For other values we use a calculator or series approximations (Maclaurin series will be studied later in this chapter).

\begin{methode}[Evaluating Hyperbolic Functions Exactly]
To find the exact value of a hyperbolic function at $x = \ln a$ ($a>0$):
\begin{enumerate}
\item Write the function in exponential form.
\item Substitute $e^{\ln a} = a$ and $e^{-\ln a} = \dfrac{1}{a}$.
\item Simplify the resulting algebraic fraction.
\end{enumerate}
\end{methode}

\begin{exemplebox}[Exact values at $\ln 2$ and $\ln 3$]
Compute $\sinh(\ln 2)$, $\cosh(\ln 2)$, $\tanh(\ln 2)$, and the same for $\ln 3$.
\begin{align*}
\sinh(\ln 2) &= \frac{e^{\ln 2} - e^{-\ln 2}}{2} = \frac{2 - \frac12}{2} = \frac{3}{4}, \\
\cosh(\ln 2) &= \frac{2 + \frac12}{2} = \frac{5}{4}, \\
\tanh(\ln 2) &= \frac{\sinh(\ln 2)}{\cosh(\ln 2)} = \frac{3/4}{5/4} = \frac{3}{5}. \\[6pt]
\sinh(\ln 3) &= \frac{3 - \frac13}{2} = \frac{4}{3}, \\
\cosh(\ln 3) &= \frac{3 + \frac13}{2} = \frac{5}{3}, \\
\tanh(\ln 3) &= \frac{4/3}{5/3} = \frac{4}{5}.
\end{align*}
\end{exemplebox}

\begin{center}
\begin{tabularx}{\linewidth}{|c|X|X|X|}\hline
\rowcolor{popblueL}
$x$ & $\sinh x$ & $\cosh x$ & $\tanh x$ \\ \hline
$0$ & $0$ & $1$ & $0$ \\ \hline
$\ln 2$ & $\dfrac{3}{4}$ & $\dfrac{5}{4}$ & $\dfrac{3}{5}$ \\ \hline
$\ln 3$ & $\dfrac{4}{3}$ & $\dfrac{5}{3}$ & $\dfrac{4}{5}$ \\ \hline
\end{tabularx}
\end{center}

For values that are not logarithms of simple numbers, we use a scientific calculator, which has dedicated \texttt{hyp} keys. Note that hyperbolic functions are \emph{not} angular functions: the argument $x$ is a pure number, so the degree/radian mode of the calculator does not affect the result of \texttt{sinh}, \texttt{cosh} or \texttt{tanh}.

\begin{exemplebox}[Approximate values]
Use a calculator to find $\sinh 1.5$, $\cosh 1.5$, $\tanh 1.5$ to four decimal places.
\[
\sinh 1.5 \approx 2.1293,\quad \cosh 1.5 \approx 2.3524,\quad \tanh 1.5 \approx 0.9051.
\]
Check: $\cosh^2 1.5 - \sinh^2 1.5 \approx 2.3524^2 - 2.1293^2 = 5.5338 - 4.5339 \approx 0.9999 \approx 1$ (the small discrepancy is due to rounding).
\end{exemplebox}

\courssub{Series Representation (Preview)}

The Maclaurin series of $e^x$ is $\displaystyle\sum_{n=0}^{\infty} \frac{x^n}{n!}$. Substituting into the definitions gives the series for $\sinh x$ and $\cosh x$. Notice the striking similarity to the series for $\sin x$ and $\cos x$, but \emph{without alternating signs}.

\begin{propriete}[Maclaurin Series]
\[
\sinh x = x + \frac{x^3}{3!} + \frac{x^5}{5!} + \cdots = \sum_{n=0}^{\infty} \frac{x^{2n+1}}{(2n+1)!},
\]
\[
\cosh x = 1 + \frac{x^2}{2!} + \frac{x^4}{4!} + \cdots = \sum_{n=0}^{\infty} \frac{x^{2n}}{(2n)!}.
\]
Both series converge for all $x\in\mathbb{R}$. (These will be derived properly in the section on Maclaurin series.)
\end{propriete}

\begin{attention}[Do Not Confuse $\sinh$ with $\sin$!]
The series for $\sin x$ is $x - \dfrac{x^3}{3!} + \dfrac{x^5}{5!} - \cdots$ (alternating signs).
The series for $\sinh x$ is $x + \dfrac{x^3}{3!} + \dfrac{x^5}{5!} + \cdots$ (all signs positive).
A common mistake is to write $\sinh x \approx x - \dfrac{x^3}{6}$ for small $x$. The correct approximation is $\sinh x \approx x + \dfrac{x^3}{6}$.
\end{attention}

\courssub{Worked Examination-Style Example}

\begin{exoresolu}[GCE Style]
Given that $y = \sinh x$, express $e^x$ in terms of $y$ and hence show that
\[
\sinh^{-1} y = \ln\bigl(y + \sqrt{y^2+1}\bigr).
\]
(We will study inverse hyperbolic functions properly later in this chapter; here we simply derive the logarithmic form.)
\tcblower
\textbf{Solution:}
We have $y = \sinh x = \dfrac{e^x - e^{-x}}{2}$.
Multiply both sides by $2e^x$ (which is positive, so no sign issues arise):
\[
2y e^x = e^{2x} - 1 \quad\Longrightarrow\quad e^{2x} - 2y e^x - 1 = 0.
\]
This is a quadratic equation in $e^x$. Applying the quadratic formula:
\[
e^x = \frac{2y \pm \sqrt{4y^2 + 4}}{2} = y \pm \sqrt{y^2+1}.
\]
Since $e^x > 0$ for all real $x$, and $\sqrt{y^2+1} > \sqrt{y^2} = |y| \ge y$, the expression $y - \sqrt{y^2+1}$ is always negative and must be rejected. We therefore take the positive sign:
\[
e^x = y + \sqrt{y^2+1} \quad\Longrightarrow\quad x = \ln\bigl(y + \sqrt{y^2+1}\bigr).
\]
Hence $\displaystyle \sinh^{-1} y = \ln\bigl(y + \sqrt{y^2+1}\bigr)$, as required.
\end{exoresolu}

\courssub{Exercises}

\begin{exercice}[Practice]
\begin{enumerate}
\item Find the exact values of $\sinh(\ln 5)$, $\cosh(\ln 5)$, $\tanh(\ln 5)$.
\item Evaluate to four decimal places: $\sinh 2.3$, $\cosh 2.3$, $\tanh 2.3$, $\operatorname{sech} 2.3$.
\item Prove that $\cosh(\ln x) = \dfrac{x^2+1}{2x}$ for $x>0$.
\item Show that $\tanh(\ln x) = \dfrac{x^2-1}{x^2+1}$ for $x>0$.
\item If $\sinh x = \dfrac{5}{12}$, find the exact values of $\cosh x$ and $\tanh x$.
\item Use the series for $\sinh x$ to approximate $\sinh 0.2$ up to and including the $x^5$ term. Compare with the calculator value.
\end{enumerate}
\end{exercice}

\begin{corrige}[Exercise 1]
\begin{enumerate}
\item $\sinh(\ln 5) = \dfrac{5 - \frac15}{2} = \dfrac{12}{5}$;\quad $\cosh(\ln 5) = \dfrac{5 + \frac15}{2} = \dfrac{13}{5}$;\quad $\tanh(\ln 5) = \dfrac{12/5}{13/5} = \dfrac{12}{13}$.
\item $\sinh 2.3 \approx 4.9370$, $\cosh 2.3 \approx 5.0372$, $\tanh 2.3 \approx 0.9801$, $\operatorname{sech} 2.3 = \dfrac{1}{\cosh 2.3} \approx 0.1985$.
\item $\cosh(\ln x) = \dfrac{e^{\ln x} + e^{-\ln x}}{2} = \dfrac{x + \frac1x}{2} = \dfrac{x^2+1}{2x}$.
\item $\tanh(\ln x) = \dfrac{x - \frac1x}{x + \frac1x} = \dfrac{x^2-1}{x^2+1}$ (multiplying numerator and denominator by $x$).
\item $\cosh^2 x = 1 + \sinh^2 x = 1 + \dfrac{25}{144} = \dfrac{169}{144}$, so $\cosh x = \dfrac{13}{12}$ (we take the positive root because $\cosh x \ge 1$ for all $x$). Then $\tanh x = \dfrac{5/12}{13/12} = \dfrac{5}{13}$.
\item $\sinh 0.2 \approx 0.2 + \dfrac{0.2^3}{6} + \dfrac{0.2^5}{120} = 0.2 + 0.0013333 + 0.0000027 \approx 0.2013360$. Calculator value: $\sinh 0.2 \approx 0.2013360$. The agreement is excellent.
\end{enumerate}
\end{corrige}

\begin{savaistu}[The Catenary and Suspension Bridges]
The shape of a perfectly flexible hanging cable under its own weight is a \emph{catenary}, with equation $y = a\cosh\!\left(\dfrac{x}{a}\right)$. This is \emph{not} a parabola — a very common misconception, held even by Galileo for some time. Suspension bridges such as the first Wouri bridge in Douala use this principle: the main cables follow a catenary, which distributes tension optimally. The hyperbolic cosine is essentially the only function equal to its own second derivative (up to scaling), which is exactly why it solves the equilibrium equation of a hanging chain. The word ``catenary'' comes from the Latin \emph{catena}, meaning chain.
\end{savaistu}

\coursec{Graphs of Hyperbolic Functions}

In the previous section we defined the six hyperbolic functions from the exponential function and learned how to compute their numerical values. A table of values, however, only gives a snapshot. To truly \emph{understand} these functions --- and later to solve equations, differentiate, integrate and sketch inverse functions --- we need their \cle{graphs}. In this section we sketch all six hyperbolic curves rigorously, using the tools you already master: domain, parity (symmetry), limits (asymptotes) and the sign of the derivative (monotonicity). We then compare these curves with their trigonometric cousins and finish with graph transformations.

\courssub{General strategy for sketching a hyperbolic curve}

Every graph in this section is obtained by the same four-step programme, which you should apply automatically at the GCE examination.

\begin{methode}[Sketching a hyperbolic curve]
\begin{enumerate}
\item \textbf{Domain.} Determine where the function is defined (watch out for division by $\sinh x$, which vanishes only at $x=0$).
\item \textbf{Symmetry.} Test $f(-x)$: if $f(-x)=f(x)$ the function is \cle{even} (symmetric about the $y$-axis); if $f(-x)=-f(x)$ it is \cle{odd} (symmetric about the origin). You then only need to study $x\ge 0$.
\item \textbf{Limits and asymptotes.} Evaluate $\lim_{x\to\pm\infty} f(x)$ and the limits at any excluded point. This reveals horizontal and vertical asymptotes.
\item \textbf{Monotonicity.} Compute $f'(x)$ using $\sinh'=\cosh$ and $\cosh'=\sinh$, study its sign, and build the variation table. Read off intercepts and extrema.
\end{enumerate}
\end{methode}

\begin{propriete}[Parity of the basic hyperbolic functions]
$\sinh$ is an \textbf{odd} function, $\cosh$ is an \textbf{even} function, and $\tanh$ is \textbf{odd}. Consequently $\operatorname{csch}$ and $\coth$ are odd, while $\operatorname{sech}$ is even.
\end{propriete}

\textbf{Proof (examinable).} Directly from the definitions:
\[
\sinh(-x)=\frac{e^{-x}-e^{x}}{2}=-\frac{e^{x}-e^{-x}}{2}=-\sinh x,
\qquad
\cosh(-x)=\frac{e^{-x}+e^{x}}{2}=\cosh x.
\]
Then $\tanh(-x)=\dfrac{\sinh(-x)}{\cosh(-x)}=-\tanh x$, and the parity of the reciprocals follows. \hfill$\blacksquare$

\courssub{The graph of $y=\cosh x$}

We apply the programme to $f(x)=\cosh x=\dfrac{e^x+e^{-x}}{2}$.

\textbf{Domain and symmetry.} $\cosh x$ is defined on $\mathbb{R}$ and is even, so we study $x\ge 0$ and reflect in the $y$-axis.

\textbf{Limits.} As $x\to+\infty$, $e^{-x}\to 0$, so $\cosh x \approx \dfrac{e^x}{2}\to +\infty$. There is no horizontal asymptote; the curve grows exponentially, hugging the curve $y=\tfrac12 e^{x}$. As $x\to-\infty$, by evenness $\cosh x\to+\infty$, asymptotic to $y=\tfrac12 e^{-x}$.

\textbf{Monotonicity.} $f'(x)=\sinh x$. Since $\sinh x>0$ for $x>0$ and $\sinh x<0$ for $x<0$, $\cosh$ is strictly decreasing on $(-\infty,0]$ and strictly increasing on $[0,+\infty)$, with a minimum at $x=0$:
\[
\cosh 0 = \frac{1+1}{2}=1.
\]

\begin{center}
\begin{tabular}{|c|ccccc|}
\hline
$x$ & $-\infty$ & & $0$ & & $+\infty$ \\
\hline
$f'(x)=\sinh x$ & $-$ & $-$ & $0$ & $+$ & $+$ \\
\hline
$f(x)=\cosh x$ & $+\infty$ & $\searrow$ & $1$ & $\nearrow$ & $+\infty$ \\
\hline
\end{tabular}
\end{center}

The curve is a smooth ``U''-shape with minimum at $(0,1)$, never touching the $x$-axis (since $\cosh x \ge 1$ for all $x$).

\begin{popfigure}
\begin{center}
\begin{tikzpicture}[scale=1.05]
\draw[->, thick, popline] (-3.4,0) -- (3.4,0) node[below left] {$x$};
\draw[->, thick, popline] (0,-0.4) -- (0,4.4) node[below left] {$y$};
\draw[very thick, popblue, domain=-2.1:2.1, samples=80] plot (\x, {(exp(\x)+exp(-\x))/2});
\draw[dashed, popmuted, domain=-2.1:2.1, samples=60] plot (\x, {0.5*exp(\x)});
\draw[dashed, popmuted, domain=-2.1:2.1, samples=60] plot (\x, {0.5*exp(-\x)});
\fill[poporange] (0,1) circle (2pt);
\node[poporange, below right, align=center] at (0.05,0.95) {minimum $(0,1)$};
\node[popblue] at (2.15,3.6) {$y=\cosh x$};
\node[popmuted] at (2.5,1.7) {$y=\tfrac12 e^{x}$};
\node[popmuted] at (-2.5,1.7) {$y=\tfrac12 e^{-x}$};
\foreach \t in {-2,-1,1,2} \draw (\t,0.06) -- (\t,-0.06) node[below] {\scriptsize $\t$};
\draw (0.06,1) -- (-0.06,1) node[left] {\scriptsize $1$};
\end{tikzpicture}
\end{center}
\legende{The curve $y=\cosh x$: even, minimum at $(0,1)$, asymptotic to $y=\tfrac12 e^{x}$ as $x\to+\infty$ and to $y=\tfrac12 e^{-x}$ as $x\to-\infty$.}
\end{popfigure}

\begin{savaistu}[The hanging chain]
The curve $y=\cosh x$ is called a \cle{catenary} (from the Latin \emph{catena}, chain). A flexible chain or cable hanging freely under its own weight --- like an electricity line between two pylons on the road from Yaound\'e to Douala --- takes exactly the shape of a catenary, not a parabola as one might guess. Engineers use $\cosh$ to compute the tension and the sag of such cables.
\end{savaistu}

\courssub{The graph of $y=\sinh x$}

\textbf{Domain and symmetry.} Defined on $\mathbb{R}$, odd: study $x\ge 0$ and use origin symmetry.

\textbf{Limits.} As $x\to+\infty$, $\sinh x\approx \tfrac12 e^{x}\to+\infty$; as $x\to-\infty$, $\sinh x\to-\infty$.

\textbf{Monotonicity.} $f'(x)=\cosh x \ge 1 >0$ for all $x$: $\sinh$ is \textbf{strictly increasing on} $\mathbb{R}$. It vanishes only at $x=0$ (the only intercept), and the tangent at the origin has gradient $\cosh 0 = 1$, i.e.\ the curve leaves the origin along the line $y=x$.

\begin{center}
\begin{tabular}{|c|ccccc|}
\hline
$x$ & $-\infty$ & & $0$ & & $+\infty$ \\
\hline
$f'(x)=\cosh x$ & $+$ & $+$ & $1$ & $+$ & $+$ \\
\hline
$f(x)=\sinh x$ & $-\infty$ & $\nearrow$ & $0$ & $\nearrow$ & $+\infty$ \\
\hline
\end{tabular}
\end{center}

\begin{attention}[Do not confuse the shapes]
$\sinh x$ has \emph{no} minimum or maximum --- it climbs steadily from $-\infty$ to $+\infty$. Only $\cosh x$ has a turning point. Also, $\sinh x$ is negative for $x<0$; never sketch it above the $x$-axis on the left.
\end{attention}

\courssub{The graph of $y=\tanh x$}

\textbf{Domain and symmetry.} Defined on $\mathbb{R}$ (since $\cosh x \ge 1 \neq 0$), odd.

\textbf{Limits --- the key feature.} Write
\[
\tanh x = \frac{e^{x}-e^{-x}}{e^{x}+e^{-x}} = \frac{1-e^{-2x}}{1+e^{-2x}}.
\]
As $x\to+\infty$, $e^{-2x}\to 0$, so $\tanh x \to 1$. As $x\to-\infty$, multiplying numerator and denominator by $e^{2x}$ gives $\tanh x = \dfrac{e^{2x}-1}{e^{2x}+1}\to -1$.

\begin{propriete}[Asymptotes of $\tanh$]
The curve $y=\tanh x$ has two \textbf{horizontal asymptotes}: $y=1$ as $x\to+\infty$ and $y=-1$ as $x\to-\infty$. Moreover $-1<\tanh x<1$ for all real $x$.
\end{propriete}

\textbf{Monotonicity.} Using the quotient rule,
\[
\frac{d}{dx}\tanh x = \frac{\cosh^2 x - \sinh^2 x}{\cosh^2 x}=\frac{1}{\cosh^2 x}=\operatorname{sech}^2 x >0,
\]
so $\tanh$ is strictly increasing on $\mathbb{R}$, passing through the origin with gradient $1$, and squeezing between its two asymptotes. The graph is an ``S''-shaped (sigmoid) curve.

\begin{center}
\begin{tabular}{|c|ccccc|}
\hline
$x$ & $-\infty$ & & $0$ & & $+\infty$ \\
\hline
$f'(x)=\operatorname{sech}^2 x$ & $+$ & $+$ & $1$ & $+$ & $+$ \\
\hline
$f(x)=\tanh x$ & $-1$ & $\nearrow$ & $0$ & $\nearrow$ & $1$ \\
\hline
\end{tabular}
\end{center}

\begin{methode}[Finding horizontal asymptotes of a hyperbolic quotient]
\begin{enumerate}
\item Express the function in terms of $e^{x}$ and $e^{-x}$.
\item For $x\to+\infty$, divide numerator and denominator by the dominant exponential ($e^{x}$ or $e^{2x}$) and let the vanishing terms go to $0$.
\item For $x\to-\infty$, divide by the dominant exponential in that direction ($e^{-x}$ or $e^{-2x}$).
\item Conclude: $y=L$ is a horizontal asymptote whenever $\lim f(x)=L$ (finite).
\end{enumerate}
\end{methode}

\courssub{The reciprocal functions: $\operatorname{sech}$, $\operatorname{csch}$, $\coth$}

The remaining three graphs are obtained as \textbf{reciprocals} of the first three, using the standard rules: where $f$ is large, $1/f$ is small; where $f$ is small (nonzero), $1/f$ is large; where $f=0$, $1/f$ has a vertical asymptote.

\textbf{The graph of $y=\operatorname{sech} x = \dfrac{1}{\cosh x}$.} Since $\cosh x \ge 1$, we have $0<\operatorname{sech} x \le 1$. It is even, has a \emph{maximum} at $(0,1)$, and $\operatorname{sech} x \to 0$ as $x\to\pm\infty$: the $x$-axis is a horizontal asymptote. The curve is a symmetric ``bell-shaped'' bump. Its derivative is $-\operatorname{sech} x \tanh x$, negative for $x>0$: decreasing on $[0,+\infty)$.

\begin{center}
\begin{tabular}{|c|ccccc|}
\hline
$x$ & $-\infty$ & & $0$ & & $+\infty$ \\
\hline
$f'(x)=-\operatorname{sech} x \tanh x$ & $+$ & $+$ & $0$ & $-$ & $-$ \\
\hline
$f(x)=\operatorname{sech} x$ & $0$ & $\nearrow$ & $1$ & $\searrow$ & $0$ \\
\hline
\end{tabular}
\end{center}

\textbf{The graph of $y=\operatorname{csch} x = \dfrac{1}{\sinh x}$.} Since $\sinh 0 = 0$, $\operatorname{csch}$ is undefined at $x=0$: the domain is $\mathbb{R}\setminus\{0\}$ and the line $x=0$ is a \textbf{vertical asymptote}, with
\[
\lim_{x\to 0^+}\operatorname{csch} x = +\infty, \qquad \lim_{x\to 0^-}\operatorname{csch} x = -\infty.
\]
It is odd, strictly decreasing on each of $(-\infty,0)$ and $(0,+\infty)$, and $\operatorname{csch} x \to 0$ as $x\to\pm\infty$ (horizontal asymptote $y=0$). Derivative: $-\operatorname{csch} x \coth x < 0$ for $x \neq 0$.

\begin{center}
\begin{tabular}{|c|ccccccc|}
\hline
$x$ & $-\infty$ & & $0^-$ & $0^+$ & & $+\infty$ \\
\hline
$f'(x)=-\operatorname{csch} x \coth x$ & $-$ & $-$ & $-\infty$ & $-\infty$ & $-$ & $-$ \\
\hline
$f(x)=\operatorname{csch} x$ & $0$ & $\searrow$ & $-\infty$ & $+\infty$ & $\searrow$ & $0$ \\
\hline
\end{tabular}
\end{center}

\textbf{The graph of $y=\coth x = \dfrac{\cosh x}{\sinh x}$.} Again undefined at $x=0$ (vertical asymptote $x=0$). It is odd. For the limits at infinity:
\[
\coth x = \frac{1+e^{-2x}}{1-e^{-2x}} \to 1 \ \text{as } x\to+\infty, \qquad \coth x \to -1 \ \text{as } x\to-\infty,
\]
so $y=1$ and $y=-1$ are horizontal asymptotes, with $|\coth x|>1$ everywhere. Its derivative is $-\operatorname{csch}^2 x <0$: strictly decreasing on each branch. Note carefully: unlike $\tanh$, the curve $\coth$ lies \emph{outside} the strip $-1<y<1$.

\begin{center}
\begin{tabular}{|c|ccccccc|}
\hline
$x$ & $-\infty$ & & $0^-$ & $0^+$ & & $+\infty$ \\
\hline
$f'(x)=-\operatorname{csch}^2 x$ & $-$ & $-$ & $-\infty$ & $-\infty$ & $-$ & $-$ \\
\hline
$f(x)=\coth x$ & $-1$ & $\searrow$ & $-\infty$ & $+\infty$ & $\searrow$ & $1$ \\
\hline
\end{tabular}
\end{center}

\begin{attention}[Vertical asymptotes and open circles]
$\coth x$ and $\operatorname{csch} x$ are \textbf{not defined at $x=0$}. Your sketch must show the vertical asymptote $x=0$ (dashed), and the two branches must never touch or cross the $y$-axis. If you mark a value at a point where the function is undefined, use an \textbf{open circle}. A common examination error is to draw $\coth$ as a single continuous S-curve through the origin --- that is the graph of $\tanh$, not $\coth$!
\end{attention}

\begin{popfigure}
\begin{center}
\begin{tikzpicture}
\begin{axis}[
  width=14cm, height=9.5cm,
  axis lines=middle,
  xlabel={$x$}, ylabel={$y$},
  xmin=-3.2, xmax=3.2, ymin=-3.2, ymax=3.2,
  xtick={-3,-2,-1,1,2,3}, ytick={-3,-2,-1,1,2,3},
  domain=-3:3, samples=120,
  legend style={at={(1.02,0.5)}, anchor=west, font=\scriptsize},
  clip=true
]
\addplot[very thick, popblue] {(exp(x)-exp(-x))/2};
\addlegendentry{$\sinh x$}
\addplot[very thick, poporange] {(exp(x)+exp(-x))/2};
\addlegendentry{$\cosh x$}
\addplot[very thick, popgreen] {tanh(x)};
\addlegendentry{$\tanh x$}
\addplot[very thick, poppurple, domain=-3:-0.35] {cosh(x)/sinh(x)};
\addplot[very thick, poppurple, domain=0.35:3] {cosh(x)/sinh(x)};
\addlegendentry{$\coth x$}
\addplot[very thick, poppink] {1/cosh(x)};
\addlegendentry{$\operatorname{sech} x$}
\addplot[very thick, popgold, domain=-3:-0.4] {1/sinh(x)};
\addplot[very thick, popgold, domain=0.4:3] {1/sinh(x)};
\addlegendentry{$\operatorname{csch} x$}
\addplot[dashed, popmuted, domain=-3:3] {1};
\addplot[dashed, popmuted, domain=-3:3] {-1};
\end{axis}
\end{tikzpicture}
\end{center}
\legende{All six hyperbolic functions on a common axis. Note the horizontal asymptotes $y=\pm1$ for $\tanh$ and $\coth$, and the vertical asymptote $x=0$ for $\coth$ and $\operatorname{csch}$.}
\end{popfigure}

\courssub{Comparison with trigonometric functions}

The names ``hyperbolic sine'', ``hyperbolic cosine'' invite a comparison with $\sin$ and $\cos$. The similarities are striking in the \emph{algebra} (identities, derivatives), but the \emph{graphs} are completely different.

\begin{center}
\begin{tabularx}{\linewidth}{|l|X|X|}
\hline
\rowcolor{popblueL}
\textbf{Feature} & \textbf{Trigonometric} & \textbf{Hyperbolic} \\
\hline
Bounded? & $-1\le \sin x,\cos x \le 1$ & $\cosh x \ge 1$; $\sinh x$ unbounded \\
\hline
Periodic? & Yes, period $2\pi$ & No period; monotone or single-turned \\
\hline
Zeros & Infinitely many & $\sinh x=0$ only at $x=0$; $\cosh x$ never zero \\
\hline
Asymptotes & None & $\tanh,\coth$: $y=\pm1$; $\operatorname{csch},\coth$: $x=0$ \\
\hline
Behaviour at infinity & Oscillates forever & Exponential growth ($\sinh,\cosh$) or tends to a limit \\
\hline
Parity & $\sin$ odd, $\cos$ even & $\sinh$ odd, $\cosh$ even \\
\hline
\end{tabularx}
\end{center}

The deep reason for the parallel \emph{names} is Euler's formula: $\cos x = \dfrac{e^{ix}+e^{-ix}}{2}$ and $\sin x = \dfrac{e^{ix}-e^{-ix}}{2i}$ --- the same exponential building blocks as $\cosh$ and $\sinh$, but with imaginary arguments. Replacing $x$ by $ix$ converts oscillation into exponential growth.

\courssub{Transformations of hyperbolic graphs}

All the transformation rules you know from pure mathematics apply unchanged. For $y = a\,f(b(x-h))+k$:
\begin{itemize}
\item $h$ is a \textbf{horizontal shift} (translation $\binom{h}{0}$);
\item $k$ is a \textbf{vertical shift};
\item $a$ is a \textbf{vertical stretch} (scale factor $a$, reflection in the $x$-axis if $a<0$);
\item $b$ is a \textbf{horizontal stretch} of scale factor $\tfrac{1}{b}$ (reflection in the $y$-axis if $b<0$).
\end{itemize}
Track the \emph{key features} --- minimum point, asymptotes, intercepts --- through the transformation.

\begin{exoresolu}[Sketching a transformed hyperbolic graph]
Sketch the curve $y = 2\cosh(x-1) - 1$, stating the coordinates of the minimum point and the range.
\tcblower
\textbf{Solution.} Start from $y=\cosh x$ (minimum at $(0,1)$, range $[1,+\infty)$) and apply the transformations in order.

\textbf{Step 1: horizontal shift.} $y=\cosh(x-1)$ is a translation by $\binom{1}{0}$: the minimum moves to $(1,1)$.

\textbf{Step 2: vertical stretch by factor 2.} $y=2\cosh(x-1)$: the minimum moves to $(1,2)$; the range becomes $[2,+\infty)$.

\textbf{Step 3: vertical shift down by 1.} $y=2\cosh(x-1)-1$: the minimum settles at $(1,1)$.

\textbf{Conclusion.} The curve is a catenary with \textbf{minimum point $(1,1)$} and \textbf{range} $y \ge 1$, i.e.\ $[1,+\infty)$. It is symmetric about the vertical line $x=1$, and $y\to+\infty$ as $x\to\pm\infty$. The $y$-intercept is $y = 2\cosh(-1)-1 = 2\cosh 1 - 1 \approx 2(1.5431)-1 = 2.086$.
\end{exoresolu}

\begin{exoresolu}[Identifying asymptotes of a transformed curve]
Determine all asymptotes of $y = 3\tanh(2x) + 1$ and state its range.
\tcblower
\textbf{Solution.} The basic curve $y=\tanh x$ has horizontal asymptotes $y=\pm 1$.

\textbf{Step 1.} The factor $2$ in $\tanh(2x)$ is a horizontal stretch of factor $\tfrac12$: it compresses the S-curve horizontally but does \emph{not} move the horizontal asymptotes.

\textbf{Step 2.} The factor $3$ (vertical stretch) sends the asymptotes $y=\pm1$ to $y=\pm 3$.

\textbf{Step 3.} The shift $+1$ sends them to $y = 3+1=4$ and $y=-3+1=-2$.

\textbf{Conclusion.} Horizontal asymptotes: $y=4$ (as $x\to+\infty$) and $y=-2$ (as $x\to-\infty$). The range is the open interval $(-2,\,4)$, since $\tanh$ never attains its limiting values. The curve passes through $(0,1)$ since $\tanh 0 = 0$.
\end{exoresolu}

\begin{aretenir}[Key facts about the six graphs]
\begin{itemize}
\item $\cosh x$: even, minimum $(0,1)$, $\cosh x \ge 1$, grows like $\tfrac12 e^{|x|}$.
\item $\sinh x$: odd, strictly increasing on $\mathbb{R}$, only zero at $x=0$, tangent $y=x$ at the origin.
\item $\tanh x$: odd, increasing, asymptotes $y=\pm 1$, range $(-1,1)$.
\item $\operatorname{sech} x$: even, maximum $(0,1)$, asymptote $y=0$, range $(0,1]$.
\item $\operatorname{csch} x$: odd, vertical asymptote $x=0$, horizontal asymptote $y=0$.
\item $\coth x$: odd, vertical asymptote $x=0$, horizontal asymptotes $y=\pm1$, $|\coth x|>1$.
\item Derivatives $\sinh'=\cosh$, $\cosh'=\sinh$ give the monotonicity directly.
\end{itemize}
\end{aretenir}

\begin{exercice}[Graphs and transformations]
\begin{enumerate}
\item Using the derivative, show that $y=\operatorname{sech} x$ is decreasing on $[0,+\infty)$ and sketch its graph, marking the maximum and the asymptote.
\item Sketch $y = \sinh(x+2)$ and $y = -\cosh x + 2$ on separate axes, stating the key point of each curve.
\item Find the horizontal asymptotes and the range of $y = 5 - 2\tanh(3x)$.
\item On the same axes, sketch $y=\tanh x$ and $y=\coth x$ for $x>0$. Explain why the two curves get closer and closer as $x\to+\infty$.
\end{enumerate}
\end{exercice}

\begin{corrige}[Exercise 1]
\textbf{1.} $\dfrac{d}{dx}\operatorname{sech} x = -\operatorname{sech} x \tanh x$. For $x>0$, $\operatorname{sech} x>0$ and $\tanh x>0$, so the derivative is negative: $\operatorname{sech}$ is strictly decreasing on $[0,+\infty)$. Even function, maximum $(0,1)$, asymptote $y=0$.

\textbf{2.} $y=\sinh(x+2)$: translation $\binom{-2}{0}$ of $y=\sinh x$; passes through $(-2,0)$ with gradient $1$, strictly increasing. $y=-\cosh x + 2$: reflection in the $x$-axis then shift up $2$; the minimum $(0,1)$ becomes a \textbf{maximum} $(0,1)$; range $(-\infty,1]$.

\textbf{3.} Asymptotes of $\tanh$ are $\pm1$; after stretch $-2$ and shift $+5$: $y = 5-2 = 3$ and $y = 5+2 = 7$. Since the stretch factor is negative, $y\to 3$ as $x\to+\infty$ and $y\to 7$ as $x\to-\infty$. Range: $(3,7)$.

\textbf{4.} For $x>0$, $\coth x = \dfrac{1}{\tanh x}$ with $0<\tanh x<1$, so $\coth x > 1 > \tanh x$. Both tend to $1$ as $x\to+\infty$ (since $e^{-2x}\to 0$), so the two branches converge towards the common asymptote $y=1$, with $\tanh$ approaching from below and $\coth$ from above.
\end{corrige}

With the graphs firmly in hand, we are ready for the next section, where the geometric symmetries we have just observed will reappear as powerful \emph{algebraic identities} --- and where Osborn's rule will let us borrow identities from trigonometry almost for free.

\coursec{Hyperbolic Identities and Osborn's Rule}

In the previous section we plotted the curves of $\sinh x$, $\cosh x$ and $\tanh x$ and observed, for instance, that the graph of $\cosh$ always lies above the line $y=1$, while the graph of $\sinh$ passes through the origin. Those pictures strongly suggest that the hyperbolic functions are bound together by algebraic relations, just as the circular functions are. In this section we make those relations precise: we derive the \cle{addition formulas}, the \cle{double-angle} and \cle{half-angle} identities, and we discover a remarkable ``translation machine'' --- \cle{Osborn's rule} --- which converts almost any trigonometric identity into a hyperbolic one.

\courssub{The fundamental identity revisited}

Before deriving new formulas, recall the cornerstone identity established in Section 1:
\[
\cosh^{2}x-\sinh^{2}x=1.
\]
Every identity in this section flows, directly or indirectly, from the exponential definitions
\[
\cosh x=\frac{e^{x}+e^{-x}}{2},\qquad \sinh x=\frac{e^{x}-e^{-x}}{2}.
\]
The strategy is always the same: \emph{replace each hyperbolic function by its exponential form, expand, and reassemble}. Let us see this strategy at work on the most important case.

\courssub{Addition formulas}

\begin{propriete}[Addition formulas (proof examinable)]
For all real numbers $A$ and $B$:
\begin{align*}
\sinh(A+B)&=\sinh A\cosh B+\cosh A\sinh B,\\
\sinh(A-B)&=\sinh A\cosh B-\cosh A\sinh B,\\
\cosh(A+B)&=\cosh A\cosh B+\sinh A\sinh B,\\
\cosh(A-B)&=\cosh A\cosh B-\sinh A\sinh B,\\
\tanh(A+B)&=\frac{\tanh A+\tanh B}{1+\tanh A\tanh B},\qquad
\tanh(A-B)=\frac{\tanh A-\tanh B}{1-\tanh A\tanh B}.
\end{align*}
\end{propriete}

\textbf{Proof of $\sinh(A+B)$.} Start from the right-hand side and use the exponential definitions:
\begin{align*}
\sinh A\cosh B+\cosh A\sinh B
&=\frac{e^{A}-e^{-A}}{2}\cdot\frac{e^{B}+e^{-B}}{2}+\frac{e^{A}+e^{-A}}{2}\cdot\frac{e^{B}-e^{-B}}{2}\\[2pt]
&=\frac{1}{4}\Big[\big(e^{A+B}+e^{A-B}-e^{-A+B}-e^{-A-B}\big)\\
&\qquad\qquad+\big(e^{A+B}-e^{A-B}+e^{-A+B}-e^{-A-B}\big)\Big]\\[2pt]
&=\frac{1}{4}\Big[2e^{A+B}-2e^{-(A+B)}\Big]
=\frac{e^{A+B}-e^{-(A+B)}}{2}=\sinh(A+B). \qquad \blacksquare
\end{align*}
Notice how the ``cross terms'' $e^{A-B}$ and $e^{-A+B}$ cancel between the two products: this cancellation is exactly why the formula is so clean. The proof of $\cosh(A+B)$ is identical except that the cross terms \emph{add} instead of cancelling, giving $2e^{A+B}+2e^{-(A+B)}$ in the bracket, hence $\cosh(A+B)$. The formulas for $A-B$ follow by replacing $B$ with $-B$ and using parity ($\cosh$ even, $\sinh$ odd). Finally,
\[
\tanh(A+B)=\frac{\sinh(A+B)}{\cosh(A+B)}
=\frac{\sinh A\cosh B+\cosh A\sinh B}{\cosh A\cosh B+\sinh A\sinh B},
\]
and dividing numerator and denominator by $\cosh A\cosh B$ yields the stated formula for $\tanh(A+B)$. $\blacksquare$

\begin{popfigure}
\begin{tikzpicture}[
  stp/.style={draw=popblue, rounded corners=2pt, fill=popblueL, text width=10.6cm, align=left, inner sep=6pt, font=\small},
  arr/.style={-{Stealth[length=2.5mm]}, thick, poporange},
  lbl/.style={font=\footnotesize\itshape, text=popdark}]
\node[stp] (s1) at (0,0) {\textbf{Step 1 --- Write the exponentials:}\quad
$\sinh A\cosh B+\cosh A\sinh B=\dfrac{(e^{A}-e^{-A})(e^{B}+e^{-B})+(e^{A}+e^{-A})(e^{B}-e^{-B})}{4}$};
\node[stp] (s2) at (0,-1.9) {\textbf{Step 2 --- Expand the products:}\quad
$=\dfrac{1}{4}\big[(e^{A+B}+e^{A-B}-e^{-A+B}-e^{-A-B})+(e^{A+B}-e^{A-B}+e^{-A+B}-e^{-A-B})\big]$};
\node[stp] (s3) at (0,-3.8) {\textbf{Step 3 --- Cancel the cross terms:}\quad
$e^{A-B}$ and $e^{-A+B}$ disappear;\quad $=\dfrac{1}{4}\big[2e^{A+B}-2e^{-(A+B)}\big]$};
\node[stp] (s4) at (0,-5.7) {\textbf{Step 4 --- Recognise the definition:}\quad
$=\dfrac{e^{A+B}-e^{-(A+B)}}{2}=\sinh(A+B)$. \quad \textbf{Done!}};
\draw[arr] (s1.south) -- (s2.north) node[lbl, midway, right=2pt] {multiply out};
\draw[arr] (s2.south) -- (s3.north) node[lbl, midway, right=2pt] {simplify};
\draw[arr] (s3.south) -- (s4.north) node[lbl, midway, right=2pt] {factor $2$};
\end{tikzpicture}
\legende{Step-by-step derivation of $\sinh(A+B)$ from the exponential definitions.}
\end{popfigure}

\begin{exemplebox}[Using an addition formula]
Compute $\sinh(\ln 2+\ln 3)$ exactly.

\emph{Solution.} $\sinh(\ln 2+\ln 3)=\sinh(\ln 2)\cosh(\ln 3)+\cosh(\ln 2)\sinh(\ln 3)$. Now $\sinh(\ln 2)=\frac{2-\frac12}{2}=\frac34$, $\cosh(\ln 2)=\frac{2+\frac12}{2}=\frac54$, $\sinh(\ln 3)=\frac{3-\frac13}{2}=\frac43$, $\cosh(\ln 3)=\frac{3+\frac13}{2}=\frac53$. Hence
\[
\sinh(\ln 2+\ln 3)=\frac34\cdot\frac53+\frac54\cdot\frac43=\frac{15}{12}+\frac{20}{12}=\frac{35}{12}.
\]
Check: $\ln 2+\ln 3=\ln 6$, and $\sinh(\ln 6)=\frac{6-\frac16}{2}=\frac{35}{12}$. $\checkmark$
\end{exemplebox}

\courssub{Double-angle and half-angle identities}

Setting $A=B=x$ in the addition formulas immediately gives the double-angle identities.

\begin{propriete}[Double-angle identities]
For all real $x$:
\[
\sinh 2x=2\sinh x\cosh x,
\]
\[
\cosh 2x=\cosh^{2}x+\sinh^{2}x=2\cosh^{2}x-1=1+2\sinh^{2}x,
\]
\[
\tanh 2x=\frac{2\tanh x}{1+\tanh^{2}x}.
\]
\end{propriete}

\textbf{Proof.} $\sinh 2x=\sinh(x+x)=\sinh x\cosh x+\cosh x\sinh x=2\sinh x\cosh x$. Similarly $\cosh 2x=\cosh^{2}x+\sinh^{2}x$; then substituting $\sinh^{2}x=\cosh^{2}x-1$ gives $\cosh 2x=2\cosh^{2}x-1$, and substituting $\cosh^{2}x=1+\sinh^{2}x$ gives $\cosh 2x=1+2\sinh^{2}x$. The $\tanh$ formula comes from $\tanh(x+x)$. $\blacksquare$

Rearranging the last two forms of $\cosh 2x$ produces the \cle{half-angle} (or ``power-lowering'') identities, which will be precious for integration in Section 6:
\[
\cosh^{2}x=\frac{\cosh 2x+1}{2},\qquad \sinh^{2}x=\frac{\cosh 2x-1}{2}.
\]

\begin{attention}[Sign trap]
For circular functions, $\cos 2x=\cos^{2}x-\sin^{2}x$, with a \emph{minus}. For hyperbolic functions it is $\cosh 2x=\cosh^{2}x\,\mathbf{+}\,\sinh^{2}x$, with a \emph{plus}. Mixing these two signs is the single most common error at the GCE: always double-check against $\cosh^{2}x-\sinh^{2}x=1$.
\end{attention}

\begin{exoresolu}[Solving with a double-angle identity]
Solve the equation $\cosh 2x-3\cosh x+4=0$, giving answers in exact logarithmic form.
\tcblower
\textbf{Solution.} Use $\cosh 2x=2\cosh^{2}x-1$:
\[
2\cosh^{2}x-1-3\cosh x+4=0
\;\Longrightarrow\;
2\cosh^{2}x-3\cosh x+3=0.
\]
The discriminant is $\Delta = (-3)^{2}-4(2)(3)=9-24=-15<0$. Hence the quadratic has no real roots for $\cosh x$. Since $\cosh x \in \mathbb{R}$ for all real $x$, the original equation has \textbf{no real solution}.

\textbf{Remark.} A classic solvable variant is $\cosh 2x-3\cosh x+2=0$. Then $2\cosh^{2}x-3\cosh x+1=0$, i.e. $(2\cosh x-1)(\cosh x-1)=0$. Since $\cosh x\ge 1$, $\cosh x=\frac12$ is impossible, so $\cosh x=1$, giving $x=0$ as the unique solution.

\emph{Lesson:} after solving the quadratic in $\cosh x$, \textbf{always reject roots below $1$}.
\end{exoresolu}

\courssub{Osborn's rule}

Compare the identities we have met so far with their circular cousins:

\begin{center}
\begin{tabularx}{\linewidth}{|X|X|}
\hline
\rowcolor{popblueL} \textbf{Trigonometric identity} & \textbf{Hyperbolic identity}\\
\hline
$\cos^{2}x+\sin^{2}x=1$ & $\cosh^{2}x-\sinh^{2}x=1$\\
\hline
$\sin(A+B)=\sin A\cos B+\cos A\sin B$ & $\sinh(A+B)=\sinh A\cosh B+\cosh A\sinh B$\\
\hline
$\cos(A+B)=\cos A\cos B-\sin A\sin B$ & $\cosh(A+B)=\cosh A\cosh B+\sinh A\sinh B$\\
\hline
$\sin 2x=2\sin x\cos x$ & $\sinh 2x=2\sinh x\cosh x$\\
\hline
$\cos 2x=\cos^{2}x-\sin^{2}x$ & $\cosh 2x=\cosh^{2}x+\sinh^{2}x$\\
\hline
$1+\tan^{2}x=\sec^{2}x$ & $1-\tanh^{2}x=\operatorname{sech}^{2}x$\\
\hline
\end{tabularx}
\end{center}

The pattern is striking: \emph{every formula is identical, except that a sign changes whenever a product of two sines is involved}. This is no coincidence. Since $\cos x=\cosh(ix)$ and $\sin x=-i\sinh(ix)$ (from the exponential definitions), any product of two sines picks up a factor $(-i)^{2}=-1$ when translated to hyperbolic form. This observation is formalised as:

\begin{propriete}[Osborn's Rule]
To convert any trigonometric identity into the corresponding hyperbolic identity:
\begin{enumerate}
\item replace $\cos$ by $\cosh$ and $\sin$ by $\sinh$ (and similarly $\tan\to\tanh$, $\sec\to\operatorname{sech}$, $\operatorname{cosec}\to\operatorname{cosech}$, $\cot\to\coth$);
\item change the sign of every term containing a product (or implied product) of two sines, i.e. replace $\sin^{2}$ by $-\sinh^{2}$.
\end{enumerate}
The rule is a consequence of $\cosh(ix)=\cos x$ and $\sinh(ix)=i\sin x$; its justification is not examinable, but \emph{using} the rule is.
\end{propriete}

\begin{methode}[Converting a trigonometric identity]
\begin{enumerate}
\item Write the identity in terms of $\sin$ and $\cos$ only (so that hidden products of sines become visible: e.g. $\tan^{2}x=\sin^{2}x/\cos^{2}x$).
\item Replace $\sin\to\sinh$, $\cos\to\cosh$.
\item Change the sign of every product of two sines ($\sin^{2}$, $\sin A\sin B$, a $\sin^{2}$ hidden in a numerator, etc.).
\item Verify the result on a numerical value (e.g. $x=0$ or $x=\ln 2$) as a safety check.
\end{enumerate}
\end{methode}

\begin{exemplebox}[Osborn's rule in action]
From $\cos 2x=1-2\sin^{2}x$ (the sine product $-2\sin^{2}x$ changes sign), we get
\[
\cosh 2x=1+2\sinh^{2}x \quad\checkmark
\]
From $1+\tan^{2}x=\sec^{2}x$, written as $1+\dfrac{\sin^{2}x}{\cos^{2}x}=\dfrac{1}{\cos^{2}x}$, the numerator $\sin^{2}x$ changes sign:
\[
1-\tanh^{2}x=\operatorname{sech}^{2}x \quad\checkmark
\]
From $\sin 3x=3\sin x-4\sin^{3}x$: here $\sin^{3}x=\sin x\cdot\sin^{2}x$ contains a product of two sines, so its sign flips:
\[
\sinh 3x=3\sinh x+4\sinh^{3}x.
\]
\end{exemplebox}

\begin{attention}[Limits of Osborn's rule]
Osborn's rule does \textbf{not} apply blindly to identities involving \emph{odd powers of sine standing alone}. A single $\sin x$ becomes $\sinh x$ with no sign change, but a lone $\sin^{3}x$ must be read as $\sin x\cdot\sin^{2}x$ to see the hidden product of two sines. Also, the rule is designed for \emph{identities}, not for equations with specific numerical coefficients that are not identities, nor for statements involving calculus (derivatives follow their own sign rules: $\frac{d}{dx}\sinh x=\cosh x$, no minus sign!). When in doubt, go back to the exponential definitions.
\end{attention}

\courssub{Applying identities to simplify and solve}

\begin{exoresolu}[Simplifying an expression]
Show that $\dfrac{\cosh 2x+\sinh 2x}{\cosh x+\sinh x}=e^{x}$, and deduce a simplification of $(\cosh x+\sinh x)^{n}$.
\tcblower
\textbf{Solution.} Using the addition formulas with $A=B=x$, or directly the definitions:
\[
\cosh x+\sinh x=\frac{e^{x}+e^{-x}}{2}+\frac{e^{x}-e^{-x}}{2}=e^{x}.
\]
Likewise $\cosh 2x+\sinh 2x=e^{2x}$. Hence the quotient equals $\dfrac{e^{2x}}{e^{x}}=e^{x}$. $\blacksquare$

It follows immediately that
\[
(\cosh x+\sinh x)^{n}=(e^{x})^{n}=e^{nx}=\cosh(nx)+\sinh(nx),
\]
a neat analogue of De Moivre's theorem, sometimes called the \emph{hyperbolic De Moivre identity}.
\end{exoresolu}

\begin{exoresolu}[Proving an identity]
Prove that $\cosh 3x=4\cosh^{3}x-3\cosh x$.
\tcblower
\textbf{Solution.} Write $\cosh 3x=\cosh(2x+x)$ and apply the addition formula:
\begin{align*}
\cosh(2x+x)&=\cosh 2x\cosh x+\sinh 2x\sinh x\\
&=(2\cosh^{2}x-1)\cosh x+(2\sinh x\cosh x)\sinh x\\
&=2\cosh^{3}x-\cosh x+2\cosh x\sinh^{2}x\\
&=2\cosh^{3}x-\cosh x+2\cosh x(\cosh^{2}x-1)\\
&=4\cosh^{3}x-3\cosh x. \qquad \blacksquare
\end{align*}
Note that this matches, via Osborn's rule, the circular identity $\cos 3x=4\cos^{3}x-3\cos x$ --- no sine product appears, so no sign changes.
\end{exoresolu}

\begin{exercice}[Consolidation]
\begin{enumerate}
\item Using the exponential definitions, prove that $\cosh(A-B)=\cosh A\cosh B-\sinh A\sinh B$.
\item Use Osborn's rule to write down the hyperbolic analogue of $\tan(A-B)=\dfrac{\tan A-\tan B}{1+\tan A\tan B}$, then verify it numerically for $A=\ln 2$, $B=\ln 3$.
\item Prove that $\sinh 3x=3\sinh x+4\sinh^{3}x$.
\item Solve $2\sinh^{2}x-3\cosh x=0$. (Hint: use $\cosh^{2}x-\sinh^{2}x=1$ first.)
\end{enumerate}
\end{exercice}

\begin{corrige}[Exercise 1]
\textbf{1.} $\cosh A\cosh B-\sinh A\sinh B=\frac{1}{4}\big[(e^{A}+e^{-A})(e^{B}+e^{-B})-(e^{A}-e^{-A})(e^{B}-e^{-B})\big]=\frac{1}{4}\big[2e^{A-B}+2e^{-(A-B)}\big]=\cosh(A-B)$.

\textbf{2.} Osborn: $\tanh(A-B)=\dfrac{\tanh A-\tanh B}{1-\tanh A\tanh B}$ (the product $\tan A\tan B$ hides $\sin A\sin B$, so the sign flips). Numerically: $\tanh(\ln 2)=\frac{2-\frac12}{2+\frac12}=\frac35$, $\tanh(\ln 3)=\frac{3-\frac13}{3+\frac13}=\frac45$; RHS $=\dfrac{\frac35-\frac45}{1-\frac{12}{25}}=\dfrac{-\frac15}{\frac{13}{25}}=-\frac{5}{13}$; LHS $=\tanh(\ln\tfrac23)=\dfrac{\frac23-\frac32}{\frac23+\frac32}=-\frac{5}{13}$. $\checkmark$

\textbf{3.} $\sinh(2x+x)=\sinh 2x\cosh x+\cosh 2x\sinh x=2\sinh x\cosh^{2}x+(1+2\sinh^{2}x)\sinh x$; replacing $\cosh^{2}x=1+\sinh^{2}x$ gives $2\sinh x(1+\sinh^{2}x)+\sinh x+2\sinh^{3}x=3\sinh x+4\sinh^{3}x$.

\textbf{4.} $2(\cosh^{2}x-1)-3\cosh x=0\Rightarrow 2\cosh^{2}x-3\cosh x-2=0\Rightarrow(2\cosh x+1)(\cosh x-2)=0$. Since $\cosh x\ge1$, only $\cosh x=2$ survives, so $x=\pm\ln(2+\sqrt{3})$.
\end{corrige}

\begin{aretenir}[Key points of Section 3]
\begin{itemize}
\item Addition formulas mirror the trigonometric ones, \emph{except} $\cosh(A+B)=\cosh A\cosh B\,\mathbf{+}\,\sinh A\sinh B$.
\item $\cosh 2x=\cosh^{2}x+\sinh^{2}x=2\cosh^{2}x-1=1+2\sinh^{2}x$; \quad $\sinh 2x=2\sinh x\cosh x$.
\item Half-angle: $\cosh^{2}x=\frac{\cosh 2x+1}{2}$, $\sinh^{2}x=\frac{\cosh 2x-1}{2}$.
\item \cle{Osborn's rule}: convert $\sin\to\sinh$, $\cos\to\cosh$, and \textbf{change the sign of every product of two sines} (including those hidden in $\tan^{2}$, $\sin^{3}$, etc.).
\item When solving equations in $\cosh x$, always reject roots with $\cosh x<1$.
\end{itemize}
\end{aretenir}

\begin{savaistu}[Why ``hyperbolic''?]
Just as $(\cos t,\sin t)$ parametrizes the unit circle $x^{2}+y^{2}=1$, the point $(\cosh t,\sinh t)$ parametrizes the unit hyperbola $x^{2}-y^{2}=1$ --- this is precisely the identity $\cosh^{2}t-\sinh^{2}t=1$ at work. The parameter $t$ even has the same geometric meaning in both cases: twice the area of the shaded sector. Osborn's rule is the algebraic shadow of this deep parallel between circle and hyperbola.
\end{savaistu}

\coursec{Inverse Hyperbolic Functions: Definitions, Graphs, and Logarithmic Forms}

In the previous section we discovered that hyperbolic identities mirror trigonometric ones, with sign changes governed by Osborn's rule. We now ask a natural question: \emph{can we undo a hyperbolic function?} If $\sinh y = 3$, what is $y$? Questions like this arise constantly at the GCE Advanced Level --- when solving hyperbolic equations, when evaluating integrals such as $\displaystyle\int \frac{dx}{\sqrt{x^2+1}}$, and when working with the catenary. The answer is the family of \cle{inverse hyperbolic functions}, and we shall see that, remarkably, each of them can be written using nothing more than the natural logarithm.

\courssub{The Idea of an Inverse Hyperbolic Function}

Recall from pure mathematics that a function $f$ has an inverse $f^{-1}$ only if $f$ is \cle{one-to-one} (bijective) on the domain considered, and that the graph of $f^{-1}$ is the reflection of the graph of $f$ in the line $y = x$.

\begin{definition}[Inverse hyperbolic functions]
For a hyperbolic function $f$, the inverse function $f^{-1}$ is defined by
\[
y = f^{-1}(x) \iff x = f(y),
\]
with the domain of $f$ restricted where necessary so that $f$ is one-to-one. The six inverse hyperbolic functions are:
\begin{center}
\begin{tabularx}{\linewidth}{|l|X|X|X|}
\hline
\rowcolor{popblueL} \textbf{Function} & \textbf{Notation} & \textbf{Domain} & \textbf{Range} \\
\hline
Inverse hyperbolic sine & $y = \operatorname{arsinh} x$ & $\mathbb{R}$ & $\mathbb{R}$ \\
\hline
Inverse hyperbolic cosine & $y = \operatorname{arcosh} x$ & $[1, +\infty)$ & $[0, +\infty)$ \\
\hline
Inverse hyperbolic tangent & $y = \operatorname{artanh} x$ & $(-1, 1)$ & $\mathbb{R}$ \\
\hline
Inverse hyperbolic cotangent & $y = \operatorname{arcoth} x$ & $(-\infty,-1)\cup(1,+\infty)$ & $\mathbb{R}\setminus\{0\}$ \\
\hline
Inverse hyperbolic secant & $y = \operatorname{arsech} x$ & $(0, 1]$ & $[0, +\infty)$ \\
\hline
Inverse hyperbolic cosecant & $y = \operatorname{arcsch} x$ & $\mathbb{R}\setminus\{0\}$ & $\mathbb{R}\setminus\{0\}$ \\
\hline
\end{tabularx}
\end{center}
\end{definition}

\begin{attention}[Why $\operatorname{arcosh}$ needs a restricted domain]
The function $\cosh$ is \emph{not} one-to-one on $\mathbb{R}$: for example $\cosh(2) = \cosh(-2)$. Its graph is symmetric about the $y$-axis, so a horizontal line $y = k$ (with $k > 1$) cuts it twice. To define an inverse we restrict $\cosh$ to $[0, +\infty)$, where it is strictly increasing. This is exactly parallel to restricting $\cos$ to $[0, \pi]$ in order to define $\arccos$.
\end{attention}

\begin{propriete}[Domain and range of $\operatorname{arcosh}$]
$\operatorname{arcosh} x$ is defined only for $x \geq 1$, and its range is $[0, +\infty)$. In particular $\operatorname{arcosh} 1 = 0$ and $\operatorname{arcosh} x > 0$ for $x > 1$. (Not examinable as a proof; it follows directly from the restriction of $\cosh$ to $[0,+\infty)$.)
\end{propriete}

\courssub{Graphs as Reflections in $y = x$}

Because $y = \sinh x$ is strictly increasing on $\mathbb{R}$, its inverse $\operatorname{arsinh}$ exists on the whole of $\mathbb{R}$ and is also strictly increasing and odd. The graph of $y = \operatorname{arsinh} x$ is obtained by reflecting $y = \sinh x$ in the line $y = x$.

\begin{popfigure}
\begin{center}
\begin{tikzpicture}
\begin{axis}[
  width=11cm, height=8.5cm,
  axis lines=middle,
  xlabel={$x$}, ylabel={$y$},
  xmin=-3.5, xmax=3.5, ymin=-3.5, ymax=3.5,
  xtick={-3,-2,-1,1,2,3}, ytick={-3,-2,-1,1,2,3},
  legend style={at={(0.03,0.97)}, anchor=north west, font=\small},
  samples=200, domain=-3:3, grid=both, grid style={popline!40},
]
\addplot[popblue, very thick] {ln(x+sqrt(max(0,x^2+1)))};
\addlegendentry{$y=\operatorname{arsinh} x$}
\addplot[poporange, very thick, dashed] {ln(x+sqrt(max(0,x^2+1)))};
\addlegendentry{$y=\ln(x+\sqrt{x^2+1})$}
\end{axis}
\end{tikzpicture}
\end{center}
\legende{The curve $y=\operatorname{arsinh} x$ (solid) coincides exactly with its logarithmic form $y=\ln(x+\sqrt{x^2+1})$ (dashed), confirming the equivalence derived below.}
\end{popfigure}

\begin{popfigure}
\begin{center}
\begin{tikzpicture}
\begin{axis}[
  width=11cm, height=8.5cm,
  axis lines=middle,
  xlabel={$x$}, ylabel={$y$},
  xmin=-3.5, xmax=4.5, ymin=-1, ymax=4.5,
  xtick={-3,-2,-1,1,2,3,4}, ytick={1,2,3,4},
  legend style={at={(0.03,0.97)}, anchor=north west, font=\small},
  samples=200, grid=both, grid style={popline!40},
]
\addplot[popblue, very thick, domain=0:2.1] {cosh(x)};
\addlegendentry{$y=\cosh x\ (x\ge 0)$}
\addplot[poporange, very thick, domain=1:4.3] {ln(x+sqrt(max(0,x^2-1)))};
\addlegendentry{$y=\operatorname{arcosh} x$}
\addplot[gray, thin, domain=-1:4.4] {x};
\addlegendentry{$y=x$}
\end{axis}
\end{tikzpicture}
\end{center}
\legende{$y=\operatorname{arcosh} x$ is the mirror image in $y=x$ of the \emph{right half} of $y=\cosh x$. The left half of $\cosh$ has no inverse on the principal branch.}
\end{popfigure}

The remaining graphs are obtained in the same way:
\begin{itemize}
\item $y = \operatorname{artanh} x$: an increasing S-shaped curve on $(-1,1)$, with vertical asymptotes $x = \pm 1$, passing through the origin (odd function).
\item $y = \operatorname{arcoth} x$: defined outside $[-1,1]$, with vertical asymptotes $x = \pm 1$, tending to $0$ as $|x| \to \infty$.
\item $y = \operatorname{arsech} x$: defined on $(0,1]$, decreasing from $+\infty$ (as $x \to 0^+$) down to $\operatorname{arsech} 1 = 0$.
\item $y = \operatorname{arcsch} x$: defined on $\mathbb{R}\setminus\{0\}$, odd, with a vertical asymptote at $x = 0$ and tending to $0$ as $|x| \to \infty$.
\end{itemize}

\courssub{Logarithmic Equivalents of the Inverse Hyperbolic Functions}

Here is the beautiful surprise of this chapter: because hyperbolic functions are built from exponentials, their inverses can be written with logarithms. These formulae are \emph{examinable} --- you must be able to derive them.

\begin{methode}[Deriving a logarithmic form]
To find the logarithmic form of an inverse hyperbolic function:
\begin{enumerate}
\item Set $y = \operatorname{arsinh} x$ (for example), so that $x = \sinh y$.
\item Write $\sinh y$ in exponentials: $x = \dfrac{e^y - e^{-y}}{2}$.
\item Multiply through by $2e^y$ to obtain a \textbf{quadratic in $e^y$}.
\item Solve the quadratic for $e^y$ using the quadratic formula.
\item Reject any root that is impossible (remember $e^y > 0$), then take $\ln$ of both sides.
\end{enumerate}
\end{methode}

\begin{propriete}[Logarithmic forms --- proofs examinable]
\begin{align*}
\operatorname{arsinh} x &= \ln\!\left(x + \sqrt{x^2+1}\right), && x \in \mathbb{R},\\
\operatorname{arcosh} x &= \ln\!\left(x + \sqrt{x^2-1}\right), && x \geq 1,\\
\operatorname{artanh} x &= \tfrac{1}{2}\ln\!\left(\frac{1+x}{1-x}\right), && -1 < x < 1,\\
\operatorname{arcoth} x &= \tfrac{1}{2}\ln\!\left(\frac{x+1}{x-1}\right), && |x| > 1,\\
\operatorname{arsech} x &= \ln\!\left(\frac{1+\sqrt{1-x^2}}{x}\right), && 0 < x \leq 1,\\
\operatorname{arcsch} x &= \ln\!\left(\frac{1}{x} + \frac{\sqrt{1+x^2}}{|x|}\right), && x \neq 0.
\end{align*}
\end{propriete}

\begin{exemplebox}[Derivation of $\operatorname{arsinh} x = \ln(x+\sqrt{x^2+1})$]
Let $y = \operatorname{arsinh} x$, so $x = \sinh y = \dfrac{e^y - e^{-y}}{2}$. Multiplying by $2e^y$:
\[
2x\,e^y = e^{2y} - 1 \quad\Longrightarrow\quad e^{2y} - 2x\,e^y - 1 = 0.
\]
This is a quadratic in $e^y$. By the quadratic formula:
\[
e^y = \frac{2x \pm \sqrt{4x^2+4}}{2} = x \pm \sqrt{x^2+1}.
\]
Since $\sqrt{x^2+1} > |x| \geq x$, the root $x - \sqrt{x^2+1}$ is negative, which is impossible because $e^y > 0$. Hence $e^y = x + \sqrt{x^2+1}$, and taking logarithms:
\[
y = \operatorname{arsinh} x = \ln\!\left(x + \sqrt{x^2+1}\right). \qquad \blacksquare
\]
\end{exemplebox}

\begin{exemplebox}[Derivation of the logarithmic form of $\operatorname{artanh}$]
Let $y = \operatorname{artanh} x$, so $x = \tanh y = \dfrac{e^y - e^{-y}}{e^y + e^{-y}} = \dfrac{e^{2y}-1}{e^{2y}+1}$. Then
\[
x\bigl(e^{2y}+1\bigr) = e^{2y} - 1
\quad\Longrightarrow\quad
e^{2y}(1 - x) = 1 + x
\quad\Longrightarrow\quad
e^{2y} = \frac{1+x}{1-x}.
\]
Taking logarithms and halving:
\[
y = \operatorname{artanh} x = \frac{1}{2}\ln\!\left(\frac{1+x}{1-x}\right), \qquad -1 < x < 1. \qquad \blacksquare
\]
Note that the condition $|x|<1$ guarantees $\frac{1+x}{1-x} > 0$, so the logarithm exists.
\end{exemplebox}

\begin{exoresolu}[Logarithmic form of $\operatorname{arcosh}$]
Prove that for $x \geq 1$,
\[
\operatorname{arcosh} x = \ln\!\left(x + \sqrt{x^2-1}\right).
\]
\tcblower
\textbf{Solution.} Let $y = \operatorname{arcosh} x$ with $x \geq 1$, so that $y \geq 0$ and $x = \cosh y = \dfrac{e^y + e^{-y}}{2}$. Multiplying by $2e^y$:
\[
e^{2y} - 2x\,e^y + 1 = 0.
\]
Solving this quadratic in $e^y$:
\[
e^y = \frac{2x \pm \sqrt{4x^2 - 4}}{2} = x \pm \sqrt{x^2 - 1}.
\]
Both roots are positive when $x \geq 1$, and they are reciprocals of each other since
\[
\bigl(x+\sqrt{x^2-1}\bigr)\bigl(x-\sqrt{x^2-1}\bigr) = x^2 - (x^2-1) = 1.
\]
The root $x - \sqrt{x^2-1} \leq 1$ would give $y = \ln(x-\sqrt{x^2-1}) \leq 0$, i.e.\ the \emph{negative} branch. Since the principal value requires $y \geq 0$, we keep the larger root:
\[
e^y = x + \sqrt{x^2-1} \quad\Longrightarrow\quad \operatorname{arcosh} x = \ln\!\left(x + \sqrt{x^2-1}\right). \qquad \blacksquare
\]
\end{exoresolu}

\begin{attention}[Discard the negative branch for $\operatorname{arcosh}$]
When solving the quadratic for $\operatorname{arcosh}$, both roots $x \pm \sqrt{x^2-1}$ are positive, so it is tempting to keep both. But the principal value of $\operatorname{arcosh}$ must satisfy $y \geq 0$. The root $x - \sqrt{x^2-1}$ is at most $1$, giving $y \leq 0$: it corresponds to $-\operatorname{arcosh} x$. Always take the $+$ sign. The same care applies to $\operatorname{arsech}$.
\end{attention}

\begin{exemplebox}[Numerical evaluation]
Evaluate $\operatorname{arsinh} 2$ and $\operatorname{arcosh} 2$ exactly and to 4 decimal places.
\[
\operatorname{arsinh} 2 = \ln\!\left(2 + \sqrt{5}\right) \approx 1.4436,
\qquad
\operatorname{arcosh} 2 = \ln\!\left(2 + \sqrt{3}\right) \approx 1.3170.
\]
Check on your calculator: $\sinh(1.4436) \approx 2$ and $\cosh(1.3170) \approx 2$. $\checkmark$
\end{exemplebox}

\begin{aretenir}[Key points of this section]
\begin{itemize}
\item $y = f^{-1}(x) \iff x = f(y)$; the graph of $f^{-1}$ is the reflection of the graph of $f$ in $y = x$.
\item $\operatorname{arsinh}$ has domain and range $\mathbb{R}$; $\operatorname{arcosh}$ has domain $[1,\infty)$ and range $[0,\infty)$; $\operatorname{artanh}$ has domain $(-1,1)$ and range $\mathbb{R}$.
\item Every inverse hyperbolic function has a logarithmic form, obtained by solving a quadratic in $e^y$.
\item For $\operatorname{arcosh}$ and $\operatorname{arsech}$, keep only the branch giving the principal (non-negative) value.
\end{itemize}
\end{aretenir}

\begin{exercice}[Logarithmic forms and evaluation]
\begin{enumerate}
\item Starting from $x = \operatorname{sech} y$, prove that $\operatorname{arsech} x = \ln\!\left(\dfrac{1+\sqrt{1-x^2}}{x}\right)$ for $0 < x \leq 1$.
\item Show that $\operatorname{artanh} x + \operatorname{artanh}\!\left(\tfrac{1}{3}\right)$ with $x = \tfrac{1}{2}$ can be written as a single logarithm, and simplify.
\item Evaluate exactly: $\operatorname{arcosh}\!\left(\tfrac{5}{4}\right)$.
\end{enumerate}
\end{exercice}

\begin{corrige}[Exercise 4.1]
\textbf{1.} Let $y = \operatorname{arsech} x$, so $x = \operatorname{sech} y = \dfrac{2}{e^y+e^{-y}}$ with $y \geq 0$. Then $x\,e^{2y} - 2e^y + x = 0$, a quadratic in $e^y$:
\[
e^y = \frac{2 \pm \sqrt{4 - 4x^2}}{2x} = \frac{1 \pm \sqrt{1-x^2}}{x}.
\]
The principal value $y \geq 0$ requires $e^y \geq 1$, so we take the $+$ sign (the other root is its reciprocal, $\leq 1$). Hence $\operatorname{arsech} x = \ln\!\left(\dfrac{1+\sqrt{1-x^2}}{x}\right)$.

\textbf{2.} $\operatorname{artanh}\tfrac12 + \operatorname{artanh}\tfrac13 = \tfrac12\ln\!\left(\tfrac{1+1/2}{1-1/2}\right) + \tfrac12\ln\!\left(\tfrac{1+1/3}{1-1/3}\right) = \tfrac12\ln 3 + \tfrac12\ln 2 = \tfrac12\ln 6 = \ln\sqrt{6}$.

\textbf{3.} $\operatorname{arcosh}\!\left(\tfrac54\right) = \ln\!\left(\tfrac54 + \sqrt{\tfrac{25}{16}-1}\right) = \ln\!\left(\tfrac54 + \tfrac34\right) = \ln 2$.
\end{corrige}

\begin{savaistu}[Why ``area'' functions?]
In many textbooks the inverse hyperbolic functions are written $\operatorname{arsinh}$, $\operatorname{arcosh}$, etc., where ``ar'' stands for \emph{area}, not ``arc''. Indeed, while $\arcsin x$ measures an \emph{arc length} on the unit circle, $\operatorname{arsinh} x$ measures an \emph{area} bounded by the rectangular hyperbola $x^2 - y^2 = 1$ --- a deep geometric parallel between circular and hyperbolic functions that also explains why these functions appear when integrating $\dfrac{1}{\sqrt{x^2 \pm 1}}$, as we shall exploit in Section 7.
\end{savaistu}

We are now fully equipped with the inverse hyperbolic functions and their logarithmic disguises. In the next section we put them to work immediately, using them as the key tool for \emph{solving hyperbolic equations}.

\coursec{Solving Hyperbolic Equations}

In the previous section we established that every inverse hyperbolic function can be expressed as a logarithm: for instance $\arsinh x = \ln\!\left(x+\sqrt{x^{2}+1}\right)$ and $\arcosh x = \ln\!\left(x+\sqrt{x^{2}-1}\right)$ for $x\geq 1$. We now put these results to work. A \cle{hyperbolic equation} is an equation in which the unknown appears inside hyperbolic functions, such as
\[
3\sinh x + 4\cosh x = 6, \qquad 2\sinh^{2}x - 3\sinh x + 1 = 0, \qquad \arcosh x = \arsinh 2.
\]
Because $\sinh$ and $\cosh$ are built from exponentials, the whole machinery of solving such equations rests on one simple idea: \emph{replace the hyperbolic functions by their exponential definitions and reduce the problem to an ordinary algebraic equation} --- usually a quadratic in $e^{x}$. Alternatively, when the equation involves inverse hyperbolic functions, we convert them into logarithms and solve a logarithmic equation. In both cases we must remain vigilant: algebraic manipulations (squaring, taking logarithms, multiplying by expressions that could be zero) can introduce \cle{extraneous solutions} that do not satisfy the original equation.

\courssub{Equations of the Form $a\sinh x + b\cosh x = c$}

\textbf{Intuition.}\quad Suppose you are told that a certain combination of $\sinh x$ and $\cosh x$ equals a given number. Since $\sinh x = \dfrac{e^{x}-e^{-x}}{2}$ and $\cosh x = \dfrac{e^{x}+e^{-x}}{2}$, the left-hand side is really a combination of $e^{x}$ and $e^{-x}$. Multiplying through by $e^{x}$ clears the negative power and produces a \emph{quadratic equation in the unknown $u = e^{x}$}. Since $e^{x}>0$ for all real $x$, only \textbf{positive} roots of this quadratic give genuine solutions, via $x = \ln u$.

\begin{methode}[Exponential substitution for $a\sinh x + b\cosh x = c$]
\begin{enumerate}
\item Replace $\sinh x = \dfrac{e^{x}-e^{-x}}{2}$ and $\cosh x = \dfrac{e^{x}+e^{-x}}{2}$.
\item Group the terms in $e^{x}$ and $e^{-x}$: the equation becomes
\[
\frac{b+a}{2}\,e^{x} + \frac{b-a}{2}\,e^{-x} = c.
\]
\item Multiply both sides by $2e^{x}$ (which is never zero):
\[
(b+a)e^{2x} - 2c\,e^{x} + (b-a) = 0.
\]
\item Set $u = e^{x}$ and solve the quadratic $(b+a)u^{2} - 2cu + (b-a) = 0$.
\item Keep only the roots with $u>0$; the solutions are $x = \ln u$.
\item \textbf{Verify} each candidate by substituting back into the original equation.
\end{enumerate}
\end{methode}

\begin{exoresolu}[Solving $3\sinh x + 4\cosh x = 6$]
Solve the equation $3\sinh x + 4\cosh x = 6$, giving your answers in exact logarithmic form.
\tcblower
\textbf{Solution.} Write the exponential forms:
\[
3\cdot\frac{e^{x}-e^{-x}}{2} + 4\cdot\frac{e^{x}+e^{-x}}{2} = 6
\;\Longrightarrow\;
\frac{7}{2}e^{x} + \frac{1}{2}e^{-x} = 6.
\]
Multiply by $2e^{x}$:
\[
7e^{2x} - 12e^{x} + 1 = 0.
\]
Let $u = e^{x}$: $\;7u^{2} - 12u + 1 = 0$. By the quadratic formula,
\[
u = \frac{12 \pm \sqrt{144 - 28}}{14} = \frac{12 \pm \sqrt{116}}{14} = \frac{6 \pm \sqrt{29}}{7}.
\]
Both roots are positive (since $\sqrt{29}\approx 5.385 < 6$), so both are admissible:
\[
\boxed{\,x = \ln\!\left(\frac{6+\sqrt{29}}{7}\right) \quad\text{or}\quad x = \ln\!\left(\frac{6-\sqrt{29}}{7}\right).\,}
\]
\textbf{Verification.} For $u_1 = \frac{6+\sqrt{29}}{7}$, $e^{-x} = 1/u_1 = \frac{7}{6+\sqrt{29}} = \frac{7(6-\sqrt{29})}{36-29} = \frac{6-\sqrt{29}}{7}$. Then $\sinh x = \frac{u_1 - 1/u_1}{2} = \frac{\sqrt{29}}{7}$, $\cosh x = \frac{u_1 + 1/u_1}{2} = \frac{6}{7}$. LHS $= 3(\sqrt{29}/7) + 4(6/7) = (3\sqrt{29}+24)/7$. Wait, check: $3\sinh x + 4\cosh x = \frac{3}{2}(u_1-1/u_1) + \frac{4}{2}(u_1+1/u_1) = \frac{7}{2}u_1 + \frac{1}{2}(1/u_1)$. Since $u_1$ satisfies $7u_1^2 -12u_1 + 1 = 0$, we have $7u_1 + 1/u_1 = 12$, so LHS $= \frac{1}{2}(12) = 6$. \checkmark\ The same algebra holds for $u_2$.
\end{exoresolu}

\begin{popfigure}
\begin{tikzpicture}[font=\small, node distance=0.55cm,
  box/.style={draw=popblue, fill=popblueL, rounded corners, align=center, text width=4.6cm, inner sep=5pt},
  dec/.style={draw=poporange, fill=white, rounded corners, align=center, text width=3.4cm, inner sep=5pt},
  arr/.style={-{Stealth}, thick, popdark}]
\node[box] (start) {Start: $a\sinh x + b\cosh x = c$};
\node[box, below=of start] (exp) {Substitute $\sinh x=\frac{e^x-e^{-x}}{2}$, $\cosh x=\frac{e^x+e^{-x}}{2}$};
\node[box, below=of exp, align=center] (quad){Multiply by $2e^{x}$:\\ quadratic in $u=e^{x}$};
\node[dec, below=of quad] (disc) {Positive roots $u>0$?};
\node[box, below left=0.7cm and -1.2cm of disc] (yes) {$x=\ln u$ for each $u>0$};
\node[box, below right=0.7cm and -1.2cm of disc] (no) {No real solution};
\node[box, below=1.6cm of disc] (check) {Substitute back to verify};
\draw[arr] (start) -- (exp);
\draw[arr] (exp) -- (quad);
\draw[arr] (quad) -- (disc);
\draw[arr] (disc) -- node[left, popdark]{yes} (yes);
\draw[arr] (disc) -- node[right, popdark]{no} (no);
\draw[arr] (yes) |- (check);
\end{tikzpicture}
\legende{Systematic strategy for solving $a\sinh x + b\cosh x = c$ by exponential substitution.}
\end{popfigure}

\textbf{An alternative: the auxiliary-argument identity.}\quad Just as $a\sin\theta + b\cos\theta$ can be written as a single sine in trigonometry, a linear combination of $\sinh$ and $\cosh$ collapses into a \emph{single} hyperbolic function.

\begin{propriete}[Auxiliary-argument identity]
If $b > |a|$, then
\[
a\sinh x + b\cosh x = \sqrt{b^{2}-a^{2}}\;\cosh(x+\alpha), \qquad \text{where } \tanh\alpha = \frac{a}{b}.
\]
(When $a > |b|$, one instead obtains $\sqrt{a^{2}-b^{2}}\sinh(x+\alpha)$ with $\tanh\alpha = b/a$.)
\end{propriete}

\textbf{Derivation (instructive).}\quad Expand $R\cosh(x+\alpha) = R\cosh x\cosh\alpha + R\sinh x\sinh\alpha$. Matching with $a\sinh x + b\cosh x$ gives $R\sinh\alpha = a$ and $R\cosh\alpha = b$. Squaring and subtracting using $\cosh^{2}\alpha - \sinh^{2}\alpha = 1$ yields $R^{2}(\cosh^{2}\alpha - \sinh^{2}\alpha) = b^{2}-a^{2}$, hence $R = \sqrt{b^{2}-a^{2}}$ (positive root since $\cosh(x+\alpha)\ge 1$). Dividing gives $\tanh\alpha = a/b$. The condition $b>|a|$ guarantees $b^{2}-a^{2}>0$ and $|\tanh\alpha|<1$, so a real $\alpha$ exists. $\blacksquare$

\begin{exemplebox}[Same equation, quicker route]
For $3\sinh x + 4\cosh x = 6$: here $b=4>|a|=3$, $R=\sqrt{16-9}=\sqrt{7}$, $\alpha = \artanh(3/4)$. The equation becomes $\cosh(x+\alpha) = 6/\sqrt{7}$. Since $6/\sqrt{7}>1$,
\[
x + \alpha = \pm\arcosh\!\left(\tfrac{6}{\sqrt7}\right)
\;\Longrightarrow\;
x = -\artanh\!\left(\tfrac34\right) \pm \arcosh\!\left(\tfrac{6}{\sqrt7}\right),
\]
which agrees with the two logarithmic answers found above.
\end{exemplebox}

\begin{attention}[Existence condition]
The equation $a\sinh x + b\cosh x = c$ (with $b>|a|$) has a real solution only if $\dfrac{c}{\sqrt{b^{2}-a^{2}}} \geq 1$, because the range of $\cosh$ is $[1,+\infty)$. For instance $3\sinh x + 4\cosh x = 2$ would require $\cosh(x+\alpha) = 2/\sqrt{7} < 1$: \textbf{no real solution}. Always check this before applying $\arcosh$.
\end{attention}

\begin{popfigure}
\begin{tikzpicture}
\begin{axis}[width=12cm, height=6.5cm, axis lines=middle, xlabel=$x$, ylabel=$y$,
  xmin=-3.2, xmax=2.2, ymin=-2, ymax=9, domain=-3.1:2.1, samples=120,
  xtick={-3,-2,-1,1,2}, legend style={at={(0.02,0.97)}, anchor=north west, font=\small}]
\addplot[popblue, thick] {3*sinh(x) + 4*cosh(x)};
\addlegendentry{$y = 3\sinh x + 4\cosh x$}
\addplot[poporange, thick, dashed] {6};
\addlegendentry{$y = 6$}
\addplot[only marks, mark=*, popdark] coordinates {(0.583,6) (-2.029,6)};
\node[popdark, font=\small, anchor=south] at (axis cs:0.583,6.1) {$x\approx 0.583$};
\node[popdark, font=\small, anchor=south] at (axis cs:-2.029,6.1) {$x\approx -2.029$};
\end{axis}
\end{tikzpicture}
\legende{Graphical verification: the curve $y=3\sinh x+4\cosh x$ meets the line $y=6$ at exactly the two solutions found algebraically.}
\end{popfigure}

\courssub{Quadratic-Type Hyperbolic Equations}

Some equations are quadratic \emph{in a hyperbolic function itself}. The strategy is to make a substitution such as $y = \sinh x$, solve the quadratic, and then solve $\sinh x = y$ for each admissible value of $y$.

\begin{aretenir}[Admissibility of roots]
Recall the ranges: $\sinh x \in \mathbb{R}$ (all reals), $\cosh x \in [1,+\infty)$, $\tanh x \in (-1,1)$. After solving the quadratic, \textbf{reject any value outside the range} of the substituted function. In particular, $\cosh x = k$ has a real solution only when $k\geq 1$.
\end{aretenir}

\begin{exoresolu}[Solving $2\sinh^{2}x - 3\sinh x + 1 = 0$]
Solve $2\sinh^{2}x - 3\sinh x + 1 = 0$.
\tcblower
\textbf{Solution.} Let $y = \sinh x$. Then $2y^{2} - 3y + 1 = 0$, which factorises:
\[
(2y - 1)(y - 1) = 0 \;\Longrightarrow\; y = \tfrac12 \ \text{or}\ y = 1.
\]
Both values lie in the range of $\sinh$ (all reals), so both are admissible.
\begin{itemize}
\item $\sinh x = 1 \;\Rightarrow\; x = \arsinh 1 = \ln(1+\sqrt{2})$.
\item $\sinh x = \tfrac12 \;\Rightarrow\; x = \arsinh\!\left(\tfrac12\right) = \ln\!\left(\tfrac{1+\sqrt{5}}{2}\right)$.
\end{itemize}
\textbf{Check} ($\sinh x = 1$): $2(1)^{2}-3(1)+1 = 0$. \checkmark\quad The solution set is
\[
\boxed{\,x = \ln\!\left(\tfrac{1+\sqrt5}{2}\right) \ \text{or}\ x = \ln(1+\sqrt2).\,}
\]
\end{exoresolu}

\begin{exoresolu}[An equation mixing $\sinh$ and $\cosh$]
Solve $2\cosh^{2}x - 3\sinh x = 3$.
\tcblower
\textbf{Solution.} The identity $\cosh^{2}x = 1 + \sinh^{2}x$ converts everything to $\sinh$:
\[
2(1+\sinh^{2}x) - 3\sinh x = 3
\;\Longrightarrow\;
2\sinh^{2}x - 3\sinh x - 1 = 0.
\]
With $y = \sinh x$: $\;y = \dfrac{3 \pm \sqrt{9+8}}{4} = \dfrac{3\pm\sqrt{17}}{4}$. Both roots are real numbers, hence in the range of $\sinh$:
\[
x = \arsinh\!\left(\frac{3+\sqrt{17}}{4}\right),
\qquad
x = \arsinh\!\left(\frac{3-\sqrt{17}}{4}\right).
\]
\textbf{Check} (using the quadratic relation): If $y$ satisfies $2y^2-3y-1=0$, then $2y^2 = 3y+1$. LHS of original equation: $2\cosh^2 x - 3\sinh x = 2(1+y^2) - 3y = 2 + (3y+1) - 3y = 3$. \checkmark
\end{exoresolu}

\begin{attention}[The $\cosh x = k$ trap]
If a quadratic in $\cosh x$ yields $\cosh x = \tfrac12$, do \textbf{not} write $x = \arcosh(1/2)$: the logarithmic form $\ln\!\left(\tfrac12 + \sqrt{\tfrac14 - 1}\right)$ involves the square root of a negative number. Since $\cosh x \geq 1$ for all real $x$, that branch simply has \emph{no real solution}. Similarly, $\tanh x = k$ requires $|k|<1$.
\end{attention}

\courssub{Equations Involving Inverse Hyperbolic Functions}

When the unknown sits inside an \emph{inverse} hyperbolic function, we use the logarithmic forms of Section 4:
\[
\arsinh x = \ln\!\left(x+\sqrt{x^{2}+1}\right),\quad
\arcosh x = \ln\!\left(x+\sqrt{x^{2}-1}\right)\ (x\geq 1),\quad
\artanh x = \tfrac12\ln\!\left(\frac{1+x}{1-x}\right)\ (|x|<1).
\]

\begin{methode}[Equations with inverse hyperbolic functions]
\begin{enumerate}
\item Note the \textbf{domain}: for $\arcosh$, the argument must satisfy argument $\geq 1$; for $\artanh$, $|\text{argument}|<1$.
\item Either apply the direct hyperbolic function to both sides (e.g.\ take $\cosh$ of both sides), or convert to logarithmic form.
\item Solve the resulting algebraic equation.
\item \textbf{Check every candidate} in the original equation, especially after squaring.
\end{enumerate}
\end{methode}

\begin{exoresolu}[Solving $\arcosh x = \arsinh 2$]
Solve $\arcosh x = \arsinh 2$.
\tcblower
\textbf{Solution.} Take $\cosh$ of both sides (valid since $\cosh$ is one-to-one on $[0,\infty)$, and both sides are $\geq 0$):
\[
x = \cosh\!\left(\arsinh 2\right).
\]
Let $\theta = \arsinh 2$, so $\sinh\theta = 2$. Then $\cosh^{2}\theta = 1 + \sinh^{2}\theta = 5$, and since $\cosh\theta>0$, $\cosh\theta = \sqrt{5}$. Hence
\[
\boxed{x = \sqrt{5}.}
\]
\textbf{Check}: $\sqrt5 \geq 1$ (domain satisfied) and $\arcosh\sqrt5 = \ln(\sqrt5 + \sqrt{5-1}) = \ln(\sqrt5+2)$, while $\arsinh 2 = \ln(2+\sqrt5)$. Equal. \checkmark
\end{exoresolu}

\begin{exoresolu}[A logarithmic-form equation with extraneous roots]
Solve $\arcosh(2x) = \ln 3$.
\tcblower
\textbf{Solution.} Using the logarithmic form,
\[
\ln\!\left(2x + \sqrt{4x^{2}-1}\right) = \ln 3
\;\Longrightarrow\;
2x + \sqrt{4x^{2}-1} = 3.
\]
Isolate the radical: $\sqrt{4x^{2}-1} = 3 - 2x$. Squaring (a step that may introduce extraneous roots):
\[
4x^{2} - 1 = 9 - 12x + 4x^{2}
\;\Longrightarrow\;
12x = 10
\;\Longrightarrow\;
x = \tfrac56.
\]
\textbf{Check}: the domain requires $2x \geq 1$, i.e.\ $x \geq \tfrac12$: satisfied. Also the unsquared equation needs $3 - 2x \geq 0$: $3 - \tfrac53 = \tfrac43 \geq 0$. \checkmark\ Finally $\arcosh\!\left(\tfrac53\right) = \ln\!\left(\tfrac53 + \sqrt{\tfrac{25}{9}-1}\right) = \ln\!\left(\tfrac53+\tfrac43\right) = \ln 3$. \checkmark\ So $x = \tfrac56$ is the unique solution.

\medskip
\emph{Why the check matters:} had we obtained, say, $x = 2$ from the squared equation, then $3-2x = -1 < 0$ while $\sqrt{4x^{2}-1} \geq 0$ --- a contradiction, and the candidate would have to be rejected as \cle{extraneous}.
\end{exoresolu}

\begin{attention}[Squaring introduces ghosts]
The implication $A = B \Rightarrow A^{2} = B^{2}$ is not reversible: $A^{2}=B^{2}$ also holds when $A = -B$. After squaring an equation containing a square root, \textbf{always} substitute each candidate back into the \emph{original} equation. The same caution applies when exponentiating a logarithmic equation: check that every logarithm's argument is strictly positive.
\end{attention}

\begin{exercice}[Consolidation]
\begin{enumerate}
\item Solve $5\sinh x - 3\cosh x = 2$, giving exact answers.
\item Solve $6\cosh^{2}x - 7\cosh x + 2 = 0$.
\item Solve $\artanh x = \ln 2$.
\item Show that the equation $\sinh x + \cosh x = 3$ has exactly one real solution and find it in exact logarithmic form. \emph{(Hint: $\sinh x + \cosh x = e^{x}$.)}
\end{enumerate}
\end{exercice}

\begin{corrige}[Exercise 5.1]
\textbf{1.} Exponential form: $\frac{5-3}{2}e^{x} + \frac{-5-3}{2}e^{-x} = 2$, i.e.\ $e^{x} - 4e^{-x} = 2$. Multiply by $e^{x}$: $e^{2x} - 2e^{x} - 4 = 0$, so $u = e^{x} = 1 \pm \sqrt{5}$. Only $u = 1+\sqrt{5} > 0$ is admissible: $x = \ln(1+\sqrt5)$.

\textbf{2.} With $y = \cosh x$: $(3y-2)(2y-1)=0$, so $y = \tfrac23$ or $y = \tfrac12$. Both are $<1$: \textbf{no real solution}.

\textbf{3.} $\tfrac12\ln\!\left(\tfrac{1+x}{1-x}\right) = \ln 2 \Rightarrow \tfrac{1+x}{1-x} = 4 \Rightarrow 1+x = 4 - 4x \Rightarrow x = \tfrac35$. Check: $|x|<1$. \checkmark

\textbf{4.} $\sinh x + \cosh x = e^{x}$, so the equation becomes $e^{x} = 3$. Since the exponential function is strictly positive and bijective from $\mathbb{R}$ to $(0,+\infty)$, there is exactly one real solution: $x = \ln 3$.
\end{corrige}

\begin{aretenir}[Key points of Section 5]
\begin{itemize}
\item Master move: replace $\sinh x, \cosh x$ by exponentials, multiply by $e^{x}$, and solve a \textbf{quadratic in $u = e^{x}$}, keeping only roots $u>0$.
\item Auxiliary form: $a\sinh x + b\cosh x = \sqrt{b^{2}-a^{2}}\cosh(x+\alpha)$ with $\tanh\alpha = a/b$ (when $b>|a|$); a solution exists only if $\dfrac{c}{\sqrt{b^{2}-a^{2}}}\geq 1$.
\item Quadratic-type equations: substitute $y = \sinh x$ (or $\cosh x$), then respect the \textbf{ranges} ($\cosh x \geq 1$, $|\tanh x|<1$).
\item Inverse-function equations: convert to logarithms or apply the direct function to both sides, respecting domains.
\item After squaring or exponentiating, \textbf{verify every candidate} in the original equation to eliminate extraneous roots.
\end{itemize}
\end{aretenir}

With equation-solving secured, we are ready for the calculus of hyperbolic functions: in the next section we differentiate and integrate $\sinh$, $\cosh$, $\tanh$ and their inverses, and discover why these functions are tailor-made for evaluating integrals involving $\sqrt{x^{2}\pm a^{2}}$.

\coursec{Calculus of Hyperbolic Functions: Differentiation and Integration}

In Section 5 we learnt how to solve equations involving hyperbolic functions, often by returning to the exponential definitions or by using the logarithmic forms of the inverse functions. We now take the natural next step: since $\sinh x$ and $\cosh x$ are built from the differentiable functions $e^x$ and $e^{-x}$, they must themselves be differentiable, and their calculus turns out to be remarkably simple --- simpler, in fact, than the calculus of trigonometric functions, because the awkward minus signs disappear. This section develops the complete toolkit: differentiation of the six hyperbolic functions, differentiation of their inverses, and the corresponding integration formulas, including the standard integrals of $\operatorname{sech} x$ and $\operatorname{csch} x$ which regularly appear in the GCE Advanced Level examination.

\courssub{Differentiation of the Basic Hyperbolic Functions}

\textbf{Intuition.} Recall the definitions:
\[
\sinh x = \frac{e^x - e^{-x}}{2}, \qquad \cosh x = \frac{e^x + e^{-x}}{2}.
\]
Differentiating $e^x$ gives $e^x$, and differentiating $e^{-x}$ gives $-e^{-x}$. So differentiating $\sinh x$ flips the sign of the second exponential, turning the \emph{difference} into a \emph{sum} --- which is exactly $\cosh x$. The two functions differentiate into each other, just as $\sin$ and $\cos$ do, but without any sign change.

\begin{propriete}[Derivatives of $\sinh$ and $\cosh$ (proof examinable)]
For all real $x$:
\[
\frac{d}{dx}(\sinh x) = \cosh x, \qquad \frac{d}{dx}(\cosh x) = \sinh x.
\]
\end{propriete}

\textbf{Proof.} Using the exponential definitions and $\dfrac{d}{dx}(e^{-x}) = -e^{-x}$:
\[
\frac{d}{dx}(\sinh x) = \frac{d}{dx}\left(\frac{e^x - e^{-x}}{2}\right) = \frac{e^x + e^{-x}}{2} = \cosh x,
\]
\[
\frac{d}{dx}(\cosh x) = \frac{d}{dx}\left(\frac{e^x + e^{-x}}{2}\right) = \frac{e^x - e^{-x}}{2} = \sinh x. \qquad \blacksquare
\]

\begin{attention}[No minus sign!]
In trigonometry, $\frac{d}{dx}(\cos x) = -\sin x$. For hyperbolic functions, $\frac{d}{dx}(\cosh x) = +\sinh x$. This is the single most common sign error in this chapter. Remember: \cle{hyperbolic calculus has no minus signs} in the basic rules for $\sinh$, $\cosh$ and $\tanh$.
\end{attention}

\textbf{Derivatives of the other four functions.} Since $\tanh x = \dfrac{\sinh x}{\cosh x}$, the quotient rule gives
\[
\frac{d}{dx}(\tanh x) = \frac{\cosh x \cdot \cosh x - \sinh x \cdot \sinh x}{\cosh^2 x} = \frac{\cosh^2 x - \sinh^2 x}{\cosh^2 x} = \frac{1}{\cosh^2 x} = \operatorname{sech}^2 x,
\]
where we used the identity $\cosh^2 x - \sinh^2 x = 1$. The same technique applied to $\coth x = \dfrac{\cosh x}{\sinh x}$, $\operatorname{sech} x = \dfrac{1}{\cosh x}$ and $\operatorname{csch} x = \dfrac{1}{\sinh x}$ yields the complete table.

\begin{aretenir}[Derivatives of the six hyperbolic functions]
\begin{center}
\begin{tabular}{|c|c|}
\hline
\rowcolor{popblueL} $f(x)$ & $f'(x)$ \\
\hline
$\sinh x$ & $\cosh x$ \\
\hline
$\cosh x$ & $\sinh x$ \\
\hline
$\tanh x$ & $\operatorname{sech}^2 x$ \\
\hline
$\coth x$ & $-\operatorname{csch}^2 x$ \\
\hline
$\operatorname{sech} x$ & $-\operatorname{sech} x \tanh x$ \\
\hline
$\operatorname{csch} x$ & $-\operatorname{csch} x \coth x$ \\
\hline
\end{tabular}
\end{center}
Note the pattern: the three ``co''-functions ($\coth$, $\operatorname{sech}$, $\operatorname{csch}$) carry the minus signs, exactly as Osborn's rule would suggest from the trigonometric analogues.
\end{aretenir}

\begin{methode}[Differentiating composite hyperbolic functions]
\begin{enumerate}
\item Identify the outer function and the inner function $u = g(x)$.
\item Apply the chain rule: $\frac{d}{dx}f(g(x)) = f'(g(x)) \cdot g'(x)$.
\item Use the standard derivatives from the table above.
\item Simplify using hyperbolic identities where appropriate (e.g., $2\sinh u \cosh u = \sinh 2u$).
\end{enumerate}
\end{methode}

\begin{exoresolu}[Differentiating composite hyperbolic functions]
Differentiate with respect to $x$:
\quad (a) $y = \sinh(3x^2)$
\quad (b) $y = \cosh^2(2x)$
\quad (c) $y = e^x \tanh x$
\quad (d) $y = \operatorname{sech}(x^3)$.
\tcblower
\textbf{Solution.}
(a) By the chain rule with $u = 3x^2$:
\[
\frac{dy}{dx} = \cosh(3x^2) \cdot 6x = 6x\cosh(3x^2).
\]
(b) Write $y = [\cosh(2x)]^2$. Applying the chain rule twice:
\[
\frac{dy}{dx} = 2\cosh(2x) \cdot \sinh(2x) \cdot 2 = 4\sinh(2x)\cosh(2x) = 2\sinh(4x),
\]
using the double-angle identity $\sinh(2u) = 2\sinh u \cosh u$.
(c) By the product rule:
\[
\frac{dy}{dx} = e^x \tanh x + e^x \operatorname{sech}^2 x = e^x\left(\tanh x + \operatorname{sech}^2 x\right).
\]
(d) By the chain rule with $u = x^3$:
\[
\frac{dy}{dx} = -\operatorname{sech}(x^3)\tanh(x^3) \cdot 3x^2 = -3x^2 \operatorname{sech}(x^3)\tanh(x^3).
\]
\end{exoresolu}

\courssub{Differentiation of the Inverse Hyperbolic Functions}

\textbf{Strategy.} There are two standard approaches, and a good candidate masters both:
\begin{enumerate}
\item \emph{Implicit differentiation}: if $y = \operatorname{arsinh} x$, then $\sinh y = x$; differentiate both sides with respect to $x$ and solve for $\dfrac{dy}{dx}$.
\item \emph{Logarithmic forms}: differentiate $\operatorname{arsinh} x = \ln\left(x + \sqrt{x^2+1}\right)$ directly.
\end{enumerate}

\begin{propriete}[Derivative of $\operatorname{arsinh} x$ (proof examinable)]
For all real $x$:
\[
\frac{d}{dx}\left(\operatorname{arsinh} x\right) = \frac{1}{\sqrt{x^2 + 1}}.
\]
\end{propriete}

\textbf{Proof (implicit differentiation).} Let $y = \operatorname{arsinh} x$, so $\sinh y = x$. Differentiating with respect to $x$:
\[
\cosh y \, \frac{dy}{dx} = 1 \quad \Longrightarrow \quad \frac{dy}{dx} = \frac{1}{\cosh y}.
\]
Since $\cosh^2 y - \sinh^2 y = 1$ and $\cosh y > 0$ for all $y$,
\[
\cosh y = \sqrt{1 + \sinh^2 y} = \sqrt{1 + x^2}.
\]
Therefore $\dfrac{dy}{dx} = \dfrac{1}{\sqrt{x^2+1}}$. $\blacksquare$

\begin{propriete}[Derivatives of $\operatorname{arcosh} x$ and $\operatorname{artanh} x$]
\[
\frac{d}{dx}\left(\operatorname{arcosh} x\right) = \frac{1}{\sqrt{x^2 - 1}} \quad (x > 1), \qquad
\frac{d}{dx}\left(\operatorname{artanh} x\right) = \frac{1}{1 - x^2} \quad (|x| < 1).
\]
\end{propriete}

\textbf{Proof of $\frac{d}{dx}(\operatorname{artanh} x)$ (via the logarithmic form).} Recall $\operatorname{artanh} x = \dfrac{1}{2}\ln\dfrac{1+x}{1-x}$ for $|x| < 1$. Differentiating:
\[
\frac{d}{dx}\left(\operatorname{artanh} x\right) = \frac{1}{2}\left(\frac{1}{1+x} + \frac{1}{1-x}\right) = \frac{1}{2} \cdot \frac{(1-x) + (1+x)}{1-x^2} = \frac{1}{1-x^2}. \qquad \blacksquare
\]

\textbf{Proof of $\frac{d}{dx}(\operatorname{arcosh} x)$ (implicit differentiation).} Let $y = \operatorname{arcosh} x$, so $\cosh y = x$ with $y \ge 0$. Differentiating:
\[
\sinh y \, \frac{dy}{dx} = 1 \quad \Longrightarrow \quad \frac{dy}{dx} = \frac{1}{\sinh y}.
\]
Since $y \ge 0$, $\sinh y \ge 0$ and $\sinh y = \sqrt{\cosh^2 y - 1} = \sqrt{x^2 - 1}$. Hence $\dfrac{dy}{dx} = \dfrac{1}{\sqrt{x^2-1}}$ for $x > 1$. $\blacksquare$

\begin{attention}[Domain restrictions]
The formula $\frac{d}{dx}(\operatorname{arcosh} x) = \frac{1}{\sqrt{x^2-1}}$ is valid only for $x > 1$, because $\operatorname{arcosh} x$ is defined only for $x \geq 1$ (and is not differentiable at $x = 1$, where the tangent is vertical). Similarly $\operatorname{artanh} x$ requires $|x| < 1$. Always state the domain when quoting these results.
\end{attention}

\begin{methode}[Differentiating inverse hyperbolic functions]
\begin{enumerate}
\item Identify which inverse function is involved ($\operatorname{arsinh}$, $\operatorname{arcosh}$, $\operatorname{artanh}$).
\item Check the domain condition ($x \in \mathbb{R}$, $x>1$, $|x|<1$ respectively).
\item Apply the chain rule if the argument is a function $u(x)$: $\frac{d}{dx}f(u) = f'(u) \cdot u'(x)$.
\item Simplify the resulting expression.
\end{enumerate}
\end{methode}

The figure below summarises the essential differentiation results of this chapter.

\begin{popfigure}
\begin{tikzpicture}[
  box/.style={draw=popblue, thick, rounded corners=4pt, fill=popblueL, minimum width=5.6cm, minimum height=1.05cm, align=center},
  ibox/.style={draw=poppurple, thick, rounded corners=4pt, fill=white, minimum width=5.6cm, minimum height=1.05cm, align=center}]
\node[box] at (0,4.5) {$\dfrac{d}{dx}(\sinh x) = \cosh x$};
\node[box] at (0,3.1) {$\dfrac{d}{dx}(\cosh x) = \sinh x$};
\node[box] at (0,1.7) {$\dfrac{d}{dx}(\tanh x) = \operatorname{sech}^2 x$};
\node[ibox] at (7.2,4.5) {$\dfrac{d}{dx}(\operatorname{arsinh} x) = \dfrac{1}{\sqrt{x^2+1}}$};
\node[ibox] at (7.2,3.1) {$\dfrac{d}{dx}(\operatorname{arcosh} x) = \dfrac{1}{\sqrt{x^2-1}}\ (x>1)$};
\node[ibox] at (7.2,1.7) {$\dfrac{d}{dx}(\operatorname{artanh} x) = \dfrac{1}{1-x^2}\ (|x|<1)$};
\node[align=center, text=popblue] at (0,5.5) {\textbf{Direct functions}};
\node[align=center, text=poppurple] at (7.2,5.5) {\textbf{Inverse functions}};
\end{tikzpicture}
\legende{Summary table of the derivatives of the direct and inverse hyperbolic functions.}
\end{popfigure}

\begin{exoresolu}[Differentiating an inverse hyperbolic function]
Find $\dfrac{dy}{dx}$ given that $y = \operatorname{arcosh}(2x)$, $x > \dfrac{1}{2}$.
\tcblower
\textbf{Solution.} By the chain rule with $u = 2x$:
\[
\frac{dy}{dx} = \frac{1}{\sqrt{(2x)^2 - 1}} \cdot 2 = \frac{2}{\sqrt{4x^2 - 1}}.
\]
\emph{Check by implicit differentiation:} $\cosh y = 2x \Rightarrow \sinh y \, \frac{dy}{dx} = 2 \Rightarrow \frac{dy}{dx} = \frac{2}{\sinh y} = \frac{2}{\sqrt{\cosh^2 y - 1}} = \frac{2}{\sqrt{4x^2-1}}$. $\blacksquare$
\end{exoresolu}

\begin{exoresolu}[More differentiation of inverse hyperbolic functions]
Differentiate with respect to $x$:
\quad (a) $y = \operatorname{arsinh}(x^2)$
\quad (b) $y = \operatorname{artanh}(\sqrt{x})$, $0 < x < 1$
\quad (c) $y = \ln(\cosh 3x)$.
\tcblower
\textbf{Solution.}
(a) $\frac{dy}{dx} = \frac{1}{\sqrt{(x^2)^2+1}} \cdot 2x = \frac{2x}{\sqrt{x^4+1}}$.
(b) $\frac{dy}{dx} = \frac{1}{1-(\sqrt{x})^2} \cdot \frac{1}{2\sqrt{x}} = \frac{1}{2\sqrt{x}(1-x)}$.
(c) $\frac{dy}{dx} = \frac{1}{\cosh 3x} \cdot \sinh 3x \cdot 3 = 3\tanh 3x$.
\end{exoresolu}

\courssub{Integration of the Basic Hyperbolic Functions}

Every differentiation formula, read backwards, gives an integration formula. Since $\frac{d}{dx}(\sinh x) = \cosh x$ and $\frac{d}{dx}(\cosh x) = \sinh x$, we immediately obtain the following.

\begin{aretenir}[Standard integrals of hyperbolic functions]
\begin{center}
\begin{tabular}{|c|c|}
\hline
\rowcolor{popblueL} $f(x)$ & $\displaystyle\int f(x)\,dx$ \\
\hline
$\sinh x$ & $\cosh x + C$ \\
\hline
$\cosh x$ & $\sinh x + C$ \\
\hline
$\operatorname{sech}^2 x$ & $\tanh x + C$ \\
\hline
$\operatorname{csch}^2 x$ & $-\coth x + C$ \\
\hline
$\operatorname{sech} x \tanh x$ & $-\operatorname{sech} x + C$ \\
\hline
$\operatorname{csch} x \coth x$ & $-\operatorname{csch} x + C$ \\
\hline
\end{tabular}
\end{center}
Also useful: $\displaystyle\int \tanh x \, dx = \ln(\cosh x) + C$ and $\displaystyle\int \coth x \, dx = \ln|\sinh x| + C$, since the numerator is the derivative of the denominator in each case.
\end{aretenir}

\begin{exemplebox}[Direct integration]
\[
\int_0^1 \cosh(2x)\,dx = \left[\frac{1}{2}\sinh(2x)\right]_0^1 = \frac{1}{2}\sinh 2 - 0 = \frac{e^2 - e^{-2}}{4}.
\]
\end{exemplebox}

\textbf{The integral of $\operatorname{sech} x$.} This is a classic examination question because the trick is not obvious.

\begin{methode}[Integrating $\operatorname{sech} x$ (standard substitution $u = e^x$)]
\begin{enumerate}
\item Write $\operatorname{sech} x = \dfrac{2}{e^x + e^{-x}} = \dfrac{2e^x}{e^{2x}+1}$.
\item Substitute $u = e^x$, so $du = e^x\,dx$ and $dx = \dfrac{du}{u}$.
\item The integral becomes $\displaystyle\int \frac{2}{u^2+1}\,du = 2\arctan u + C$.
\item Back-substitute: $\displaystyle\int \operatorname{sech} x \, dx = 2\arctan(e^x) + C$.
\end{enumerate}
\end{methode}

\begin{propriete}[Standard integrals of $\operatorname{sech}$ and $\operatorname{csch}$]
\[
\int \operatorname{sech} x \, dx = 2\arctan(e^x) + C = \arctan(\sinh x) + C' = \arcsin(\tanh x) + C'',
\]
\[
\int \operatorname{csch} x \, dx = \ln\left|\tanh\frac{x}{2}\right| + C.
\]
\end{propriete}

\textbf{Derivation for $\operatorname{csch} x$.} Using $\operatorname{csch} x = \dfrac{2}{e^x - e^{-x}} = \dfrac{2e^x}{e^{2x}-1}$ and the substitution $u = e^x$:
\[
\int \operatorname{csch} x \, dx = \int \frac{2}{u^2 - 1}\,du = \ln\left|\frac{u-1}{u+1}\right| + C = \ln\left|\frac{e^x - 1}{e^x + 1}\right| + C = \ln\left|\tanh\frac{x}{2}\right| + C,
\]
since $\tanh\dfrac{x}{2} = \dfrac{e^{x/2} - e^{-x/2}}{e^{x/2} + e^{-x/2}} = \dfrac{e^x - 1}{e^x + 1}$. $\blacksquare$

\begin{attention}[$\operatorname{csch}$ versus $\csc$]
Do not confuse
\[
\int \operatorname{csch} x \, dx = \ln\left|\tanh\frac{x}{2}\right| + C
\quad\text{with}\quad
\int \csc x \, dx = \ln\left|\tan\frac{x}{2}\right| + C = \ln|\csc x - \cot x| + C.
\]
The hyperbolic integral involves $\tanh(x/2)$; the trigonometric one involves $\tan(x/2)$. Mixing them up costs marks every year.
\end{attention}

\begin{popfigure}
\begin{tikzpicture}
\begin{axis}[
  width=11cm, height=6.5cm,
  axis lines=middle,
  xlabel={$x$}, ylabel={$y$},
  xmin=-3.5, xmax=3.5, ymin=-0.2, ymax=1.4,
  xtick={-3,-2,-1,1,2,3}, ytick={1},
  samples=120, domain=-3.4:3.4,
  legend style={at={(0.98,0.95)}, anchor=north east, draw=popline}]
\addplot[draw=none, name path=axis] {0};
\addplot[popblue, very thick, name path=curve] {1/cosh(x)};
\addplot[popblueL] fill between[of=curve and axis];
\addlegendentry{$y = \operatorname{sech} x$}
\node[align=center, text=popdark] at (axis cs:1.9,0.55) {$\displaystyle\int \operatorname{sech} x\,dx = 2\arctan(e^x) + C$};
\end{axis}
\end{tikzpicture}
\legende{The curve $y = \operatorname{sech} x$; the shaded area under the curve is measured by $\int \operatorname{sech} x\,dx = 2\arctan(e^x) + C$.}
\end{popfigure}

\begin{exoresolu}[Integrating $\operatorname{sech} x$ and $\operatorname{csch} x$]
Evaluate: \quad (a) $\displaystyle\int_0^{\ln 3} \operatorname{sech} x \, dx$ \quad (b) $\displaystyle\int \operatorname{csch}(2x) \, dx$.
\tcblower
\textbf{Solution.}
(a) $\displaystyle\int_0^{\ln 3} \operatorname{sech} x \, dx = \Big[2\arctan(e^x)\Big]_0^{\ln 3} = 2\arctan(3) - 2\arctan(1) = 2\arctan 3 - \frac{\pi}{2}$.
(b) Let $u = 2x$, $du = 2\,dx$. Then $\displaystyle\int \operatorname{csch}(2x) \, dx = \frac{1}{2}\int \operatorname{csch} u \, du = \frac{1}{2}\ln\left|\tanh\frac{u}{2}\right| + C = \frac{1}{2}\ln\left|\tanh x\right| + C$.
\end{exoresolu}

\courssub{Chain Rule and Substitution in Integration}

Composite hyperbolic expressions are handled exactly as in ordinary integration: look for a function and its derivative, and substitute.

\begin{exoresolu}[Integration by substitution]
Evaluate:
\quad (a) $\displaystyle\int \sinh^3 x \cosh x \, dx$
\quad (b) $\displaystyle\int_0^{\ln 2} \frac{\sinh x}{\cosh^2 x}\,dx$
\quad (c) $\displaystyle\int x\cosh(x^2)\,dx$
\quad (d) $\displaystyle\int \tanh x \operatorname{sech}^2 x \, dx$.
\tcblower
\textbf{Solution.}
(a) Put $u = \sinh x$, so $du = \cosh x \, dx$:
\[
\int \sinh^3 x \cosh x \, dx = \int u^3\,du = \frac{u^4}{4} + C = \frac{1}{4}\sinh^4 x + C.
\]
(b) Put $u = \cosh x$, so $du = \sinh x\,dx$. When $x = 0$, $u = 1$; when $x = \ln 2$, $u = \cosh(\ln 2) = \dfrac{2 + \frac{1}{2}}{2} = \dfrac{5}{4}$:
\[
\int_0^{\ln 2} \frac{\sinh x}{\cosh^2 x}\,dx = \int_1^{5/4} \frac{du}{u^2} = \left[-\frac{1}{u}\right]_1^{5/4} = -\frac{4}{5} + 1 = \frac{1}{5}.
\]
(c) Put $u = x^2$, so $du = 2x\,dx$:
\[
\int x\cosh(x^2)\,dx = \frac{1}{2}\int \cosh u \, du = \frac{1}{2}\sinh(x^2) + C.
\]
(d) Put $u = \tanh x$, so $du = \operatorname{sech}^2 x\,dx$:
\[
\int \tanh x \operatorname{sech}^2 x \, dx = \int u\,du = \frac{u^2}{2} + C = \frac{1}{2}\tanh^2 x + C.
\]
\end{exoresolu}

\begin{exemplebox}[Using the inverse derivatives in reverse]
The results of Section 6.2 give three more standard integrals, extremely useful for the next section:
\[
\int \frac{dx}{\sqrt{x^2+1}} = \operatorname{arsinh} x + C, \qquad
\int \frac{dx}{\sqrt{x^2-1}} = \operatorname{arcosh} x + C \ (x>1), \qquad
\int \frac{dx}{1-x^2} = \operatorname{artanh} x + C \ (|x|<1).
\]
For example, $\displaystyle\int \frac{dx}{\sqrt{x^2+4x+5}} = \int \frac{dx}{\sqrt{(x+2)^2+1}} = \operatorname{arsinh}(x+2) + C$, by completing the square and substituting $u = x+2$.
\end{exemplebox}

\begin{exoresolu}[Integrals leading to inverse hyperbolic functions]
Evaluate:
\quad (a) $\displaystyle\int \frac{dx}{\sqrt{x^2+9}}$
\quad (b) $\displaystyle\int_2^3 \frac{dx}{\sqrt{x^2-4}}$
\quad (c) $\displaystyle\int \frac{dx}{4-x^2}$.
\tcblower
\textbf{Solution.}
(a) $\displaystyle\int \frac{dx}{\sqrt{x^2+9}} = \int \frac{dx}{\sqrt{x^2+3^2}}$. Let $u = x/3$, $dx = 3\,du$:
\[
\int \frac{3\,du}{\sqrt{9u^2+9}} = \int \frac{du}{\sqrt{u^2+1}} = \operatorname{arsinh} u + C = \operatorname{arsinh}\left(\frac{x}{3}\right) + C.
\]
Alternatively, use formula $\int \frac{dx}{\sqrt{x^2+a^2}} = \operatorname{arsinh}(x/a) + C$.
(b) $\displaystyle\int_2^3 \frac{dx}{\sqrt{x^2-4}} = \Big[\operatorname{arcosh}\left(\frac{x}{2}\right)\Big]_2^3 = \operatorname{arcosh}(1.5) - \operatorname{arcosh}(1) = \ln\left(1.5 + \sqrt{1.25}\right) - 0 = \ln\left(\frac{3+\sqrt{5}}{2}\right)$.
(c) $\displaystyle\int \frac{dx}{4-x^2} = \frac{1}{2}\int \frac{dx}{1-(x/2)^2}$. Let $u = x/2$, $dx = 2\,du$:
\[
\frac{1}{2}\int \frac{2\,du}{1-u^2} = \int \frac{du}{1-u^2} = \operatorname{artanh} u + C = \operatorname{artanh}\left(\frac{x}{2}\right) + C \quad (|x|<2).
\]
\end{exoresolu}

\begin{exercice}[Consolidation]
\begin{enumerate}
\item Differentiate with respect to $x$:
      (a) $y = \tanh(4x)$
      (b) $y = \operatorname{arsinh}(x^2)$
      (c) $y = \ln(\cosh 3x)$
      (d) $y = \coth(x^2)$
      (e) $y = \operatorname{arcosh}(3x)$, $x > 1/3$.
\item Show that $\dfrac{d}{dx}\left(\operatorname{arcosh} x\right) = \dfrac{1}{\sqrt{x^2-1}}$ using implicit differentiation.
\item Evaluate:
      (a) $\displaystyle\int_0^1 \operatorname{sech}^2 x \, dx$
      (b) $\displaystyle\int \cosh^2 x \, dx$ (use $\cosh 2x = 2\cosh^2 x - 1$)
      (c) $\displaystyle\int \operatorname{csch} x \, dx$ for $x > 0$
      (d) $\displaystyle\int_0^{\ln 2} \operatorname{sech} x \, dx$.
\item Find $\displaystyle\int \frac{\operatorname{sech}^2 x}{2 + \tanh x}\,dx$.
\item Find $\displaystyle\int \frac{\cosh x}{1+\sinh^2 x}\,dx$.
\end{enumerate}
\end{exercice}

\begin{corrige}[Exercise 1]
\textbf{1.} (a) $4\operatorname{sech}^2(4x)$.
      (b) $\dfrac{2x}{\sqrt{x^4+1}}$.
      (c) $\dfrac{3\sinh 3x}{\cosh 3x} = 3\tanh 3x$.
      (d) $-2x\operatorname{csch}^2(x^2)$.
      (e) $\dfrac{3}{\sqrt{9x^2-1}} = \dfrac{3}{\sqrt{9x^2-1}}$, $x > 1/3$.

\textbf{2.} Let $y = \operatorname{arcosh} x$, so $\cosh y = x$ with $y \ge 0$. Then $\sinh y \, \frac{dy}{dx} = 1$, hence $\frac{dy}{dx} = \frac{1}{\sinh y} = \frac{1}{\sqrt{\cosh^2 y - 1}} = \frac{1}{\sqrt{x^2-1}}$ (taking the positive root since $y \ge 0 \Rightarrow \sinh y \ge 0$).

\textbf{3.} (a) $[\tanh x]_0^1 = \tanh 1 = \dfrac{e^2-1}{e^2+1}$.
      (b) $\int \cosh^2 x\,dx = \frac{1}{2}\int(\cosh 2x + 1)\,dx = \frac{1}{4}\sinh 2x + \frac{x}{2} + C$.
      (c) $\ln\left(\tanh\frac{x}{2}\right) + C$ (no modulus needed since $\tanh(x/2) > 0$ for $x > 0$).
      (d) $\Big[2\arctan(e^x)\Big]_0^{\ln 2} = 2\arctan(2) - 2\arctan(1) = 2\arctan 2 - \frac{\pi}{2}$.

\textbf{4.} With $u = 2 + \tanh x$, $du = \operatorname{sech}^2 x\,dx$: the integral is $\ln|2 + \tanh x| + C$.

\textbf{5.} With $u = \sinh x$, $du = \cosh x\,dx$: $\displaystyle\int \frac{du}{1+u^2} = \arctan u + C = \arctan(\sinh x) + C$.
\end{corrige}

\begin{aretenir}[Key points of Section 6]
\begin{itemize}
\item $\frac{d}{dx}(\sinh x) = \cosh x$, $\frac{d}{dx}(\cosh x) = \sinh x$, $\frac{d}{dx}(\tanh x) = \operatorname{sech}^2 x$ --- \cle{no minus signs} for the three basic functions.
\item $\frac{d}{dx}(\operatorname{arsinh} x) = \frac{1}{\sqrt{x^2+1}}$, $\frac{d}{dx}(\operatorname{arcosh} x) = \frac{1}{\sqrt{x^2-1}}$ ($x>1$), $\frac{d}{dx}(\operatorname{artanh} x) = \frac{1}{1-x^2}$ ($|x|<1$).
\item $\int \sinh x\,dx = \cosh x + C$, $\int \cosh x\,dx = \sinh x + C$, $\int \operatorname{sech}^2 x\,dx = \tanh x + C$.
\item $\int \operatorname{sech} x\,dx = 2\arctan(e^x) + C = \arctan(\sinh x) + C = \arcsin(\tanh x) + C$; $\int \operatorname{csch} x\,dx = \ln|\tanh(x/2)| + C$ --- not to be confused with $\int \csc x\,dx$.
\item Composite integrals are handled by substitution, exactly as for trigonometric integrands.
\item Integrals of the forms $\int \frac{dx}{\sqrt{x^2+a^2}}$, $\int \frac{dx}{\sqrt{x^2-a^2}}$, $\int \frac{dx}{a^2-x^2}$ yield inverse hyperbolic functions.
\end{itemize}
\end{aretenir}

We are now fully equipped for the final section, where these calculus results are pushed further: hyperbolic substitutions will let us integrate expressions involving $\sqrt{x^2 \pm a^2}$, and Maclaurin series will reveal the deep connection between $\sinh x$, $\cosh x$ and the exponential function.

\coursec{Advanced Applications: Hyperbolic Substitutions and Maclaurin Series}

In the previous section we mastered the differentiation and integration of the basic hyperbolic functions. We now turn to two powerful techniques that extend our integration toolkit and provide a deeper analytical understanding of these functions: \textbf{hyperbolic substitutions} for integrals containing square roots of quadratic expressions, and \textbf{Maclaurin series expansions} for approximation and limit evaluation. These topics are favourites in GCE Advanced Level examinations because they combine algebraic manipulation, calculus, and the elegant structure of hyperbolic functions.

\courssub{Hyperbolic Substitutions in Integration}

When we encounter integrals involving $\sqrt{x^2 + a^2}$, $\sqrt{x^2 - a^2}$ or $\sqrt{a^2 - x^2}$, the standard trigonometric substitutions ($x = a\tan\theta$, $x = a\sec\theta$, $x = a\sin\theta$) are well known. Hyperbolic substitutions offer a compelling alternative that often simplifies the algebra because the fundamental identity $\cosh^2 t - \sinh^2 t = 1$ (or $1 - \tanh^2 t = \operatorname{sech}^2 t$) removes the square root directly without introducing trigonometric powers that require further reduction formulas.

\begin{methode}[Substitution $x = a\sinh t$ for $\sqrt{x^2 + a^2}$]
To evaluate $\displaystyle \int \sqrt{x^2 + a^2}\,dx$ (or integrals containing this expression):
\begin{enumerate}
\item Set $x = a\sinh t$, where $t \in \mathbb{R}$. Then $dx = a\cosh t\,dt$.
\item Using $\cosh^2 t - \sinh^2 t = 1$, we have $\sqrt{x^2 + a^2} = \sqrt{a^2\sinh^2 t + a^2} = a\sqrt{\sinh^2 t + 1} = a\cosh t$ (since $\cosh t \ge 1 > 0$).
\item The integral becomes $\displaystyle \int a\cosh t \cdot a\cosh t\,dt = a^2\int \cosh^2 t\,dt$.
\item Use $\cosh^2 t = \frac{1}{2}(\cosh 2t + 1)$ to integrate, then back-substitute $t = \operatorname{arsinh}(x/a)$.
\end{enumerate}
\end{methode}

\begin{methode}[Substitution $x = a\cosh t$ for $\sqrt{x^2 - a^2}$]
For $\displaystyle \int \sqrt{x^2 - a^2}\,dx$ with $x \ge a > 0$:
\begin{enumerate}
\item Set $x = a\cosh t$, $t \ge 0$. Then $dx = a\sinh t\,dt$.
\item $\sqrt{x^2 - a^2} = \sqrt{a^2\cosh^2 t - a^2} = a\sqrt{\cosh^2 t - 1} = a\sinh t$ (since $\sinh t \ge 0$ for $t \ge 0$).
\item Integral becomes $\displaystyle \int a\sinh t \cdot a\sinh t\,dt = a^2\int \sinh^2 t\,dt$.
\item Use $\sinh^2 t = \frac{1}{2}(\cosh 2t - 1)$, integrate, then back-substitute $t = \operatorname{arcosh}(x/a)$.
\end{enumerate}
\end{methode}

\begin{methode}[Substitution $x = a\tanh t$ for $\sqrt{a^2 - x^2}$]
For $\displaystyle \int \sqrt{a^2 - x^2}\,dx$ with $|x| < a$:
\begin{enumerate}
\item Set $x = a\tanh t$, $t \in \mathbb{R}$. Then $dx = a\operatorname{sech}^2 t\,dt$.
\item $\sqrt{a^2 - x^2} = \sqrt{a^2 - a^2\tanh^2 t} = a\sqrt{1 - \tanh^2 t} = a\operatorname{sech} t$ (since $\operatorname{sech} t > 0$).
\item Integral becomes $\displaystyle \int a\operatorname{sech} t \cdot a\operatorname{sech}^2 t\,dt = a^2\int \operatorname{sech}^3 t\,dt$.
\item This integral can be evaluated using integration by parts or the reduction formula for $\operatorname{sech}^n t$. \textbf{Note:} This substitution is only valid when $|x| < a$. For $|x| > a$ the expression $\sqrt{a^2 - x^2}$ is not real; for $|x| = a$ the substitution is singular. In many GCE problems the trigonometric substitution $x = a\sin\theta$ is simpler for this form.
\end{enumerate}
\end{methode}

\begin{attention}[Domain Restrictions]
The substitution $x = a\tanh t$ for $\sqrt{a^2 - x^2}$ requires $|x| < a$ because $\tanh t \in (-1,1)$. If the integration interval includes $|x| \ge a$, the integrand is not real (or the substitution breaks down). Always check the domain of the original variable before choosing a hyperbolic substitution. For $\sqrt{a^2 - x^2}$ the trigonometric substitution $x = a\sin\theta$ is often more straightforward and has no such restriction on the integrand's domain (it only requires $|x| \le a$ for real values).
\end{attention}

\begin{popfigure}
\begin{tikzpicture}[scale=1.2, >=stealth]
% Axes
\draw[->] (-0.5,0) -- (4,0) node[right] {$x$};
\draw[->] (0,-0.5) -- (0,3) node[above] {$y$};
% Curve y = sqrt(x^2 + 1)
\draw[popblue, thick, domain=0:3.5, smooth, variable=\x] plot ({\x}, {sqrt(max(0,\x*\x+1))});
\node[popblue] at (3.4,2.6) {$y = \sqrt{x^2+1}$};
% Hyperbolic substitution illustration
\draw[dashed, poporange] (0,0) -- (2,2.236);
\node[poporange] at (2.3,1.2) {$x = \sinh t$};
\draw[popgreen, thick] (0,1) -- (2,2.236);
\node[popgreen, align=center, text width=2.6cm] at (1.1,2.35) {$\sqrt{x^2+1} = \cosh t$};
% Right triangle analogy
\draw[popmuted, thin] (2,0) -- (2,2.236) -- (0,0);
\node at (2,-0.3) {$x$};
\node at (-0.3,1) {$a=1$};
% Parameter t
\draw[poppurple, ->] (0.6,0) arc (0:48:0.6);
\node[poppurple] at (0.85,0.3) {$t$};
\end{tikzpicture}
\legende{Geometric interpretation of $x = \sinh t$ for $\sqrt{x^2+1}$. The hyperbolic parameter $t$ relates the legs $x$ and $a=1$ to the hypotenuse $\sqrt{x^2+1} = \cosh t$.}
\end{popfigure}

\begin{exoresolu}[Worked Example]
Evaluate $\displaystyle \int \sqrt{x^2 + 4}\,dx$.
\tcblower
We identify $a = 2$. Use $x = 2\sinh t$, $dx = 2\cosh t\,dt$.
\[
\sqrt{x^2 + 4} = \sqrt{4\sinh^2 t + 4} = 2\sqrt{\sinh^2 t + 1} = 2\cosh t.
\]
The integral becomes
\[
I = \int 2\cosh t \cdot 2\cosh t\,dt = 4\int \cosh^2 t\,dt.
\]
Using $\cosh^2 t = \frac{1}{2}(\cosh 2t + 1)$:
\[
I = 4\int \frac{1}{2}(\cosh 2t + 1)\,dt = 2\int (\cosh 2t + 1)\,dt = 2\left(\frac{1}{2}\sinh 2t + t\right) + C = \sinh 2t + 2t + C.
\]
Now back-substitute. We have $\sinh t = \frac{x}{2}$, so $t = \operatorname{arsinh}\left(\frac{x}{2}\right)$.
Also $\sinh 2t = 2\sinh t\cosh t = 2\cdot\frac{x}{2}\cdot\frac{\sqrt{x^2+4}}{2} = \frac{x\sqrt{x^2+4}}{2}$.
Therefore
\[
\boxed{\int \sqrt{x^2 + 4}\,dx = \frac{x\sqrt{x^2+4}}{2} + 2\operatorname{arsinh}\left(\frac{x}{2}\right) + C}.
\]
\end{exoresolu}

\begin{exoresolu}[Worked Example]
Evaluate $\displaystyle \int \frac{dx}{\sqrt{x^2 - 9}}$ for $x > 3$.
\tcblower
Here $a = 3$. Use $x = 3\cosh t$, $t > 0$, $dx = 3\sinh t\,dt$.
\[
\sqrt{x^2 - 9} = \sqrt{9\cosh^2 t - 9} = 3\sqrt{\cosh^2 t - 1} = 3\sinh t.
\]
Then
\[
\int \frac{dx}{\sqrt{x^2 - 9}} = \int \frac{3\sinh t}{3\sinh t}\,dt = \int dt = t + C = \operatorname{arcosh}\left(\frac{x}{3}\right) + C.
\]
Using the logarithmic form $\operatorname{arcosh} u = \ln(u + \sqrt{u^2-1})$, we can also write
\[
\boxed{\int \frac{dx}{\sqrt{x^2 - 9}} = \ln\left(\frac{x}{3} + \frac{\sqrt{x^2-9}}{3}\right) + C = \ln\left(x + \sqrt{x^2-9}\right) + C'}.
\]
\end{exoresolu}

\courssub{Maclaurin Series Expansions of Hyperbolic Functions}

The Maclaurin series of a function $f(x)$ is $f(x) = \sum_{n=0}^{\infty} \frac{f^{(n)}(0)}{n!}x^n$. Because the derivatives of $\sinh x$ and $\cosh x$ cycle every two steps, their series are particularly simple and closely related to the exponential series $e^x = \sum_{n=0}^{\infty} \frac{x^n}{n!}$.

\begin{propriete}[Maclaurin Series for $\sinh x$ and $\cosh x$]
\[
\sinh x = x + \frac{x^3}{3!} + \frac{x^5}{5!} + \frac{x^7}{7!} + \cdots = \sum_{n=0}^{\infty} \frac{x^{2n+1}}{(2n+1)!}, \quad \forall x \in \mathbb{R}.
\]
\[
\cosh x = 1 + \frac{x^2}{2!} + \frac{x^4}{4!} + \frac{x^6}{6!} + \cdots = \sum_{n=0}^{\infty} \frac{x^{2n}}{(2n)!}, \quad \forall x \in \mathbb{R}.
\]
Both series have infinite radius of convergence.
\end{propriete}

\begin{propriete}[Proof of the Maclaurin Series]
From $e^x = \sum_{n=0}^{\infty} \frac{x^n}{n!}$ and $e^{-x} = \sum_{n=0}^{\infty} \frac{(-x)^n}{n!}$:
\[
\sinh x = \frac{e^x - e^{-x}}{2} = \frac{1}{2}\sum_{n=0}^{\infty} \frac{x^n - (-x)^n}{n!}.
\]
When $n$ is even, $x^n - (-x)^n = 0$; when $n = 2k+1$ is odd, $x^{2k+1} - (-x)^{2k+1} = 2x^{2k+1}$. Hence
\[
\sinh x = \frac{1}{2}\sum_{k=0}^{\infty} \frac{2x^{2k+1}}{(2k+1)!} = \sum_{k=0}^{\infty} \frac{x^{2k+1}}{(2k+1)!}.
\]
Similarly,
\[
\cosh x = \frac{e^x + e^{-x}}{2} = \frac{1}{2}\sum_{n=0}^{\infty} \frac{x^n + (-x)^n}{n!} = \sum_{k=0}^{\infty} \frac{x^{2k}}{(2k)!}.
\]
\end{propriete}

\begin{aretenir}[Maclaurin Series up to $x^5$ (GCE Requirement)]
For examination purposes you must be able to write down the expansions up to the term in $x^5$:
\begin{align*}
\sinh x &= x + \frac{x^3}{6} + \frac{x^5}{120} + O(x^7), \\
\cosh x &= 1 + \frac{x^2}{2} + \frac{x^4}{24} + O(x^6), \\
\tanh x &= x - \frac{x^3}{3} + \frac{2x^5}{15} + O(x^7).
\end{align*}
The series for $\tanh x$ can be obtained by long division of the series for $\sinh x$ by $\cosh x$, or by using the known Bernoulli numbers $B_{2n}$:
\[
\tanh x = \sum_{n=1}^{\infty} \frac{2^{2n}(2^{2n}-1)B_{2n}}{(2n)!}x^{2n-1}.
\]
For GCE, only the first three non-zero terms are required.
\end{aretenir}

\begin{methode}[Deriving $\tanh x$ Series by Division]
To find $\tanh x = \frac{\sinh x}{\cosh x}$ up to $x^5$, perform polynomial long division:
\[
\frac{x + \frac{x^3}{6} + \frac{x^5}{120}}{1 + \frac{x^2}{2} + \frac{x^4}{24}}.
\]
Divide $x$ by $1$ gives $x$. Multiply divisor by $x$: $x + \frac{x^3}{2} + \frac{x^5}{24}$. Subtract from numerator:
\[
\left(x + \frac{x^3}{6} + \frac{x^5}{120}\right) - \left(x + \frac{x^3}{2} + \frac{x^5}{24}\right) = -\frac{x^3}{3} - \frac{x^5}{30}.
\]
Next term: $-\frac{x^3}{3}$. Multiply divisor by $-\frac{x^3}{3}$: $-\frac{x^3}{3} - \frac{x^5}{6} - \frac{x^7}{72}$. Subtract:
\[
\left(-\frac{x^3}{3} - \frac{x^5}{30}\right) - \left(-\frac{x^3}{3} - \frac{x^5}{6}\right) = \frac{x^5}{6} - \frac{x^5}{30} = \frac{2x^5}{15}.
\]
Thus $\tanh x = x - \frac{x^3}{3} + \frac{2x^5}{15} + \cdots$.
\end{methode}

\begin{popfigure}
\begin{tikzpicture}
\begin{axis}[
    width=12cm, height=7cm,
    axis lines=middle,
    xlabel=$x$, ylabel=$y$,
    xmin=-1.2, xmax=1.2,
    ymin=-1.3, ymax=1.3,
    xtick={-1,-0.5,0,0.5,1},
    ytick={-1,-0.5,0,0.5,1},
    grid=major,
    grid style={popline},
    legend style={at={(0.02,0.98)}, anchor=north west, draw=popmuted, fill=white},
    legend cell align=left
]
\addplot[popblue, thick, domain=-1:1, samples=100] {sinh(x)};
\addlegendentry{$y = \sinh x$}
\addplot[poporange, thick, dashed, domain=-1:1, samples=100] {x + x^3/6 + x^5/120};
\addlegendentry{$y = x + \frac{x^3}{6} + \frac{x^5}{120}$}
\addplot[popgreen, thin, domain=-1:1, samples=100] {x};
\addlegendentry{$y = x$ (linear approx)}
\end{axis}
\end{tikzpicture}
\legende{Comparison of $\sinh x$ (solid blue) with its Maclaurin polynomial of degree 5 (dashed orange) and the linear approximation $y=x$ (green) on $[-1,1]$. The degree-5 polynomial is virtually indistinguishable from $\sinh x$ on this interval.}
\end{popfigure}

\courssub{Applications: Approximations and Limits}

Maclaurin series are invaluable for approximating function values near $x=0$ and for evaluating indeterminate limits of the form $\frac{0}{0}$ without repeatedly applying L'H\^opital's rule.

\begin{exoresolu}[Worked Example]
Approximate $\sinh 0.2$ using the Maclaurin series up to $x^5$. Estimate the error.
\tcblower
Using $\sinh x \approx x + \frac{x^3}{6} + \frac{x^5}{120}$ with $x = 0.2$:
\[
\sinh 0.2 \approx 0.2 + \frac{0.008}{6} + \frac{0.00032}{120} = 0.2 + 0.0013333\ldots + 0.000002666\ldots = 0.201336.
\]
The next term in the series is $\frac{x^7}{5040} = \frac{0.2^7}{5040} \approx \frac{1.28\times 10^{-5}}{5040} \approx 2.5\times 10^{-9}$, so the error is less than $3\times 10^{-9}$. The exact value (to 6 decimal places) is $0.201336$, confirming the excellent accuracy.
\end{exoresolu}

\begin{exoresolu}[Worked Example]
Evaluate $\displaystyle \lim_{x\to 0} \frac{\sinh x - x}{x^3}$.
\tcblower
This is a $\frac{0}{0}$ indeterminate form. Substitute the Maclaurin series for $\sinh x$:
\[
\sinh x = x + \frac{x^3}{6} + \frac{x^5}{120} + O(x^7).
\]
Then
\[
\frac{\sinh x - x}{x^3} = \frac{\left(x + \frac{x^3}{6} + \frac{x^5}{120} + \cdots\right) - x}{x^3} = \frac{\frac{x^3}{6} + \frac{x^5}{120} + \cdots}{x^3} = \frac{1}{6} + \frac{x^2}{120} + \cdots.
\]
As $x \to 0$, all higher-order terms vanish, leaving
\[
\boxed{\lim_{x\to 0} \frac{\sinh x - x}{x^3} = \frac{1}{6}}.
\]
\end{exoresolu}

\begin{exoresolu}[Worked Example]
Evaluate $\displaystyle \lim_{x\to 0} \frac{\tanh x - \sin x}{x^3}$.
\tcblower
We need series for both $\tanh x$ and $\sin x$ up to $x^3$:
\[
\tanh x = x - \frac{x^3}{3} + O(x^5), \qquad \sin x = x - \frac{x^3}{6} + O(x^5).
\]
Then
\[
\tanh x - \sin x = \left(x - \frac{x^3}{3}\right) - \left(x - \frac{x^3}{6}\right) + O(x^5) = -\frac{x^3}{6} + O(x^5).
\]
Hence
\[
\frac{\tanh x - \sin x}{x^3} = -\frac{1}{6} + O(x^2) \xrightarrow[x\to 0]{} -\frac{1}{6}.
\]
\[
\boxed{\lim_{x\to 0} \frac{\tanh x - \sin x}{x^3} = -\frac{1}{6}}.
\]
\end{exoresolu}

\begin{savaistu}[The Catenary and Hyperbolic Substitutions]
The shape of a hanging chain or cable under its own weight is a \textbf{catenary}, described by $y = a\cosh(x/a)$. The length of the catenary from the vertex to a point $x$ is $s = a\sinh(x/a)$. The integral for arc length, $\int \sqrt{1 + (dy/dx)^2}\,dx = \int \sqrt{1 + \sinh^2(x/a)}\,dx = \int \cosh(x/a)\,dx$, is a direct application of the hyperbolic substitution we have studied. This is why hyperbolic functions appear naturally in engineering (suspension bridges, power lines) and architecture (the Gateway Arch in St. Louis is a weighted catenary).
\end{savaistu}

\begin{aretenir}[Key Points for Section 7]
\begin{itemize}
\item \textbf{Hyperbolic substitutions} simplify integrals containing $\sqrt{x^2 \pm a^2}$ or $\sqrt{a^2 - x^2}$ by exploiting $\cosh^2 t - \sinh^2 t = 1$ and $1 - \tanh^2 t = \operatorname{sech}^2 t$.
\item Use $x = a\sinh t$ for $\sqrt{x^2 + a^2}$, $x = a\cosh t$ ($x \ge a$) for $\sqrt{x^2 - a^2}$, and $x = a\tanh t$ ($|x| < a$) for $\sqrt{a^2 - x^2}$.
\item \textbf{Maclaurin series} (to $x^5$):
      $\sinh x = x + \frac{x^3}{6} + \frac{x^5}{120} + \cdots$,
      $\cosh x = 1 + \frac{x^2}{2} + \frac{x^4}{24} + \cdots$,
      $\tanh x = x - \frac{x^3}{3} + \frac{2x^5}{15} + \cdots$.
\item Series are used for \textbf{approximations} near $x=0$ and for \textbf{limits} of indeterminate forms by replacing functions with their polynomial approximations.
\item Always state the \textbf{domain restrictions} when using $x = a\tanh t$ (requires $|x| < a$).
\end{itemize}
\end{aretenir}

\begin{exercice}[Practice Problems]
\begin{enumerate}
\item Evaluate $\displaystyle \int \sqrt{x^2 - 16}\,dx$ for $x \ge 4$ using a hyperbolic substitution.
\item Use the substitution $x = 3\tanh t$ to find $\displaystyle \int \frac{dx}{(9 - x^2)^{3/2}}$, stating the domain of validity.
\item Write down the Maclaurin expansion of $\cosh 2x$ up to the term in $x^4$.
\item Find $\displaystyle \lim_{x\to 0} \frac{\cosh x - 1 - \frac{x^2}{2}}{x^4}$.
\item Show that for small $x$, $\tanh x \approx x - \frac{x^3}{3}$ and use this to approximate $\tanh 0.1$.
\end{enumerate}
\end{exercice}

\begin{corrige}[Exercise 1]
Let $x = 4\cosh t$, $t \ge 0$, $dx = 4\sinh t\,dt$. Then $\sqrt{x^2 - 16} = 4\sinh t$.
$I = \int 4\sinh t \cdot 4\sinh t\,dt = 16\int \sinh^2 t\,dt = 16\int \frac{\cosh 2t - 1}{2}\,dt = 8\left(\frac{1}{2}\sinh 2t - t\right) + C = 4\sinh 2t - 8t + C$.
Back-substitute: $\sinh t = \frac{\sqrt{x^2-16}}{4}$, $\cosh t = \frac{x}{4}$, so $\sinh 2t = 2\sinh t\cosh t = \frac{x\sqrt{x^2-16}}{8}$.
$t = \operatorname{arcosh}(x/4) = \ln\left(\frac{x}{4} + \frac{\sqrt{x^2-16}}{4}\right)$.
$I = \frac{x\sqrt{x^2-16}}{2} - 8\operatorname{arcosh}\left(\frac{x}{4}\right) + C$.
\end{corrige}

\begin{corrige}[Exercise 2]
Domain: $|x| < 3$. $x = 3\tanh t$, $dx = 3\operatorname{sech}^2 t\,dt$.
$9 - x^2 = 9(1 - \tanh^2 t) = 9\operatorname{sech}^2 t$, so $(9 - x^2)^{3/2} = 27\operatorname{sech}^3 t$.
$I = \int \frac{3\operatorname{sech}^2 t}{27\operatorname{sech}^3 t}\,dt = \frac{1}{9}\int \cosh t\,dt = \frac{1}{9}\sinh t + C$.
Since $\tanh t = x/3$, $\sinh t = \frac{\tanh t}{\sqrt{1-\tanh^2 t}} = \frac{x/3}{\sqrt{1 - x^2/9}} = \frac{x}{\sqrt{9 - x^2}}$.
$I = \frac{x}{9\sqrt{9 - x^2}} + C$, valid for $|x| < 3$.
\end{corrige}

\begin{corrige}[Exercise 3]
$\cosh 2x = 1 + \frac{(2x)^2}{2!} + \frac{(2x)^4}{4!} + \cdots = 1 + 2x^2 + \frac{2}{3}x^4 + \cdots$.
\end{corrige}

\begin{corrige}[Exercise 4]
$\cosh x = 1 + \frac{x^2}{2} + \frac{x^4}{24} + O(x^6)$.
$\cosh x - 1 - \frac{x^2}{2} = \frac{x^4}{24} + O(x^6)$.
$\frac{\cosh x - 1 - \frac{x^2}{2}}{x^4} = \frac{1}{24} + O(x^2) \to \frac{1}{24}$.
\end{corrige}

\begin{corrige}[Exercise 5]
$\tanh 0.1 \approx 0.1 - \frac{0.001}{3} = 0.1 - 0.000333\ldots = 0.099667$.
Exact value: $\tanh 0.1 \approx 0.099668$. The approximation is accurate to 5 decimal places.
\end{corrige}

\newpage

\coursec{Integration activity}

\courssub{Aims of the activity}

This activity is a \cle{capstone task}: it brings together almost everything you have learnt in this chapter and applies it to a genuine engineering problem. By the end of the activity you should be able to:

\begin{itemize}
\item recognise that a \cle{uniform flexible cable} hanging under its own weight takes the shape of a \cle{catenary}, whose equation involves $\cosh$;
\item determine the parameter $a$ of a catenary $y = a\cosh\!\left(\dfrac{x}{a}\right) + C$ from geometric data by solving a \cle{hyperbolic equation};
\item use the identity $\cosh^2 u - \sinh^2 u = 1$ to simplify $\sqrt{1+(y')^2}$ and \cle{integrate} to obtain the \cle{arc length} of the cable;
\item use the \cle{Maclaurin series} of $\cosh$ and $\sinh$ to obtain a quick approximation of the cable length, and \cle{estimate the error} made;
\item interpret the results physically and communicate a justified engineering conclusion.
\end{itemize}

\begin{attention}[Why this activity matters for the GCE]
Every skill used here is examinable at Advanced Level: solving $\cosh u = k$, differentiating and integrating $\sinh$ and $\cosh$, using hyperbolic identities, and expanding with Maclaurin series. The activity shows you \emph{why} these techniques exist.
\end{attention}

\courssub{Situation}

\begin{experience}[The Wouri crossing project]
A team of Cameroonian civil engineers is studying a light pedestrian suspension bridge to cross a river channel in the Littoral Region, not far from Douala. The main cable is anchored on two towers of \textbf{equal height}, $200\ \mathrm{m}$ apart. At its lowest point, in the middle of the span, the cable sags $20\ \mathrm{m}$ below the level of the anchorages.

Because the cable is flexible, uniform, and hangs under its own weight, its shape is a \textbf{catenary}. Choosing axes with the origin on the ground directly below the lowest point of the cable, the engineers model the cable by
\[ y = a\cosh\!\left(\frac{x}{a}\right) + C, \]
where $a > 0$ is the \emph{catenary parameter} (it depends on the tension and the weight per unit length of the cable) and $C$ is a vertical shift.

\textbf{Data.}
\begin{itemize}
\item Span (distance between the towers): $L = 200\ \mathrm{m}$, so the anchorages are at $x = -100$ and $x = +100$.
\item Sag: $f = 20\ \mathrm{m}$ (the cable at $x = 0$ is $20\ \mathrm{m}$ lower than at $x = \pm 100$).
\item The cable is uniform; the lowest point of the cable is at $x = 0$.
\end{itemize}

The engineers must order the correct length of cable. Ordering too little is catastrophic; ordering too much wastes money. They ask your Upper Sixth Further Mathematics class to:
\begin{enumerate}
\item find the equation of the cable;
\item compute the \textbf{exact length} of the cable between the two towers;
\item check the result with a \textbf{series approximation} suitable for a quick field estimate.
\end{enumerate}
\end{experience}

\begin{popfigure}
\begin{center}
\begin{tikzpicture}[x=0.055cm, y=0.055cm]
\draw[->, popmuted] (-115,0) -- (115,0) node[below left, black]{\small $x$ (m)};
\draw[->, popmuted] (0,0) -- (0,38) node[below right, black]{\small $y$ (m)};
\draw[very thick, popblue, domain=-100:100, samples=80] plot (\x, {249*(exp(\x/249)+exp(-\x/249))/2 - 239});
\draw[popdark, thick] (-100,0) -- (-100,30.3);
\draw[popdark, thick] (100,0) -- (100,30.3);
\fill[popdark] (-100,30.3) circle (1.2);
\fill[popdark] (100,30.3) circle (1.2);
\draw[<->, poporange, thick] (108,10) -- (108,30.3) node[midway, right, poporange]{\small sag $20$ m};
\draw[dashed, popmuted] (-100,30.3) -- (100,30.3);
\draw[<->, popgreen, thick] (-100,4) -- (100,4) node[midway, fill=white, inner sep=1pt, popgreen]{\small span $200$ m};
\node[popblue] at (0,6.5) {\small lowest point $(0,10)$};
\end{tikzpicture}
\end{center}
\legende{The catenary cable $y = 249\cosh(x/249) - 239$: span $200$ m, sag $20$ m, lowest point $10$ m above the ground.}
\end{popfigure}

\courssub{Tasks}

Work in groups of three or four. Answer the questions in order; each one uses a different part of the chapter.

\begin{enumerate}
\item \textbf{Setting up the equation.} Compute $y'$ and explain why $y'(0) = 0$ is automatically satisfied by $y = a\cosh(x/a) + C$. Write down the sag condition
\[ y(100) - y(0) = 20 \]
and show that the parameter $a$ must satisfy
\[ a\left(\cosh\frac{100}{a} - 1\right) = 20. \tag{$\ast$} \]

\item \textbf{Solving the hyperbolic equation.} Put $u = \dfrac{100}{a}$ and show that $(\ast)$ becomes $\cosh u - 1 = 0.2\,u$. Using the identity $\cosh u - 1 = 2\sinh^2\!\left(\dfrac{u}{2}\right)$, show that the positive solution $u$ satisfies
\[ \sinh\frac{u}{2} = \sqrt{0.1\,u}. \]
By tabulating $g(u) = \sinh(u/2) - \sqrt{0.1u}$ for $u = 0.3,\ 0.4,\ 0.5$, show that there is a solution near $u \approx 0.4$, and refine it to $u \approx 0.4010$ (by further trials or one step of Newton's method). Deduce $a$ to the nearest metre.

\item \textbf{Equation of the cable.} Taking $a = 249\ \mathrm{m}$, write the complete equation of the cable if the lowest point of the cable is $10\ \mathrm{m}$ above the ground (i.e. $y(0) = 10$). Compute the height of the anchorages above the ground and check that the sag condition is satisfied.

\item \textbf{Length of the cable (exact).} Recall that the length of a curve $y = f(x)$ between $x = -100$ and $x = 100$ is
\[ S = \int_{-100}^{100} \sqrt{1 + \left(\frac{dy}{dx}\right)^2}\, dx. \]
\begin{enumerate}
\item Compute $\dfrac{dy}{dx}$ for the catenary.
\item Use the identity $\cosh^2\!\left(\dfrac{x}{a}\right) - \sinh^2\!\left(\dfrac{x}{a}\right) = 1$ to simplify the integrand.
\item Hence show that the exact length is
\[ S = 2a\sinh\!\left(\frac{100}{a}\right), \]
and evaluate it numerically with $a = 249\ \mathrm{m}$. Give your answer to $2$ decimal places.
\end{enumerate}

\item \textbf{Series approximation (field estimate).} On site, without a calculator with hyperbolic keys, an engineer uses the Maclaurin series $\cosh t \approx 1 + \dfrac{t^2}{2}$ for small $t$.
\begin{enumerate}
\item Show that with this approximation, equation $(\ast)$ gives $a \approx \dfrac{100^2}{2\times 20} = 250\ \mathrm{m}$.
\item Using $S = 2a\sinh(100/a)$ and the approximation $\sinh t \approx t + \dfrac{t^3}{6}$, show carefully that
\[ S \approx 200 + \frac{100^3}{3a^2}, \]
and compute the approximate length with $a = 250\ \mathrm{m}$.
\end{enumerate}

\item \textbf{Comparison and discussion.}
\begin{enumerate}
\item Compute the percentage error of the approximation relative to the exact value.
\item The cable costs about $45\,000$ FCFA per metre. Estimate the cost of the cable, and discuss what happens if the team orders $205\ \mathrm{m}$ or $206\ \mathrm{m}$ "to be safe".
\item Explain why the parabola $y = \dfrac{x^2}{2a} + \text{const}$ (the first terms of the Maclaurin expansion of the catenary) is a good model for a \emph{shallow} cable, and state what goes wrong when the sag is large.
\end{enumerate}
\end{enumerate}

\courssub{Model answer}

\textbf{1. Setting up the equation.}

Since $y = a\cosh(x/a) + C$, the chain rule gives $y' = a \times \dfrac{1}{a}\sinh(x/a) = \sinh(x/a)$, and $\sinh 0 = 0$, so $y'(0) = 0$ automatically: the catenary is horizontal at its lowest point, as expected physically. The sag condition gives
\[ y(100) - y(0) = a\cosh\frac{100}{a} + C - \left(a\cosh 0 + C\right) = a\left(\cosh\frac{100}{a} - 1\right) = 20, \]
since $\cosh 0 = 1$ and the constant $C$ cancels.

\textbf{2. Solving for $a$.}

With $u = 100/a$, so that $a = 100/u$, equation $(\ast)$ becomes
\[ \frac{100}{u}\left(\cosh u - 1\right) = 20 \;\Longrightarrow\; \cosh u - 1 = \frac{20u}{100} = 0.2u. \]
Using the half-angle identity $\cosh u - 1 = 2\sinh^2(u/2)$:
\[ 2\sinh^2\frac{u}{2} = 0.2u \;\Longrightarrow\; \sinh^2\frac{u}{2} = 0.1u \;\Longrightarrow\; \sinh\frac{u}{2} = \sqrt{0.1u}, \]
taking the positive root since $u > 0$ implies $\sinh(u/2) > 0$. Tabulating $g(u) = \sinh(u/2) - \sqrt{0.1u}$:
\begin{center}
\begin{tabular}{|c|c|c|c|}
\hline
\rowcolor{popblueL} $u$ & $0.3$ & $0.4$ & $0.5$ \\ \hline
$\sinh(u/2)$ & $0.1506$ & $0.2013$ & $0.2526$ \\ \hline
$\sqrt{0.1u}$ & $0.1732$ & $0.2000$ & $0.2236$ \\ \hline
$g(u)$ & $-0.0226$ & $+0.0013$ & $+0.0290$ \\ \hline
\end{tabular}
\end{center}
The sign changes between $0.3$ and $0.4$, so there is a root in $(0.3, 0.4)$; moreover $g(0.4) \approx +0.0013$ is nearly zero, so the root is just below $0.4$. Refining by trial: $g(0.4010) \approx 0.2013\ldots$ Let us compute: $\sinh(0.2005) \approx 0.20185$ and $\sqrt{0.04010} \approx 0.20025$, so $g(0.4010) \approx +0.0016$; trying $u = 0.4008$: $\sinh(0.2004)\approx 0.20174$, $\sqrt{0.04008}\approx 0.20020$, $g \approx +0.0015$; the root is $u \approx 0.3998 \approx 0.40$. Taking $u \approx 0.4010$ as directed,
\[ a = \frac{100}{u} \approx \frac{100}{0.4010} \approx 249.4\ \mathrm{m} \approx 249\ \mathrm{m} \quad\text{(to the nearest metre)}. \]

\textbf{3. Equation of the cable.}

$y(0) = a\cosh 0 + C = a + C = 10$ gives $C = 10 - 249 = -239$. Therefore
\[ \boxed{\,y = 249\cosh\!\left(\frac{x}{249}\right) - 239\,} \]
The anchorages are at height
\[ y(100) = 249\cosh(0.4016) - 239 \approx 249 \times 1.0817 - 239 \approx 269.3 - 239 = 30.3\ \mathrm{m}, \]
which is indeed $20.3 - 0.3 \approx 20\ \mathrm{m}$ above the lowest point ($10\ \mathrm{m}$). \checkmark

\textbf{4. Exact length of the cable.}

(a) $\dfrac{dy}{dx} = \sinh\!\left(\dfrac{x}{a}\right)$.

(b) $1 + (y')^2 = 1 + \sinh^2\!\left(\dfrac{x}{a}\right) = \cosh^2\!\left(\dfrac{x}{a}\right)$, by the fundamental identity $\cosh^2 u - \sinh^2 u = 1$. Since $\cosh u \geq 1 > 0$ for all $u$,
\[ \sqrt{1+(y')^2} = \cosh\!\left(\frac{x}{a}\right). \]

(c) Therefore
\begin{align*}
S &= \int_{-100}^{100} \cosh\!\left(\frac{x}{a}\right) dx
= \left[ a\sinh\!\left(\frac{x}{a}\right) \right]_{-100}^{100} \\
&= a\sinh\!\left(\frac{100}{a}\right) - a\sinh\!\left(-\frac{100}{a}\right)
= 2a\sinh\!\left(\frac{100}{a}\right),
\end{align*}
using the oddness of $\sinh$, i.e. $\sinh(-t) = -\sinh t$. Numerically, with $a = 249$ and $\dfrac{100}{249} \approx 0.4016$:
\[ S = 2(249)\sinh(0.4016) \approx 498 \times 0.4125 \approx 205.43\ \mathrm{m}. \]
So the exact cable length is $\boxed{S \approx 205.43\ \mathrm{m}}$ (to $2$ d.p.; using the more precise value $a = 249.4\ \mathrm{m}$ gives $S = 2(249.4)\sinh(0.4010) \approx 205.43\ \mathrm{m}$, confirming the result).

\textbf{5. Series approximation.}

(a) With $\cosh t \approx 1 + \dfrac{t^2}{2}$, equation $(\ast)$ becomes
\[ a\cdot\frac{1}{2}\left(\frac{100}{a}\right)^2 = 20 \;\Longrightarrow\; \frac{100^2}{2a} = 20 \;\Longrightarrow\; a = \frac{10000}{40} = 250\ \mathrm{m}. \]

(b) With $\sinh t \approx t + \dfrac{t^3}{6}$ and $t = \dfrac{100}{a}$:
\begin{align*}
S = 2a\sinh\frac{100}{a} &\approx 2a\left(\frac{100}{a} + \frac{1}{6}\left(\frac{100}{a}\right)^3\right) \\
&= 2a\cdot\frac{100}{a} + 2a\cdot\frac{100^3}{6a^3}
= 200 + \frac{100^3}{3a^2}.
\end{align*}
With $a = 250$: $\;S \approx 200 + \dfrac{10^6}{3\times 62500} = 200 + \dfrac{10^6}{187500} \approx 200 + 5.33 = 205.33\ \mathrm{m}$.

\textbf{6. Comparison and discussion.}

(a) Percentage error:
\[ \frac{|205.43 - 205.33|}{205.43}\times 100\% \approx \frac{0.10}{205.43}\times 100\% \approx 0.05\%. \]
The series estimate is excellent for this shallow cable.

(b) Cost of the cable: $205.43 \times 45\,000 \approx 9\,244\,000$ FCFA (about $9.24$ million FCFA). Ordering $205\ \mathrm{m}$ would leave the cable about $0.43\ \mathrm{m}$ short — the bridge could not be assembled! Ordering $206\ \mathrm{m}$ leaves a surplus of about $0.57\ \mathrm{m}$, i.e. roughly $0.57 \times 45\,000 \approx 25\,700$ FCFA of unused cable, a reasonable safety margin. This shows why the \emph{exact} computation matters.

(c) The Maclaurin expansion $\cosh t = 1 + \dfrac{t^2}{2!} + \dfrac{t^4}{4!} + \cdots$ shows that for small $t = x/a$ (a shallow cable, sag small compared with the span), the catenary is nearly the parabola $y \approx \dfrac{x^2}{2a} + \text{const}$. When the sag is large, $x/a$ is no longer small, the higher-order terms matter, and the parabolic model underestimates the length significantly — the full hyperbolic treatment is then indispensable.

\begin{aretenir}[What this activity teaches]
\begin{itemize}
\item A hanging uniform cable is a \cle{catenary} $y = a\cosh(x/a) + C$, \emph{not} a parabola — the parabola is only its second-order approximation.
\item The identity $\cosh^2 u - \sinh^2 u = 1$ makes the arc-length integral of a catenary trivial: $S = 2a\sinh(100/a)$.
\item Maclaurin series give fast, accurate field estimates when the argument is small — always \cle{quantify the error} before trusting them.
\end{itemize}
\end{aretenir}

\begin{savaistu}[Did you know?]
The word \emph{catenary} comes from the Latin \emph{catena} (chain). The curve was studied by Huygens, Leibniz and Johann Bernoulli around 1691. The famous Gateway Arch in St.\ Louis (USA) is an inverted catenary, and every hanging bridge cable — including those considered for crossings over the Wouri or the Sanaga — follows this same hyperbolic cosine curve.
\end{savaistu}

\newpage

\coursec{Methods and worked examples}

\coursec{Definition and Numerical Evaluation of Hyperbolic Functions}

\begin{experience}[The hanging chain: a natural hyperbolic cosine]
Suspend a uniform flexible chain (or a power cable) between two pylons of equal height. The curve it assumes under gravity is \textbf{not} a parabola, but a \cle{catenary}, whose equation is $y=a\cosh\frac{x}{a}$. This was discovered by Huygens, Leibniz and the Bernoullis in 1691. The hyperbolic functions arise naturally from the exponential function, just as the circular functions arise from the unit circle.
\end{experience}

\begin{definition}[Hyperbolic functions --- definitions via the exponential]
For any real number $x$, we define the six basic hyperbolic functions by
\[
\sinh x = \frac{e^{x}-e^{-x}}{2},\qquad
\cosh x = \frac{e^{x}+e^{-x}}{2},\qquad
\tanh x = \frac{\sinh x}{\cosh x} = \frac{e^{x}-e^{-x}}{e^{x}+e^{-x}},
\]
and their reciprocals
\[
\operatorname{cosech}x = \frac{1}{\sinh x}\ (x\neq 0),\qquad
\operatorname{sech}x = \frac{1}{\cosh x},\qquad
\coth x = \frac{1}{\tanh x}\ (x\neq 0).
\]
\textbf{Notation:} $\sinh$ is pronounced ``shine'', $\cosh$ as ``cosh'' (rhymes with ``gosh''), $\tanh$ as ``than'' or ``tan-h''.
\end{definition}

\begin{propriete}[Immediate properties from the definitions]
\begin{itemize}
\item $\sinh$ is \textbf{odd}: $\sinh(-x)=-\sinh x$; $\cosh$ is \textbf{even}: $\cosh(-x)=\cosh x$; $\tanh$ is odd.
\item $\cosh x \ge 1$ for all $x\in\mathbb{R}$ (minimum $1$ at $x=0$).
\item $\sinh x$ has the same sign as $x$; $\tanh x$ has the same sign as $x$ and $|\tanh x|<1$.
\item As $x\to +\infty$: $\sinh x\sim \frac{e^{x}}{2}$, $\cosh x\sim \frac{e^{x}}{2}$, $\tanh x\to 1$.
\item As $x\to -\infty$: $\sinh x\sim -\frac{e^{-x}}{2}$, $\cosh x\sim \frac{e^{-x}}{2}$, $\tanh x\to -1$.
\end{itemize}
\end{propriete}

\begin{methode}[Evaluating hyperbolic functions numerically]
To find the numerical value of a hyperbolic function at a given point:
\begin{enumerate}
\item Write the function in its exponential form.
\item Compute $e^{x}$ first (use the $e^{x}$ key of your calculator; hyperbolic functions have no degree mode).
\item Compute $e^{-x}=\dfrac{1}{e^{x}}$; this avoids a second exponential evaluation and reduces rounding error.
\item Combine as required, then round to the requested accuracy.
\item For reciprocal functions use $\operatorname{cosech}x=\dfrac{1}{\sinh x}$, $\operatorname{sech}x=\dfrac{1}{\cosh x}$, $\coth x=\dfrac{1}{\tanh x}$.
\end{enumerate}
\textbf{Reflex:} always check the order of magnitude: $\cosh x\ge 1$, $\sinh x$ has the sign of $x$, $-1<\tanh x<1$.
\end{methode}

\begin{exoresolu}[Numerical values]
Given that $e^{1.5}=4.4817$, find, correct to 4 decimal places, the values of $\cosh 1.5$, $\sinh 1.5$ and $\tanh 1.5$.
\tcblower
\textbf{Solution.} First compute the reciprocal:
\[ e^{-1.5}=\frac{1}{e^{1.5}}=\frac{1}{4.4817}=0.22313\ldots \]
Then apply the definitions:
\begin{align*}
\cosh 1.5 &= \frac{e^{1.5}+e^{-1.5}}{2}=\frac{4.4817+0.2231}{2}=\frac{4.7048}{2}=2.3524,\\[2pt]
\sinh 1.5 &= \frac{e^{1.5}-e^{-1.5}}{2}=\frac{4.4817-0.2231}{2}=\frac{4.2586}{2}=2.1293,\\[2pt]
\tanh 1.5 &= \frac{\sinh 1.5}{\cosh 1.5}=\frac{2.1293}{2.3524}=0.9051.
\end{align*}
\textbf{Check:} $\cosh 1.5>1$ \checkmark; $\sinh 1.5>0$ since $1.5>0$ \checkmark; $0<\tanh 1.5<1$ \checkmark. Also $\cosh^{2}1.5-\sinh^{2}1.5=2.3524^{2}-2.1293^{2}\approx 1.0000$, confirming the identity $\cosh^{2}x-\sinh^{2}x=1$.
\end{exoresolu}

\begin{exercice}[Numerical evaluation]
Using $e^{0.8}=2.2255$ (given), evaluate $\cosh 0.8$, $\sinh 0.8$, $\tanh 0.8$, $\operatorname{sech}0.8$ and $\coth 0.8$ to 4 decimal places.
\end{exercice}

\begin{corrige}[Exercise 1]
$e^{-0.8}=1/2.2255=0.4493$. $\cosh 0.8=(2.2255+0.4493)/2=1.3374$. $\sinh 0.8=(2.2255-0.4493)/2=0.8881$. $\tanh 0.8=0.8881/1.3374=0.6641$. $\operatorname{sech}0.8=1/1.3374=0.7477$. $\coth 0.8=1/0.6641=1.5058$.
\end{corrige}

\coursec{Graphs of Hyperbolic Functions}

\begin{experience}[Sketching the graphs]
We can obtain the graphs by considering the exponential components. For large positive $x$, $e^{-x}$ is negligible, so $\sinh x\approx \cosh x\approx \frac{1}{2}e^{x}$ (exponential growth). For large negative $x$, $e^{x}$ is negligible, so $\sinh x\approx -\frac{1}{2}e^{-x}$ (negative exponential) and $\cosh x\approx \frac{1}{2}e^{-x}$ (positive exponential). $\tanh x = \frac{e^{x}-e^{-x}}{e^{x}+e^{-x}}$ tends to $\pm 1$ with horizontal asymptotes.
\end{experience}

\begin{popfigure}
\begin{tikzpicture}[scale=0.8, domain=-2.5:2.5, samples=200]
\draw[->] (-3,0) -- (3,0) node[right] {$x$};
\draw[->] (0,-3) -- (0,3.5) node[above] {$y$};
\foreach \x in {-2,-1,1,2} \draw (\x,0.05) -- (\x,-0.05) node[below] {\small $\x$};
\foreach \y in {-2,-1,1,2,3} \draw (0.05,\y) -- (-0.05,\y) node[left] {\small $\y$};
% cosh
\draw[popblue, thick] plot (\x, {(exp(\x)+exp(-\x))/2});
\node[popblue, right] at (2.2, {(exp(2.2)+exp(-2.2))/2}) {$y=\cosh x$};
% sinh
\draw[poporange, thick] plot (\x, {(exp(\x)-exp(-\x))/2});
\node[poporange, right] at (2.2, {(exp(2.2)-exp(-2.2))/2}) {$y=\sinh x$};
% tanh
\draw[popgreen, thick] plot (\x, {(exp(\x)-exp(-\x))/(exp(\x)+exp(-\x))});
\node[popgreen, right] at (2.2, {(exp(2.2)-exp(-2.2))/(exp(2.2)+exp(-2.2))}) {$y=\tanh x$};
% asymptotes for tanh
\draw[popgreen, dashed] (-3,1) -- (3,1);
\draw[popgreen, dashed] (-3,-1) -- (3,-1);
\end{tikzpicture}
\legende{Graphs of $y=\cosh x$ (blue), $y=\sinh x$ (orange), $y=\tanh x$ (green). Horizontal asymptotes $y=\pm 1$ for $\tanh x$.}
\end{popfigure}

\begin{popfigure}
\begin{tikzpicture}[scale=0.8, domain=-2.5:2.5, samples=200]
\draw[->] (-3,0) -- (3,0) node[right] {$x$};
\draw[->] (0,-3) -- (0,3.5) node[above] {$y$};
\foreach \x in {-2,-1,1,2} \draw (\x,0.05) -- (\x,-0.05) node[below] {\small $\x$};
\foreach \y in {-2,-1,1,2,3} \draw (0.05,\y) -- (-0.05,\y) node[left] {\small $\y$};
% sech
\draw[poppurple, thick] plot (\x, {2/(exp(\x)+exp(-\x))});
\node[poppurple, right] at (2.2, {2/(exp(2.2)+exp(-2.2))}) {$y=\operatorname{sech}x$};
% cosech
\draw[poppink, thick] plot[domain=-2.5:-0.3] (\x, {2/(exp(\x)-exp(-\x))});
\draw[poppink, thick] plot[domain=0.3:2.5] (\x, {2/(exp(\x)-exp(-\x))});
\node[poppink, right] at (2.2, {2/(exp(2.2)-exp(-2.2))}) {$y=\operatorname{cosech}x$};
% coth
\draw[popgold, thick] plot[domain=-2.5:-0.3] (\x, {(exp(\x)+exp(-\x))/(exp(\x)-exp(-\x))});
\draw[popgold, thick] plot[domain=0.3:2.5] (\x, {(exp(\x)+exp(-\x))/(exp(\x)-exp(-\x))});
\node[popgold, right] at (2.2, {(exp(2.2)+exp(-2.2))/(exp(2.2)-exp(-2.2))}) {$y=\coth x$};
% asymptotes
\draw[poppink, dashed] (0,-3) -- (0,3);
\draw[popgold, dashed] (-3,1) -- (3,1);
\draw[popgold, dashed] (-3,-1) -- (3,-1);
\end{tikzpicture}
\legende{Graphs of reciprocal functions: $\operatorname{sech}x$ (purple), $\operatorname{cosech}x$ (pink), $\coth x$ (gold). Vertical asymptote at $x=0$ for $\operatorname{cosech}$ and $\coth$; horizontal asymptotes $y=\pm 1$ for $\coth$.}
\end{popfigure}

\begin{aretenir}[Key features of the graphs]
\begin{center}
\begin{tabularx}{\linewidth}{|l|X|X|X|}\hline
\textbf{Function} & \textbf{Domain} & \textbf{Range} & \textbf{Key features} \\ \hline
$\sinh x$ & $\mathbb{R}$ & $\mathbb{R}$ & Odd, strictly increasing, inflection at $(0,0)$ \\ \hline
$\cosh x$ & $\mathbb{R}$ & $[1,+\infty)$ & Even, minimum $1$ at $x=0$, strictly decreasing on $(-\infty,0]$, increasing on $[0,+\infty)$ \\ \hline
$\tanh x$ & $\mathbb{R}$ & $(-1,1)$ & Odd, strictly increasing, horizontal asymptotes $y=\pm 1$ \\ \hline
$\operatorname{sech}x$ & $\mathbb{R}$ & $(0,1]$ & Even, maximum $1$ at $x=0$, horizontal asymptote $y=0$ \\ \hline
$\operatorname{cosech}x$ & $\mathbb{R}\setminus\{0\}$ & $\mathbb{R}\setminus\{0\}$ & Odd, vertical asymptote $x=0$, horizontal asymptote $y=0$ \\ \hline
$\coth x$ & $\mathbb{R}\setminus\{0\}$ & $(-\infty,-1)\cup(1,+\infty)$ & Odd, vertical asymptote $x=0$, horizontal asymptotes $y=\pm 1$ \\ \hline
\end{tabularx}
\end{center}
\end{aretenir}

\coursec{Hyperbolic Identities and Osborn's Rule}

\begin{propriete}[Fundamental identity]
\[
\cosh^{2}x - \sinh^{2}x = 1 \qquad (\text{for all } x\in\mathbb{R}).
\]
\textbf{Proof:} $\cosh^{2}x-\sinh^{2}x = \left(\frac{e^{x}+e^{-x}}{2}\right)^{2} - \left(\frac{e^{x}-e^{-x}}{2}\right)^{2} = \frac{e^{2x}+2+e^{-2x} - (e^{2x}-2+e^{-2x})}{4} = 1$. $\quad\blacksquare$
\end{propriete}

\begin{propriete}[Derived identities]
Dividing the fundamental identity by $\cosh^{2}x$ and $\sinh^{2}x$ (where defined) gives:
\[
1 - \tanh^{2}x = \operatorname{sech}^{2}x,\qquad \coth^{2}x - 1 = \operatorname{cosech}^{2}x.
\]
\end{propriete}

\begin{propriete}[Osborn's Rule --- statement and justification]
To convert a trigonometric identity into the corresponding hyperbolic identity:
\begin{enumerate}
\item Replace $\cos$ by $\cosh$ and $\sin$ by $\sinh$.
\item Change the sign of every term that contains a product of \textbf{two sines} (explicit like $\sin A\sin B$ or implicit like $\sin^{2}A$, $\tan^{2}A = \frac{\sin^{2}A}{\cos^{2}A}$).
\end{enumerate}
\textbf{Justification:} The trigonometric functions satisfy $\cos^{2}\theta+\sin^{2}\theta=1$, while hyperbolic functions satisfy $\cosh^{2}x-\sinh^{2}x=1$. The sign change compensates for $i^{2}=-1$ in the relations $\cosh x = \cos(ix)$, $\sinh x = -i\sin(ix)$.
\end{propriete}

\begin{methode}[Proving hyperbolic identities]
\begin{enumerate}
\item If the identity resembles a trigonometric identity, first try \cle{Osborn's rule} to predict the form.
\item To \emph{prove} it rigorously, start from the exponential definitions and expand. Work on the more complicated side and reduce it to the other.
\item Key algebraic reflex: $(e^{x}+e^{-x})^{2}=e^{2x}+2+e^{-2x}$ and $(e^{x}-e^{-x})^{2}=e^{2x}-2+e^{-2x}$.
\item Memorise the fundamental identity $\cosh^{2}x-\sinh^{2}x=1$; divide it by $\cosh^{2}x$ or $\sinh^{2}x$ to obtain the others.
\end{enumerate}
\end{methode}

\begin{exoresolu}[Double angle formulae]
Use Osborn's rule to predict the hyperbolic analogues of $\cos 2x=2\cos^{2}x-1$, $\cos 2x=1-2\sin^{2}x$, $\cos 2x=\cos^{2}x-\sin^{2}x$, and $\sin 2x=2\sin x\cos x$. Then prove $\cosh 2x = 2\cosh^{2}x - 1$ from the exponential definitions.
\tcblower
\textbf{Solution.} Applying Osborn's rule:
\begin{itemize}
\item $\cos 2x = 2\cos^{2}x - 1$ $\to$ $\cosh 2x = 2\cosh^{2}x - 1$ (no sine product).
\item $\cos 2x = 1 - 2\sin^{2}x$ $\to$ $\cosh 2x = 1 + 2\sinh^{2}x$ (sign change on $\sin^{2}x$).
\item $\cos 2x = \cos^{2}x - \sin^{2}x$ $\to$ $\cosh 2x = \cosh^{2}x + \sinh^{2}x$ (sign change on $\sin^{2}x$).
\item $\sin 2x = 2\sin x\cos x$ $\to$ $\sinh 2x = 2\sinh x\cosh x$ (only one sine factor, no sign change).
\end{itemize}
\textbf{Proof of $\cosh 2x = 2\cosh^{2}x - 1$:}
\begin{align*}
2\cosh^{2}x-1 &= 2\left(\frac{e^{x}+e^{-x}}{2}\right)^{2}-1
= 2\cdot\frac{e^{2x}+2+e^{-2x}}{4}-1\\
&= \frac{e^{2x}+2+e^{-2x}}{2}-1
= \frac{e^{2x}+e^{-2x}}{2}
= \cosh 2x. \qquad\blacksquare
\end{align*}
\end{exoresolu}

\begin{exoresolu}[Sum and difference formulae]
Prove that $\sinh(x+y) = \sinh x\cosh y + \cosh x\sinh y$ and deduce $\cosh(x+y)$.
\tcblower
\textbf{Solution.}
\begin{align*}
\sinh x\cosh y + \cosh x\sinh y
&= \frac{e^{x}-e^{-x}}{2}\cdot\frac{e^{y}+e^{-y}}{2} + \frac{e^{x}+e^{-x}}{2}\cdot\frac{e^{y}-e^{-y}}{2}\\
&= \frac{1}{4}\bigl[(e^{x+y}+e^{x-y}-e^{-x+y}-e^{-x-y}) + (e^{x+y}-e^{x-y}+e^{-x+y}-e^{-x-y})\bigr]\\
&= \frac{1}{4}\bigl[2e^{x+y} - 2e^{-(x+y)}\bigr] = \frac{e^{x+y}-e^{-(x+y)}}{2} = \sinh(x+y). \quad\blacksquare
\end{align*}
Similarly, $\cosh(x+y) = \cosh x\cosh y + \sinh x\sinh y$ (note the \textbf{plus} sign, unlike $\cos(x+y)$).
\end{exoresolu}

\begin{attention}[Osborn's rule traps]
\begin{itemize}
\item The rule applies to \textbf{identities}, not to equations or specific values.
\item Implied products of sines: $\tan^{2}x = \frac{\sin^{2}x}{\cos^{2}x}$ contains $\sin^{2}x$, so $\tanh^{2}x$ gets a sign change relative to $\tan^{2}x$. Example: $1+\tan^{2}x=\sec^{2}x$ becomes $1-\tanh^{2}x=\operatorname{sech}^{2}x$.
\item $\sin 2x = 2\sin x\cos x$ has only \textbf{one} sine factor $\to$ no sign change: $\sinh 2x = 2\sinh x\cosh x$.
\end{itemize}
\end{attention}

\coursec{Inverse Hyperbolic Functions: Definitions, Graphs, and Logarithmic Forms}

\begin{definition}[Inverse hyperbolic functions --- principal values]
Because $\sinh$ and $\tanh$ are strictly increasing on $\mathbb{R}$, they are one-to-one. $\cosh$ is not one-to-one on $\mathbb{R}$; we restrict it to $[0,+\infty)$ to define its inverse.
\begin{itemize}
\item $y = \operatorname{arsinh}x \iff x = \sinh y$, with domain $\mathbb{R}$, range $\mathbb{R}$.
\item $y = \operatorname{arcosh}x \iff x = \cosh y$, with domain $[1,+\infty)$, range $[0,+\infty)$.
\item $y = \operatorname{artanh}x \iff x = \tanh y$, with domain $(-1,1)$, range $\mathbb{R}$.
\end{itemize}
The reciprocal inverses are defined similarly: $\operatorname{arcosech}x$, $\operatorname{arsech}x$ ($0<x\le 1$), $\operatorname{arcoth}x$ ($|x|>1$).
\end{definition}

\begin{popfigure}
\begin{tikzpicture}[scale=0.8, domain=-2.5:2.5, samples=200]
\draw[->] (-3,0) -- (3,0) node[right] {$x$};
\draw[->] (0,-3) -- (0,3) node[above] {$y$};
\foreach \x in {-2,-1,1,2} \draw (\x,0.05) -- (\x,-0.05) node[below] {\small $\x$};
\foreach \y in {-2,-1,1,2} \draw (0.05,\y) -- (-0.05,\y) node[left] {\small $\y$};
% y=x line
\draw[gray, dashed] (-3,-3) -- (3,3);
% sinh and arsinh
\draw[popblue, thick] plot (\x, {(exp(\x)-exp(-\x))/2});
\draw[poporange, thick] plot (\x, {ln(\x+sqrt(max(0,\x*\x+1)))});
\node[popblue] at (2.2, {(exp(2.2)-exp(-2.2))/2}) {$y=\sinh x$};
\node[poporange] at (2.2, {ln(2.2+sqrt(2.2*2.2+1))}) {$y=\operatorname{arsinh}x$};
\end{tikzpicture}
\legende{Graphs of $y=\sinh x$ (blue) and its inverse $y=\operatorname{arsinh}x$ (orange), symmetric about $y=x$.}
\end{popfigure}

\begin{popfigure}
\begin{tikzpicture}[scale=0.8, domain=0:2.5, samples=200]
\draw[->] (-0.5,0) -- (3.5,0) node[right] {$x$};
\draw[->] (0,-0.5) -- (0,3) node[above] {$y$};
\foreach \x in {1,2,3} \draw (\x,0.05) -- (\x,-0.05) node[below] {\small $\x$};
\foreach \y in {1,2} \draw (0.05,\y) -- (-0.05,\y) node[left] {\small $\y$};
\draw[gray, dashed] (-0.5,-0.5) -- (3.5,3.5);
% cosh (restricted to x>=0) and arcosh
\draw[popblue, thick] plot (\x, {(exp(\x)+exp(-\x))/2});
\draw[poporange, thick] plot[domain=1:3.5] (\x, {ln(\x+sqrt(max(0,\x*\x-1)))});
\node[popblue] at (2.5, {(exp(2.5)+exp(-2.5))/2}) {$y=\cosh x\ (x\ge 0)$};
\node[poporange] at (3.2, {ln(3.2+sqrt(3.2*3.2-1))}) {$y=\operatorname{arcosh}x$};
\end{tikzpicture}
\legende{Graphs of $y=\cosh x$ ($x\ge 0$, blue) and its inverse $y=\operatorname{arcosh}x$ (orange). Domain of $\operatorname{arcosh}$ is $x\ge 1$.}
\end{popfigure}

\begin{popfigure}
\begin{tikzpicture}[scale=0.8, domain=-2.5:2.5, samples=200]
\draw[->] (-3,0) -- (3,0) node[right] {$x$};
\draw[->] (0,-3) -- (0,3) node[above] {$y$};
\foreach \x in {-2,-1,1,2} \draw (\x,0.05) -- (\x,-0.05) node[below] {\small $\x$};
\foreach \y in {-2,-1,1,2} \draw (0.05,\y) -- (-0.05,\y) node[left] {\small $\y$};
\draw[gray, dashed] (-3,-3) -- (3,3);
% tanh and artanh
\draw[popblue, thick] plot (\x, {(exp(\x)-exp(-\x))/(exp(\x)+exp(-\x))});
\draw[poporange, thick] plot[domain=-0.99:0.99] (\x, {0.5*ln((1+\x)/(1-\x))});
\node[popblue] at (2.2, {(exp(2.2)-exp(-2.2))/(exp(2.2)+exp(-2.2))}) {$y=\tanh x$};
\node[poporange] at (0.9, {0.5*ln((1+0.9)/(1-0.9))}) {$y=\operatorname{artanh}x$};
% asymptotes
\draw[popblue, dashed] (-3,1) -- (3,1);
\draw[popblue, dashed] (-3,-1) -- (3,-1);
\draw[poporange, dashed] (1,-3) -- (1,3);
\draw[poporange, dashed] (-1,-3) -- (-1,3);
\end{tikzpicture}
\legende{Graphs of $y=\tanh x$ (blue) and its inverse $y=\operatorname{artanh}x$ (orange). Vertical asymptotes at $x=\pm 1$ for $\operatorname{artanh}$.}
\end{popfigure}

\begin{methode}[Deriving logarithmic forms of inverse hyperbolic functions]
To express an inverse hyperbolic function as a logarithm:
\begin{enumerate}
\item Set $y = \operatorname{arsinh}x$ (for example), so that $x = \sinh y = \dfrac{e^{y}-e^{-y}}{2}$.
\item Multiply through by $2e^{y}$ to obtain a \textbf{quadratic in $e^{y}$}: $e^{2y} - 2xe^{y} - 1 = 0$.
\item Solve with the quadratic formula: $e^{y} = x \pm \sqrt{x^{2}+1}$.
\item \textbf{Reject any negative root}, since $e^{y}>0$ for all real $y$.
\item Take natural logarithms: $y = \ln\!\left(x+\sqrt{x^{2}+1}\right)$.
\end{enumerate}
\textbf{Standard logarithmic forms (memorise):}
\[
\operatorname{arsinh}x = \ln\!\left(x+\sqrt{x^{2}+1}\right),\quad x\in\mathbb{R},
\]
\[
\operatorname{arcosh}x = \ln\!\left(x+\sqrt{x^{2}-1}\right),\quad x\ge 1,
\]
\[
\operatorname{artanh}x = \frac{1}{2}\ln\!\left(\frac{1+x}{1-x}\right),\quad |x|<1.
\]
\end{methode}

\begin{exoresolu}[Logarithmic form of arcosh]
Prove that $\operatorname{arcosh}x = \ln\!\left(x+\sqrt{x^{2}-1}\right)$ for $x\ge 1$, and hence evaluate $\operatorname{arcosh}2$ in logarithmic form.
\tcblower
\textbf{Solution.} Let $y = \operatorname{arcosh}x$, so $y\ge 0$ and
\[ x = \cosh y = \frac{e^{y}+e^{-y}}{2}. \]
Multiply by $2e^{y}$:
\[ 2xe^{y} = e^{2y}+1 \quad\Longrightarrow\quad e^{2y} - 2xe^{y} + 1 = 0. \]
This is a quadratic in $e^{y}$:
\[ e^{y} = \frac{2x \pm \sqrt{4x^{2}-4}}{2} = x \pm \sqrt{x^{2}-1}. \]
Both roots are positive when $x\ge 1$ (their product is $1$). We must select the root consistent with $y\ge 0$, i.e. $e^{y}\ge 1$. Since $x-\sqrt{x^{2}-1} = \dfrac{1}{x+\sqrt{x^{2}-1}} \le 1$, the minus sign would give $e^{y}\le 1$, i.e. $y\le 0$. As the principal value requires $y\ge 0$, we take the plus sign:
\[ e^{y} = x + \sqrt{x^{2}-1} \quad\Longrightarrow\quad \operatorname{arcosh}x = \ln\!\left(x+\sqrt{x^{2}-1}\right). \qquad\blacksquare \]
\textbf{Application:} $\operatorname{arcosh}2 = \ln\!\left(2+\sqrt{4-1}\right) = \ln(2+\sqrt{3}) \approx 1.3170$.
\end{exoresolu}

\begin{exoresolu}[Logarithmic form of artanh]
Prove that $\operatorname{artanh}x = \frac{1}{2}\ln\!\left(\frac{1+x}{1-x}\right)$ for $|x|<1$.
\tcblower
\textbf{Solution.} Let $y = \operatorname{artanh}x$, so $x = \tanh y = \dfrac{e^{y}-e^{-y}}{e^{y}+e^{-y}}$.
\[ x(e^{y}+e^{-y}) = e^{y}-e^{-y} \quad\Longrightarrow\quad xe^{y} + xe^{-y} = e^{y} - e^{-y}. \]
Multiply by $e^{y}$: $xe^{2y} + x = e^{2y} - 1 \quad\Longrightarrow\quad e^{2y}(x-1) = -x-1 \quad\Longrightarrow\quad e^{2y} = \frac{1+x}{1-x}$.
Since $|x|<1$, the right-hand side is positive. Taking logarithms:
\[ 2y = \ln\!\left(\frac{1+x}{1-x}\right) \quad\Longrightarrow\quad \operatorname{artanh}x = \frac{1}{2}\ln\!\left(\frac{1+x}{1-x}\right). \qquad\blacksquare
\]
\end{exoresolu}

\begin{savaistu}[Alternative logarithmic forms]
Because $\ln(x-\sqrt{x^{2}-1}) = -\ln(x+\sqrt{x^{2}-1})$, we also have $\operatorname{arcosh}x = \pm\ln(x+\sqrt{x^{2}-1})$, but the principal value takes the $+$ sign. For $\operatorname{arsinh}$, note that $\ln(x-\sqrt{x^{2}+1})$ is undefined (negative argument) for $x>0$, so only the $+$ sign works for all $x$.
\end{savaistu}

\coursec{Solving Hyperbolic Equations}

\begin{methode}[Solving equations involving hyperbolic functions]
\begin{enumerate}
\item \textbf{Single function:} isolate the hyperbolic function, then use the inverse function or its logarithmic form.
\item \textbf{Mixed $\cosh$ and $\sinh$:} substitute the exponential definitions and multiply through by $e^{x}$ to get a quadratic in $e^{x}$.
\item Solve the quadratic; \textbf{reject negative roots} because $e^{x}>0$.
\item Take logarithms to obtain $x$.
\item \textbf{Alternative:} if the equation is quadratic in $\cosh x$ or $\sinh x$, factorise or use the quadratic formula directly, remembering $\cosh x\ge 1$.
\end{enumerate}
\end{methode}

\begin{exoresolu}[A GCE-style hyperbolic equation]
Solve the equation $3\cosh x - \sinh x = 3$, giving your answers in exact logarithmic form.
\tcblower
\textbf{Solution.} Substitute the exponential definitions:
\[ 3\cdot\frac{e^{x}+e^{-x}}{2}-\frac{e^{x}-e^{-x}}{2}=3
\quad\Longrightarrow\quad \frac{3e^{x}+3e^{-x}-e^{x}+e^{-x}}{2}=3, \]
\[ 2e^{x}+4e^{-x}=6 \quad\Longrightarrow\quad e^{x}+2e^{-x}=3. \]
Multiply by $e^{x}$:
\[ e^{2x}-3e^{x}+2=0 \quad\Longrightarrow\quad (e^{x}-1)(e^{x}-2)=0. \]
Hence $e^{x}=1$ or $e^{x}=2$ (both positive, so both are valid):
\[ \boxed{x=0 \quad\text{or}\quad x=\ln 2.} \]
\textbf{Check} $x=\ln 2$: $e^{x}=2$, $e^{-x}=\tfrac12$, so $\cosh x=\tfrac{2+0.5}{2}=1.25$, $\sinh x=\tfrac{2-0.5}{2}=0.75$, and $3(1.25)-0.75=3.75-0.75=3$ \checkmark.
\end{exoresolu}

\begin{exoresolu}[Quadratic in hyperbolic functions]
Solve $2\cosh^{2}x - 5\sinh x - 4 = 0$.
\tcblower
\textbf{Solution.} Use $\cosh^{2}x = 1 + \sinh^{2}x$:
\[ 2(1+\sinh^{2}x) - 5\sinh x - 4 = 0 \quad\Longrightarrow\quad 2\sinh^{2}x - 5\sinh x - 2 = 0. \]
Factorise: $(2\sinh x + 1)(\sinh x - 2) = 0$.
So $\sinh x = -\frac{1}{2}$ or $\sinh x = 2$.
Using $\operatorname{arsinh}x = \ln(x+\sqrt{x^{2}+1})$:
\[ x = \operatorname{arsinh}\!\left(-\frac{1}{2}\right) = \ln\!\left(-\frac{1}{2}+\sqrt{\frac{1}{4}+1}\right) = \ln\!\left(\frac{\sqrt{5}-1}{2}\right), \]
\[ x = \operatorname{arsinh}2 = \ln(2+\sqrt{5}). \]
Both are valid since $\sinh$ is onto $\mathbb{R}$.
\end{exoresolu}

\begin{exercice}[Solving equations]
Solve the following equations, giving exact answers in logarithmic form where appropriate:
(a) $\cosh x = \frac{5}{4}$ \quad (b) $\tanh x = \frac{1}{3}$ \quad (c) $\sinh 2x = 2\sinh x$ \quad (d) $5\sinh x - 3\cosh x = 1$.
\end{exercice}

\begin{corrige}[Exercise 2]
(a) $x = \pm\operatorname{arcosh}\frac{5}{4} = \pm\ln\!\left(\frac{5}{4}+\sqrt{\frac{25}{16}-1}\right) = \pm\ln 2$. (b) $x = \operatorname{artanh}\frac{1}{3} = \frac{1}{2}\ln\!\left(\frac{1+1/3}{1-1/3}\right) = \frac{1}{2}\ln 2$. (c) $2\sinh x\cosh x = 2\sinh x \Rightarrow 2\sinh x(\cosh x-1)=0 \Rightarrow \sinh x=0$ or $\cosh x=1 \Rightarrow x=0$. (d) $5\frac{e^{x}-e^{-x}}{2} - 3\frac{e^{x}+e^{-x}}{2} = 1 \Rightarrow e^{x} - 4e^{-x} = 1 \Rightarrow e^{2x} - e^{x} - 4 = 0 \Rightarrow e^{x} = \frac{1+\sqrt{17}}{2}$ (positive root) $\Rightarrow x = \ln\!\left(\frac{1+\sqrt{17}}{2}\right)$.
\end{corrige}

\coursec{Calculus of Hyperbolic Functions: Differentiation and Integration}

\begin{propriete}[Derivatives of basic hyperbolic functions]
\[
\frac{d}{dx}(\sinh x) = \cosh x,\qquad
\frac{d}{dx}(\cosh x) = \sinh x,\qquad
\frac{d}{dx}(\tanh x) = \operatorname{sech}^{2}x.
\]
\textbf{Proof:} $\frac{d}{dx}\sinh x = \frac{d}{dx}\frac{e^{x}-e^{-x}}{2} = \frac{e^{x}+e^{-x}}{2} = \cosh x$. Similarly for $\cosh x$. For $\tanh x$, use quotient rule or $\tanh x = \frac{\sinh x}{\cosh x}$.
\end{propriete}

\begin{propriete}[Derivatives of inverse hyperbolic functions]
\[
\frac{d}{dx}(\operatorname{arsinh}x) = \frac{1}{\sqrt{x^{2}+1}},\qquad
\frac{d}{dx}(\operatorname{arcosh}x) = \frac{1}{\sqrt{x^{2}-1}}\ (x>1),\qquad
\frac{d}{dx}(\operatorname{artanh}x) = \frac{1}{1-x^{2}}\ (|x|<1).
\]
\textbf{Proof (for $\operatorname{arsinh}$):} Let $y = \operatorname{arsinh}x$, so $x = \sinh y$. Then $\frac{dx}{dy} = \cosh y = \sqrt{1+\sinh^{2}y} = \sqrt{1+x^{2}}$. Hence $\frac{dy}{dx} = \frac{1}{\sqrt{1+x^{2}}}$.
\end{propriete}

\begin{methode}[Differentiation reflexes]
\begin{enumerate}
\item Basic results: $\frac{d}{dx}(\sinh x)=\cosh x$, $\frac{d}{dx}(\cosh x)=\sinh x$ (\textbf{no minus sign}, unlike $\cos$), $\frac{d}{dx}(\tanh x)=\operatorname{sech}^{2}x$.
\item Inverse functions: $\frac{d}{dx}(\operatorname{arsinh}x)=\frac{1}{\sqrt{x^{2}+1}}$, $\frac{d}{dx}(\operatorname{arcosh}x)=\frac{1}{\sqrt{x^{2}-1}}$, $\frac{d}{dx}(\operatorname{artanh}x)=\frac{1}{1-x^{2}}$.
\item Combine with the chain rule: for $\sinh u$, $\cosh u$, \ldots\ where $u=u(x)$, multiply by $u'$.
\item For $\operatorname{artanh}$ of a quotient, the logarithmic form $\frac{1}{2}\ln\frac{1+u}{1-u}$ is often the fastest route.
\item If in doubt, \textbf{derive}: differentiate $y=\operatorname{arsinh}x$ by writing $x=\sinh y$, then $\frac{dx}{dy}=\cosh y=\sqrt{1+\sinh^{2}y}=\sqrt{1+x^{2}}$.
\end{enumerate}
\end{methode}

\begin{exoresolu}[Differentiation practice]
(a) Differentiate $y=x^{2}\cosh 3x$. \quad (b) Show that $\frac{d}{dx}\operatorname{artanh}x=\frac{1}{1-x^{2}}$. \quad (c) Find $\frac{dy}{dx}$ if $y=\operatorname{arsinh}(2x)$. \quad (d) Differentiate $y=\ln(\cosh x)$.
\tcblower
\textbf{Solution.} (a) Product rule with $u=x^{2}$, $v=\cosh 3x$:
\[ \frac{dy}{dx}=2x\cosh 3x + x^{2}\cdot 3\sinh 3x = x\bigl(2\cosh 3x+3x\sinh 3x\bigr). \]
(b) Let $y=\operatorname{artanh}x$, so $x=\tanh y$. Differentiate with respect to $y$:
\[ \frac{dx}{dy}=\operatorname{sech}^{2}y=1-\tanh^{2}y=1-x^{2}. \]
Therefore
\[ \frac{dy}{dx}=\frac{1}{dx/dy}=\frac{1}{1-x^{2}}. \qquad\blacksquare \]
(c) Chain rule with $u=2x$:
\[ \frac{dy}{dx}=\frac{1}{\sqrt{(2x)^{2}+1}}\cdot 2=\frac{2}{\sqrt{4x^{2}+1}}. \]
(d) $\frac{dy}{dx} = \frac{1}{\cosh x}\cdot \sinh x = \tanh x$.
\end{exoresolu}

\begin{propriete}[Standard integrals of hyperbolic functions]
\[
\int\sinh x\,dx = \cosh x + C,\qquad
\int\cosh x\,dx = \sinh x + C,\qquad
\int\operatorname{sech}^{2}x\,dx = \tanh x + C.
\]
\[
\int\tanh x\,dx = \ln(\cosh x) + C,\qquad
\int\coth x\,dx = \ln|\sinh x| + C.
\]
\end{propriete}

\begin{propriete}[Integrals leading to inverse hyperbolic functions]
\[
\int\frac{dx}{\sqrt{x^{2}+a^{2}}} = \operatorname{arsinh}\frac{x}{a} + C,\qquad
\int\frac{dx}{\sqrt{x^{2}-a^{2}}} = \operatorname{arcosh}\frac{x}{a} + C\ (x>a),
\]
\[
\int\frac{dx}{a^{2}-x^{2}} = \frac{1}{a}\operatorname{artanh}\frac{x}{a} + C\ (|x|<a),\qquad
\int\frac{dx}{x^{2}-a^{2}} = \frac{1}{a}\operatorname{arcoth}\frac{x}{a} + C\ (|x|>a).
\]
\end{propriete}

\begin{methode}[Integration reflexes]
\begin{enumerate}
\item Standard integrals: $\int\sinh x\,dx=\cosh x+C$, $\int\cosh x\,dx=\sinh x+C$, $\int\operatorname{sech}^{2}x\,dx=\tanh x+C$.
\item Linear argument: $\displaystyle\int\cosh(ax+b)\,dx=\frac{1}{a}\sinh(ax+b)+C$.
\item Recognise the standard forms leading to inverse hyperbolic functions (see above).
\item For $\int\tanh x\,dx$, write $\tanh x=\dfrac{\sinh x}{\cosh x}$ and use substitution $u=\cosh x$: the answer is $\ln(\cosh x)+C$.
\item Powers: use identities, e.g. $\cosh^{2}x=\dfrac{\cosh 2x+1}{2}$, $\sinh^{2}x=\dfrac{\cosh 2x-1}{2}$, exactly as for trigonometric integrals.
\end{enumerate}
\end{methode}

\begin{exoresolu}[Two standard integrals]
Evaluate (a) $\displaystyle\int_{0}^{1}\cosh 2x\,dx$; \quad (b) $\displaystyle\int\frac{dx}{\sqrt{x^{2}+9}}$, expressing the answer in logarithmic form.
\tcblower
\textbf{Solution.} (a) Since $\dfrac{d}{dx}\sinh 2x=2\cosh 2x$:
\[ \int_{0}^{1}\cosh 2x\,dx=\left[\frac{1}{2}\sinh 2x\right]_{0}^{1}=\frac{1}{2}\bigl(\sinh 2-\sinh 0\bigr)=\frac{1}{2}\sinh 2=\frac{e^{2}-e^{-2}}{4}\approx 1.8134. \]
(b) This matches the standard form with $a=3$:
\[ \int\frac{dx}{\sqrt{x^{2}+9}}=\operatorname{arsinh}\frac{x}{3}+C. \]
Using the logarithmic form $\operatorname{arsinh}u=\ln\!\left(u+\sqrt{u^{2}+1}\right)$:
\[ \operatorname{arsinh}\frac{x}{3}=\ln\!\left(\frac{x}{3}+\sqrt{\frac{x^{2}}{9}+1}\right)=\ln\!\left(\frac{x+\sqrt{x^{2}+9}}{3}\right). \]
Since $\ln 3$ is a constant, it can be absorbed into $C$:
\[ \int\frac{dx}{\sqrt{x^{2}+9}}=\ln\!\left(x+\sqrt{x^{2}+9}\right)+C'. \qquad\blacksquare \]
\end{exoresolu}

\begin{exoresolu}[Integration of powers]
Evaluate $\displaystyle\int\cosh^{2}3x\,dx$.
\tcblower
\textbf{Solution.} Use $\cosh^{2}\theta = \frac{\cosh 2\theta + 1}{2}$ with $\theta = 3x$:
\[ \int\cosh^{2}3x\,dx = \int\frac{\cosh 6x + 1}{2}\,dx = \frac{1}{2}\left(\frac{\sinh 6x}{6} + x\right) + C = \frac{\sinh 6x}{12} + \frac{x}{2} + C. \]
\end{exoresolu}

\coursec{Advanced Applications: Hyperbolic Substitutions and Maclaurin Series}

\begin{methode}[Choosing the right hyperbolic substitution]
Match the form of the integrand to the identity that simplifies it:
\begin{enumerate}
\item $\sqrt{x^{2}+a^{2}}$: put $x=a\sinh t$, because $a^{2}\sinh^{2}t+a^{2}=a^{2}\cosh^{2}t$ and $dx=a\cosh t\,dt$.
\item $\sqrt{x^{2}-a^{2}}$: put $x=a\cosh t$, because $a^{2}\cosh^{2}t-a^{2}=a^{2}\sinh^{2}t$ and $dx=a\sinh t\,dt$ (take $t\ge 0$ so $\sinh t\ge 0$).
\item $a^{2}-x^{2}$ (with $|x|<a$): put $x=a\tanh t$, because $a^{2}-a^{2}\tanh^{2}t=a^{2}\operatorname{sech}^{2}t$ and $dx=a\operatorname{sech}^{2}t\,dt$.
\end{enumerate}
Then: compute $dx$ in terms of $dt$, simplify the square root using the identity, integrate in $t$, and convert back to $x$ using inverse functions or logarithmic forms.
\end{methode}

\begin{exoresolu}[Hyperbolic substitution]
Use the substitution $x=3\sinh t$ to evaluate $\displaystyle\int\sqrt{x^{2}+9}\,dx$.
\tcblower
\textbf{Solution.} With $x=3\sinh t$: $dx=3\cosh t\,dt$ and
\[ \sqrt{x^{2}+9}=\sqrt{9\sinh^{2}t+9}=3\sqrt{\sinh^{2}t+1}=3\cosh t \quad(\cosh t>0). \]
Hence
\[ \int\sqrt{x^{2}+9}\,dx=\int 3\cosh t\cdot 3\cosh t\,dt=9\int\cosh^{2}t\,dt. \]
Use the identity $\cosh^{2}t=\dfrac{\cosh 2t+1}{2}$:
\[ 9\int\cosh^{2}t\,dt=\frac{9}{2}\int(\cosh 2t+1)\,dt=\frac{9}{2}\left(\frac{\sinh 2t}{2}+t\right)+C=\frac{9}{4}\sinh 2t+\frac{9}{2}t+C. \]
Convert back: $t=\operatorname{arsinh}\dfrac{x}{3}$ and
\[ \sinh 2t=2\sinh t\cosh t=2\cdot\frac{x}{3}\cdot\frac{\sqrt{x^{2}+9}}{3}=\frac{2x\sqrt{x^{2}+9}}{9}. \]
Therefore
\[ \int\sqrt{x^{2}+9}\,dx=\frac{x\sqrt{x^{2}+9}}{2}+\frac{9}{2}\operatorname{arsinh}\frac{x}{3}+C
=\frac{x\sqrt{x^{2}+9}}{2}+\frac{9}{2}\ln\!\left(x+\sqrt{x^{2}+9}\right)+C'. \qquad\blacksquare \]
\textbf{Check (differentiate):} $\dfrac{d}{dx}\!\left[\tfrac{x}{2}\sqrt{x^{2}+9}\right]=\tfrac{1}{2}\sqrt{x^{2}+9}+\tfrac{x^{2}}{2\sqrt{x^{2}+9}}=\tfrac{2x^{2}+9}{2\sqrt{x^{2}+9}}$, and $\dfrac{d}{dx}\!\left[\tfrac{9}{2}\operatorname{arsinh}\tfrac{x}{3}\right]=\tfrac{9}{2}\cdot\tfrac{1}{\sqrt{x^{2}+9}}$. Sum: $\tfrac{2x^{2}+9+9}{2\sqrt{x^{2}+9}}=\tfrac{x^{2}+9}{\sqrt{x^{2}+9}}=\sqrt{x^{2}+9}$ \checkmark.
\end{exoresolu}

\begin{exoresolu}[Another substitution]
Evaluate $\displaystyle\int\frac{dx}{\sqrt{x^{2}-16}}$ for $x>4$ using a hyperbolic substitution.
\tcblower
\textbf{Solution.} Put $x=4\cosh t$ ($t>0$). Then $dx=4\sinh t\,dt$ and $\sqrt{x^{2}-16}=\sqrt{16\cosh^{2}t-16}=4\sinh t$.
\[ \int\frac{dx}{\sqrt{x^{2}-16}} = \int\frac{4\sinh t}{4\sinh t}\,dt = \int dt = t + C = \operatorname{arcosh}\frac{x}{4} + C = \ln\!\left(\frac{x}{4}+\sqrt{\frac{x^{2}}{16}-1}\right)+C = \ln\!\left(x+\sqrt{x^{2}-16}\right)+C'. \]
\end{exoresolu}

\begin{methode}[Building a Maclaurin series for hyperbolic functions]
\begin{enumerate}
\item Recall the Maclaurin formula: $f(x)=f(0)+f'(0)x+\dfrac{f''(0)}{2!}x^{2}+\dfrac{f'''(0)}{3!}x^{3}+\cdots$
\item Exploit the \textbf{cyclic derivatives}: for $\sinh$ and $\cosh$, derivatives alternate between the two functions, and at $x=0$: $\sinh 0=0$, $\cosh 0=1$.
\item Hence even powers survive for $\cosh$, odd powers for $\sinh$:
\[ \sinh x = x + \frac{x^{3}}{3!} + \frac{x^{5}}{5!} + \cdots = \sum_{n=0}^{\infty}\frac{x^{2n+1}}{(2n+1)!},\qquad
\cosh x = 1 + \frac{x^{2}}{2!} + \frac{x^{4}}{4!} + \cdots = \sum_{n=0}^{\infty}\frac{x^{2n}}{(2n)!}. \]
\item \textbf{Shortcut:} take the series of $e^{x}=\sum\dfrac{x^{n}}{n!}$ and use the definitions; the odd/even terms separate automatically:
\[ \sinh x = \frac{e^{x}-e^{-x}}{2} = \frac{1}{2}\sum_{n=0}^{\infty}\frac{x^{n}-(-x)^{n}}{n!} = \sum_{n\text{ odd}}\frac{x^{n}}{n!}. \]
\item For approximations, state clearly the order of the neglected terms, e.g. $\cosh x\approx 1+\dfrac{x^{2}}{2}$ for small $x$.
\end{enumerate}
\end{methode}

\begin{exoresolu}[Series expansion and approximation]
(a) Obtain the Maclaurin expansion of $\cosh x$ up to the term in $x^{4}$. (b) Hence show that for small $x$, $\cosh 2x - 1 - 2x^{2}\approx \dfrac{2x^{4}}{3}$, and use this to estimate $\displaystyle\lim_{x\to 0}\frac{\cosh 2x-1}{x^{2}}$.
\tcblower
\textbf{Solution.} (a) Let $f(x)=\cosh x$. The derivatives cycle:
\[ f'(x)=\sinh x,\quad f''(x)=\cosh x,\quad f'''(x)=\sinh x,\quad f^{(4)}(x)=\cosh x. \]
At $x=0$: $f(0)=1$, $f'(0)=0$, $f''(0)=1$, $f'''(0)=0$, $f^{(4)}(0)=1$. Therefore
\[ \cosh x = 1 + \frac{x^{2}}{2!}+\frac{x^{4}}{4!}+\cdots = 1+\frac{x^{2}}{2}+\frac{x^{4}}{24}+\cdots \]
(b) Replace $x$ by $2x$:
\[ \cosh 2x = 1+\frac{(2x)^{2}}{2}+\frac{(2x)^{4}}{24}+\cdots = 1+2x^{2}+\frac{16x^{4}}{24}+\cdots = 1+2x^{2}+\frac{2x^{4}}{3}+\cdots \]
Hence $\cosh 2x-1-2x^{2}\approx \dfrac{2x^{4}}{3}$ for small $x$, as required. For the limit:
\[ \frac{\cosh 2x-1}{x^{2}}=\frac{2x^{2}+\frac{2}{3}x^{4}+\cdots}{x^{2}}=2+\frac{2}{3}x^{2}+\cdots \xrightarrow[x\to 0]{} 2. \]
\[ \boxed{\displaystyle\lim_{x\to 0}\frac{\cosh 2x-1}{x^{2}}=2.} \]
\textbf{Consistency check:} the same result follows from the identity $\cosh 2x-1=2\sinh^{2}x$ and $\sinh x\approx x$, giving $\dfrac{2\sinh^{2}x}{x^{2}}\to 2$ \checkmark.
\end{exoresolu}

\begin{exoresolu}[Maclaurin series of $\tanh x$]
Find the Maclaurin expansion of $\tanh x$ up to the term in $x^{5}$.
\tcblower
\textbf{Solution.} Use $\tanh x = \frac{\sinh x}{\cosh x} = \frac{x + \frac{x^{3}}{6} + \frac{x^{5}}{120} + \cdots}{1 + \frac{x^{2}}{2} + \frac{x^{4}}{24} + \cdots}$.
Perform polynomial long division or use $(1+u)^{-1}\approx 1-u+u^{2}$ with $u=\frac{x^{2}}{2}+\frac{x^{4}}{24}$:
\[ \tanh x = \left(x + \frac{x^{3}}{6} + \frac{x^{5}}{120}\right)\left(1 - \frac{x^{2}}{2} - \frac{x^{4}}{24} + \frac{x^{4}}{4} + \cdots\right) = x - \frac{x^{3}}{3} + \frac{2x^{5}}{15} + \cdots \]
\end{exoresolu}

\begin{attention}[Frequent errors to avoid in GCE examinations]
\begin{itemize}
\item Writing $\dfrac{d}{dx}(\cosh x)=-\sinh x$: the minus sign belongs to $\cos$, \textbf{not} to $\cosh$.
\item Forgetting to reject negative roots of $e^{x}$ when solving hyperbolic equations.
\item Using $\operatorname{arcosh}$ formulas for $x<1$ or $\operatorname{artanh}$ for $|x|\ge 1$: always check domains.
\item Applying Osborn's rule but forgetting to change the sign of \emph{implied} products of sines such as $\tan^{2}x=\dfrac{\sin^{2}x}{\cos^{2}x}$.
\item Confusing $\operatorname{arsinh}$ with $\dfrac{1}{\sinh}$: the inverse function is \textbf{not} the reciprocal.
\item In integration, forgetting the factor $\frac{1}{a}$ when the argument is $ax+b$.
\item In hyperbolic substitutions, forgetting to change the differential $dx$ and the limits (if definite integral).
\item In Maclaurin series, mixing up the factorial denominators: $\sinh$ has odd factorials, $\cosh$ has even factorials.
\end{itemize}
\end{attention}

\begin{savaistu}[Hyperbolic functions in Cameroon and beyond]
The catenary curve $y=a\cosh(x/a)$ appears in the design of suspension bridges (like the Wouri Bridge in Douala) and power lines. The shape of a hanging chain minimizes potential energy. In relativity, rapidity $\phi$ is defined by $\tanh\phi = v/c$, and Lorentz transformations become hyperbolic rotations. The Gateway Arch in St. Louis is a weighted catenary, not a parabola.
\end{savaistu}

\begin{exercice}[Mixed practice --- GCE style]
1. Prove that $\sinh 3x = 3\sinh x + 4\sinh^{3}x$.
2. Solve $\cosh 2x - 3\sinh x = 0$, giving answers in exact logarithmic form.
3. Find $\displaystyle\int_{0}^{\ln 2}\frac{e^{x}}{\sqrt{e^{2x}+4}}\,dx$.
4. Use the substitution $x=2\tanh t$ to evaluate $\displaystyle\int\frac{dx}{(4-x^{2})^{3/2}}$.
5. Obtain the Maclaurin series of $\sinh x \cosh x$ up to $x^{5}$ and hence deduce the series for $\sinh 2x$.
\end{exercice}

\begin{corrige}[Exercise 3]
1. $\sinh 3x = \sinh(2x+x) = \sinh 2x\cosh x + \cosh 2x\sinh x = (2\sinh x\cosh x)\cosh x + (1+2\sinh^{2}x)\sinh x = 2\sinh x(1+\sinh^{2}x) + \sinh x + 2\sinh^{3}x = 3\sinh x + 4\sinh^{3}x$.
2. $\cosh 2x = 1+2\sinh^{2}x$, so $1+2\sinh^{2}x - 3\sinh x = 0 \Rightarrow (2\sinh x - 1)(\sinh x - 1) = 0$. $\sinh x = \frac{1}{2}$ or $1$. $x = \operatorname{arsinh}\frac{1}{2} = \ln\!\left(\frac{1+\sqrt{5}}{2}\right)$ or $x = \operatorname{arsinh}1 = \ln(1+\sqrt{2})$.
3. Substitute $u=e^{x}$, $du=e^{x}dx$. Limits: $x=0\to u=1$, $x=\ln 2\to u=2$. $\int_{1}^{2}\frac{du}{\sqrt{u^{2}+4}} = \left[\operatorname{arsinh}\frac{u}{2}\right]_{1}^{2} = \operatorname{arsinh}1 - \operatorname{arsinh}\frac{1}{2} = \ln(1+\sqrt{2}) - \ln\!\left(\frac{1+\sqrt{5}}{2}\right) = \ln\!\left(\frac{2(1+\sqrt{2})}{1+\sqrt{5}}\right)$.
4. $x=2\tanh t$, $dx=2\operatorname{sech}^{2}t\,dt$, $4-x^{2}=4\operatorname{sech}^{2}t$, $(4-x^{2})^{3/2}=8\operatorname{sech}^{3}t$. Integral $= \int\frac{2\operatorname{sech}^{2}t}{8\operatorname{sech}^{3}t}\,dt = \frac{1}{4}\int\cosh t\,dt = \frac{1}{4}\sinh t + C$. Since $\tanh t = x/2$, $\sinh t = \frac{x/2}{\sqrt{1-x^{2}/4}} = \frac{x}{\sqrt{4-x^{2}}}$. Answer: $\frac{x}{4\sqrt{4-x^{2}}}+C$.
5. $\sinh x\cosh x = \left(x+\frac{x^{3}}{6}+\frac{x^{5}}{120}\right)\left(1+\frac{x^{2}}{2}+\frac{x^{4}}{24}\right) = x + \frac{2x^{3}}{3} + \frac{17x^{5}}{120} + \cdots$. But $\sinh x\cosh x = \frac{1}{2}\sinh 2x$, so $\sinh 2x = 2x + \frac{4x^{3}}{3} + \frac{17x^{5}}{60} + \cdots$. Check: direct substitution $2x$ into $\sinh$ series gives $2x + \frac{8x^{3}}{6} + \frac{32x^{5}}{120} = 2x + \frac{4x^{3}}{3} + \frac{4x^{5}}{15} = 2x + \frac{4x^{3}}{3} + \frac{16x^{5}}{60}$. Discrepancy? Re-check multiplication: $\frac{1}{6}\cdot\frac{1}{2}=\frac{1}{12}$, $\frac{1}{120}\cdot1=\frac{1}{120}$, $\frac{1}{6}\cdot\frac{1}{24}=\frac{1}{144}$... Better: $\sinh 2x = 2x + \frac{(2x)^{3}}{3!} + \frac{(2x)^{5}}{5!} = 2x + \frac{8x^{3}}{6} + \frac{32x^{5}}{120} = 2x + \frac{4x^{3}}{3} + \frac{4x^{5}}{15}$. So $\sinh x\cosh x = x + \frac{2x^{3}}{3} + \frac{2x^{5}}{15} + \cdots$. The product expansion must yield this. Let's do carefully: $(x + x^3/6 + x^5/120)(1 + x^2/2 + x^4/24) = x + x^3/2 + x^5/24 + x^3/6 + x^5/12 + x^5/120 = x + (1/2+1/6)x^3 + (1/24+1/12+1/120)x^5 = x + (2/3)x^3 + (5/120+10/120+1/120)x^5 = x + 2x^3/3 + 16x^5/120 = x + 2x^3/3 + 2x^5/15$. Correct.
\end{corrige}

\newpage

\coursec{Exercises}

\courssub{I check my knowledge}

\begin{exercice}[Multiple choice]
For each question, choose the single correct answer, justifying briefly.
\begin{enumerate}
\item $\cosh x$ is equal to:
\begin{enumerate}
\item $\dfrac{e^x - e^{-x}}{2}$
\item $\dfrac{e^x + e^{-x}}{2}$
\item $\dfrac{e^x + e^{-x}}{e^x - e^{-x}}$
\item $\dfrac{2}{e^x + e^{-x}}$
\end{enumerate}
\item The domain of the function $x \mapsto \operatorname{arcosh} x$ is:
\begin{enumerate}
\item $\mathbb{R}$
\item $]0, +\infty[$
\item $[1, +\infty[$
\item $]-1, 1[$
\end{enumerate}
\item $\cosh^2 x - \sinh^2 x$ equals:
\begin{enumerate}
\item $-1$
\item $0$
\item $1$
\item $\cosh 2x$
\end{enumerate}
\item $\displaystyle\int \sinh x\, dx$ equals:
\begin{enumerate}
\item $\cosh x + C$
\item $-\cosh x + C$
\item $\operatorname{sech}^2 x + C$
\item $\ln(\cosh x) + C$
\end{enumerate}
\end{enumerate}
\end{exercice}

\begin{exercice}[True or false]
State whether each assertion is true or false, and justify your answer carefully.
\begin{enumerate}
\item The function $\sinh$ is an even function.
\item For all real $x$, $\cosh x \geq 1$.
\item $\operatorname{arsinh} x = \ln\left(x + \sqrt{x^2 + 1}\right)$ for all real $x$.
\item $\dfrac{d}{dx}(\tanh x) = 1 + \tanh^2 x$.
\end{enumerate}
\end{exercice}

\begin{exercice}[Definitions]
\begin{enumerate}
\item Give the definitions of $\sinh x$, $\cosh x$ and $\tanh x$ in terms of the exponential function.
\item State Osborn's rule and illustrate it by writing the hyperbolic identity corresponding to $\cos 2x = 1 - 2\sin^2 x$.
\item Define $\operatorname{artanh} x$ and state its domain.
\end{enumerate}
\end{exercice}

\begin{exercice}[Direct calculations]
\begin{enumerate}
\item Given that $\sinh x = \dfrac{3}{4}$, find the exact values of $\cosh x$ and $\tanh x$.
\item Evaluate $\cosh(\ln 2)$ and $\sinh(\ln 3)$ exactly.
\item Differentiate $f(x) = \sinh(2x) + \cosh^2 x$ with respect to $x$.
\item Evaluate $\displaystyle\int_0^{\ln 2} \cosh x\, dx$, giving your answer as a rational number.
\end{enumerate}
\end{exercice}

\courssub{I practise}

\begin{exercice}[Proving identities]
\begin{enumerate}
\item Prove that $\cosh 2x = 2\cosh^2 x - 1$.
\item Prove that $\sinh(x + y) = \sinh x \cosh y + \cosh x \sinh y$.
\item Hence express $\sinh 2x$ in terms of $\sinh x$ and $\cosh x$.
\end{enumerate}
\end{exercice}

\begin{exercice}[Solving a hyperbolic equation]
Solve for $x$, giving your answers in exact logarithmic form:
\[ 3\sinh x - 4\cosh x = 5. \]
\end{exercice}

\begin{exercice}[Logarithmic form of an inverse function]
\begin{enumerate}
\item Starting from $y = \operatorname{arcosh} x$, show that $e^{2y} - 2xe^y + 1 = 0$.
\item Hence prove that $\operatorname{arcosh} x = \ln\left(x + \sqrt{x^2 - 1}\right)$ for $x \geq 1$, explaining why the other root of the quadratic is rejected.
\item Use this result to evaluate $\operatorname{arcosh} 2$ in logarithmic form.
\end{enumerate}
\end{exercice}

\begin{exercice}[Integration by hyperbolic substitution]
\begin{enumerate}
\item Using the substitution $x = 3\cosh t$, evaluate
\[ \int \frac{dx}{\sqrt{x^2 - 9}}, \qquad x > 3, \]
and express the result in terms of an inverse hyperbolic function.
\item Using a suitable hyperbolic substitution, evaluate $\displaystyle\int \frac{dx}{\sqrt{x^2 + 4}}$.
\end{enumerate}
\end{exercice}

\courssub{I prepare for the GCE}

\begin{exercice}[Maclaurin series and approximation \hfill (14 marks)]
\begin{enumerate}
\item[(a)] Write down the Maclaurin expansions of $\sinh x$ and $\cosh x$ up to the term in $x^5$. \hfill (4 marks)
\item[(b)] Hence, or otherwise, find the Maclaurin expansion of $\tanh x$ up to the term in $x^5$. \hfill (5 marks)
\item[(c)] Use your expansion to estimate $\tanh 0.2$, giving your answer to 5 decimal places. \hfill (3 marks)
\item[(d)] Given that the true value of $\tanh 0.2$ is $0.197375$, comment on the accuracy of your estimate. \hfill (2 marks)
\end{enumerate}
\end{exercice}

\begin{exercice}[Calculus with hyperbolic functions \hfill (16 marks)]
\begin{enumerate}
\item[(a)] Using the substitution $u = e^x$, prove that
\[ \int \operatorname{sech} x\, dx = 2\arctan(e^x) + C. \] \hfill (5 marks)
\item[(b)] Show that $\dfrac{d}{dx}\left(\operatorname{artanh} x\right) = \dfrac{1}{1 - x^2}$ for $|x| < 1$. \hfill (4 marks)
\item[(c)] Hence evaluate $\displaystyle\int_0^{1/2} \frac{dx}{1 - x^2}$, giving your answer in exact logarithmic form. \hfill (3 marks)
\item[(d)] Solve the equation $\operatorname{artanh} x = \ln 2$, giving $x$ as an exact fraction. \hfill (4 marks)
\end{enumerate}
\end{exercice}

\begin{exercice}[The catenary: a hanging chain \hfill (18 marks)]
A uniform cable hanging between two pylons in Douala is modelled by the curve
\[ y = 10\cosh\left(\frac{x}{10}\right), \qquad -10 \leq x \leq 10, \]
where $x$ and $y$ are measured in metres.
\begin{enumerate}
\item[(a)] Show that $1 + \left(\dfrac{dy}{dx}\right)^2 = \cosh^2\left(\dfrac{x}{10}\right)$. \hfill (4 marks)
\item[(b)] The arc length of a curve $y = f(x)$ between $x = a$ and $x = b$ is given by
\[ L = \int_a^b \sqrt{1 + \left(\frac{dy}{dx}\right)^2}\, dx. \]
Show that the length of the cable between $x = -10$ and $x = 10$ is $20\sinh(1)$ metres, and find this length correct to 2 decimal places. \hfill (6 marks)
\item[(c)] Find the height of the cable above its lowest point at the pylons $x = \pm 10$, giving your answer in terms of $\cosh 1$. \hfill (3 marks)
\item[(d)] Sketch the curve, indicating its axis of symmetry, its minimum point and its behaviour as $|x|$ increases. \hfill (5 marks)
\end{enumerate}
\end{exercice}

\newpage

\coursec{Solutions to the exercises}

\begin{corrige}[Exercise 1 — Multiple choice]
\begin{enumerate}
\item \textbf{Answer (b).} By definition, $\cosh x = \dfrac{e^x + e^{-x}}{2}$. Option (a) is $\sinh x$, option (c) is $\coth x$, option (d) is $\operatorname{sech} x$.
\item \textbf{Answer (c).} Since $\cosh : [0, +\infty[ \to [1, +\infty[$ is a bijection, its inverse $\operatorname{arcosh}$ has domain $[1, +\infty[$.
\item \textbf{Answer (c).} The fundamental identity gives $\cosh^2 x - \sinh^2 x = 1$ for all real $x$. (Note: $\cosh 2x = \cosh^2 x + \sinh^2 x$, so (d) is wrong.)
\item \textbf{Answer (a).} Since $\dfrac{d}{dx}(\cosh x) = \sinh x$, an antiderivative of $\sinh x$ is $\cosh x$, so $\displaystyle\int \sinh x\, dx = \cosh x + C$.
\end{enumerate}
\end{corrige}

\begin{corrige}[Exercise 2 — True or false]
\begin{enumerate}
\item \textbf{False.} $\sinh(-x) = \dfrac{e^{-x} - e^{x}}{2} = -\sinh x$, so $\sinh$ is \emph{odd}, not even. (It is $\cosh$ that is even.)
\item \textbf{True.} $\cosh x - 1 = \dfrac{e^x + e^{-x} - 2}{2} = \dfrac{(e^{x/2} - e^{-x/2})^2}{2} \geq 0$, with equality only at $x = 0$. Hence $\cosh x \geq 1$ for all $x \in \mathbb{R}$.
\item \textbf{True.} Setting $y = \operatorname{arsinh} x$, we have $x = \sinh y$, so $e^{2y} - 2xe^y - 1 = 0$, giving $e^y = x + \sqrt{x^2+1}$ (the other root $x - \sqrt{x^2+1} < 0$ is rejected since $e^y > 0$). Hence $y = \ln\!\left(x + \sqrt{x^2+1}\right)$, valid for all $x \in \mathbb{R}$.
\item \textbf{False.} $\dfrac{d}{dx}(\tanh x) = \operatorname{sech}^2 x = 1 - \tanh^2 x$, \emph{not} $1 + \tanh^2 x$ (that is the \emph{trigonometric} formula for $\tan$).
\end{enumerate}
\end{corrige}

\begin{corrige}[Exercise 3 — Definitions]
\begin{enumerate}
\item $\sinh x = \dfrac{e^x - e^{-x}}{2}$, \quad $\cosh x = \dfrac{e^x + e^{-x}}{2}$, \quad $\tanh x = \dfrac{\sinh x}{\cosh x} = \dfrac{e^x - e^{-x}}{e^x + e^{-x}}$.
\item \textbf{Osborn's rule:} to obtain a hyperbolic identity from a trigonometric one, replace each trigonometric function by the corresponding hyperbolic function ($\sin \to \sinh$, $\cos \to \cosh$, $\tan \to \tanh$), and change the sign of every product (or implied product, such as a square) of two sines. Applying this to $\cos 2x = 1 - 2\sin^2 x$: the term $\sin^2 x$ is a product of two sines, so its sign changes:
\[ \cosh 2x = 1 + 2\sinh^2 x. \]
\item $\operatorname{artanh} x$ is the inverse function of $\tanh : \mathbb{R} \to ]-1, 1[$: for $x \in ]-1, 1[$, $y = \operatorname{artanh} x \iff x = \tanh y$. Its domain is $]-1, 1[$.
\end{enumerate}
\end{corrige}

\begin{corrige}[Exercise 4 — Direct calculations]
\begin{enumerate}
\item From $\cosh^2 x - \sinh^2 x = 1$: $\cosh^2 x = 1 + \dfrac{9}{16} = \dfrac{25}{16}$, so $\cosh x = \dfrac{5}{4}$ (positive root since $\cosh x \geq 1$). Then $\tanh x = \dfrac{\sinh x}{\cosh x} = \dfrac{3/4}{5/4} = \dfrac{3}{5}$.
\item $\cosh(\ln 2) = \dfrac{2 + \tfrac{1}{2}}{2} = \dfrac{5}{4}$; \quad $\sinh(\ln 3) = \dfrac{3 - \tfrac{1}{3}}{2} = \dfrac{4}{3}$.
\item Let $f(x) = \sinh 2x + \cosh^2 x$. Then
\[ f'(x) = 2\cosh(2x) + 2\cosh x \sinh x = 2\cosh 2x + \sinh 2x \quad (\text{using } 2\sinh x\cosh x = \sinh 2x). \]
\item $\displaystyle\int_0^{\ln 2} \cosh x\, dx = \big[\sinh x\big]_0^{\ln 2} = \sinh(\ln 2) - 0 = \frac{2 - \tfrac{1}{2}}{2} = \frac{3}{4}$.
\end{enumerate}
\end{corrige}

\begin{corrige}[Exercise 5 — Proving identities]
\begin{enumerate}
\item Starting from the right-hand side:
\begin{align*}
2\cosh^2 x - 1 &= 2\left(\frac{e^x + e^{-x}}{2}\right)^2 - 1 = \frac{e^{2x} + 2 + e^{-2x}}{2} - 1 = \frac{e^{2x} + e^{-2x}}{2} = \cosh 2x. \qquad \blacksquare
\end{align*}
\item From the definitions:
\begin{align*}
\sinh x \cosh y + \cosh x \sinh y &= \frac{(e^x - e^{-x})(e^y + e^{-y}) + (e^x + e^{-x})(e^y - e^{-y})}{4} \\
&= \frac{2e^{x+y} - 2e^{-(x+y)}}{4} = \frac{e^{x+y} - e^{-(x+y)}}{2} = \sinh(x+y). \qquad \blacksquare
\end{align*}
\item Setting $y = x$ in part (b): $\sinh 2x = \sinh x \cosh x + \cosh x \sinh x = 2\sinh x \cosh x$.
\end{enumerate}
\end{corrige}

\begin{corrige}[Exercise 6 — Solving a hyperbolic equation]
Solve $3\sinh x - 4\cosh x = 5$ for $x \in \mathbb{R}$.

Write in exponential form:
\[ 3\cdot\frac{e^x - e^{-x}}{2} - 4\cdot\frac{e^x + e^{-x}}{2} = 5 \;\Longrightarrow\; -e^x - 7e^{-x} = 10. \]
Multiplying by $e^x$ (with $u = e^x > 0$):
\[ -u^2 - 7 = 10u \;\Longrightarrow\; u^2 + 10u + 7 = 0. \]
The discriminant is $\Delta = 100 - 28 = 72$, so
\[ u = \frac{-10 \pm \sqrt{72}}{2} = -5 \pm 3\sqrt{2}. \]
Since $3\sqrt{2} \approx 4.243 < 5$, both roots $-5 + 3\sqrt{2}$ and $-5 - 3\sqrt{2}$ are strictly negative. But $u = e^x > 0$ for all real $x$. Hence the quadratic has no admissible root, and the original equation has \textbf{no real solution}.

\textbf{Verification (optional):} The function $f(x) = 3\sinh x - 4\cosh x$ has derivative $f'(x) = 3\cosh x - 4\sinh x$. Setting $f'(x)=0$ gives $\tanh x = \frac{3}{4}$, so $\sinh x = \frac{3}{\sqrt{7}}$, $\cosh x = \frac{4}{\sqrt{7}}$ (using $\cosh^2 x - \sinh^2 x = 1$). The maximum value is $f_{\max} = \frac{9}{\sqrt{7}} - \frac{16}{\sqrt{7}} = -\sqrt{7} \approx -2.65$. Since $f_{\max} < 5$, the equation $f(x)=5$ is indeed impossible in $\mathbb{R}$.

\textbf{Conclusion:} The solution set is $\varnothing$.
\end{corrige}

\begin{corrige}[Exercise 7 — Logarithmic form of an inverse function]
\begin{enumerate}
\item $y = \operatorname{arcosh} x \iff x = \cosh y = \dfrac{e^y + e^{-y}}{2}$. Multiplying by $2e^y$: $2xe^y = e^{2y} + 1$, i.e. $e^{2y} - 2xe^y + 1 = 0$. $\blacksquare$
\item This is a quadratic in $e^y$: $e^y = \dfrac{2x \pm \sqrt{4x^2 - 4}}{2} = x \pm \sqrt{x^2 - 1}$. The two roots are positive and their product is $(x + \sqrt{x^2-1})(x - \sqrt{x^2-1}) = 1$, so $x - \sqrt{x^2 - 1} = \dfrac{1}{x + \sqrt{x^2-1}} \leq 1$, giving $y = \ln(x - \sqrt{x^2-1}) \leq 0$. Since $\operatorname{arcosh}$ takes values in $[0, +\infty[$ (we chose the principal branch $y \geq 0$), we reject this root and keep
\[ \operatorname{arcosh} x = \ln\!\left(x + \sqrt{x^2 - 1}\right), \qquad x \geq 1. \qquad \blacksquare \]
\item $\operatorname{arcosh} 2 = \ln\!\left(2 + \sqrt{4 - 1}\right) = \ln(2 + \sqrt{3})$.
\end{enumerate}
\end{corrige}

\begin{corrige}[Exercise 8 — Integration by hyperbolic substitution]
\begin{enumerate}
\item Let $x = 3\cosh t$, $t > 0$ (valid since $x > 3$). Then $dx = 3\sinh t\, dt$ and $\sqrt{x^2 - 9} = \sqrt{9\cosh^2 t - 9} = 3\sinh t$ (positive as $t > 0$). Hence
\[ \int \frac{dx}{\sqrt{x^2 - 9}} = \int \frac{3\sinh t}{3\sinh t}\, dt = \int dt = t + C = \operatorname{arcosh}\!\left(\frac{x}{3}\right) + C. \]
\item The form $\sqrt{x^2 + a^2}$ calls for a \emph{sinh} substitution: let $x = 2\sinh t$, $dx = 2\cosh t\, dt$, $\sqrt{x^2 + 4} = 2\cosh t$. Then
\[ \int \frac{dx}{\sqrt{x^2 + 4}} = \int \frac{2\cosh t}{2\cosh t}\, dt = t + C = \operatorname{arsinh}\!\left(\frac{x}{2}\right) + C = \ln\!\left(\frac{x + \sqrt{x^2+4}}{2}\right) + C. \]
\end{enumerate}
\end{corrige}

\begin{corrige}[Exercise 9 — Maclaurin series and approximation \hfill (14 marks)]
\textbf{(a) (4 marks)} Since $\dfrac{d}{dx}\sinh x = \cosh x$ and $\dfrac{d}{dx}\cosh x = \sinh x$, with $\sinh 0 = 0$, $\cosh 0 = 1$, the Maclaurin series are:
\[ \sinh x = x + \frac{x^3}{3!} + \frac{x^5}{5!} + \cdots = x + \frac{x^3}{6} + \frac{x^5}{120} + \cdots \quad\text{(M1 A1)} \]
\[ \cosh x = 1 + \frac{x^2}{2!} + \frac{x^4}{4!} + \cdots = 1 + \frac{x^2}{2} + \frac{x^4}{24} + \cdots \quad\text{(M1 A1)} \]

\textbf{(b) (5 marks)} Write $\tanh x = \dfrac{\sinh x}{\cosh x}$ and use long division (or set $\tanh x = a_1 x + a_3 x^3 + a_5 x^5$, since $\tanh$ is odd, and solve $\sinh x = \cosh x \cdot \tanh x$):
\[ x + \frac{x^3}{6} + \frac{x^5}{120} = \left(1 + \frac{x^2}{2} + \frac{x^4}{24}\right)\left(a_1 x + a_3 x^3 + a_5 x^5\right). \]
Comparing coefficients: $a_1 = 1$; $\dfrac{a_1}{2} + a_3 = \dfrac{1}{6} \Rightarrow a_3 = -\dfrac{1}{3}$; $\dfrac{a_1}{24} + \dfrac{a_3}{2} + a_5 = \dfrac{1}{120} \Rightarrow a_5 = \dfrac{1}{120} - \dfrac{1}{24} + \dfrac{1}{6} = \dfrac{2}{15}$. (M1 M1 A1 A1 A1)
\[ \tanh x = x - \frac{x^3}{3} + \frac{2x^5}{15} + \cdots \]

\textbf{(c) (3 marks)} With $x = 0.2$: (M1 for substitution)
\[ \tanh 0.2 \approx 0.2 - \frac{0.008}{3} + \frac{2(0.00032)}{15} = 0.2 - 0.0026667 + 0.0000427 = 0.19738 \text{ (5 d.p.)}. \quad\text{(A1 A1)} \]

\textbf{(d) (2 marks)} The estimate $0.19738$ agrees with the true value $0.197375$ to 5 decimal places (error $\approx 10^{-6}$). (B1) The approximation is excellent because $|x| = 0.2$ is small, so the neglected term in $x^7$ is of order $10^{-7}$. (B1)
\end{corrige}

\begin{corrige}[Exercise 10 — Calculus with hyperbolic functions \hfill (16 marks)]
\textbf{(a) (5 marks)} $\operatorname{sech} x = \dfrac{2}{e^x + e^{-x}} = \dfrac{2e^x}{e^{2x} + 1}$. Let $u = e^x$, $du = e^x dx$: (M1 M1)
\[ \int \operatorname{sech} x\, dx = \int \frac{2e^x}{e^{2x}+1}\, dx = \int \frac{2\, du}{u^2 + 1} = 2\arctan u + C = 2\arctan(e^x) + C. \quad\text{(M1 A1 A1)} \;\blacksquare \]

\textbf{(b) (4 marks)} $y = \operatorname{artanh} x \iff x = \tanh y$. Differentiating with respect to $y$: $\dfrac{dx}{dy} = \operatorname{sech}^2 y = 1 - \tanh^2 y = 1 - x^2$. (M1 A1) Hence $\dfrac{dy}{dx} = \dfrac{1}{1 - x^2}$, valid for $|x| < 1$ (where $1 - x^2 \neq 0$). (M1 A1) $\blacksquare$

\textbf{(c) (3 marks)} From (b), an antiderivative of $\dfrac{1}{1-x^2}$ is $\operatorname{artanh} x$: (M1)
\[ \int_0^{1/2} \frac{dx}{1 - x^2} = \big[\operatorname{artanh} x\big]_0^{1/2} = \operatorname{artanh}\!\frac{1}{2} = \frac{1}{2}\ln\!\left(\frac{1 + \tfrac{1}{2}}{1 - \tfrac{1}{2}}\right) = \frac{1}{2}\ln 3. \quad\text{(M1 A1)} \]

\textbf{(d) (4 marks)} $\operatorname{artanh} x = \dfrac{1}{2}\ln\!\dfrac{1+x}{1-x} = \ln 2 \Rightarrow \ln\dfrac{1+x}{1-x} = \ln 4 \Rightarrow \dfrac{1+x}{1-x} = 4$. (M1 A1) Then $1 + x = 4 - 4x$, so $5x = 3$ and $x = \dfrac{3}{5}$. (M1 A1) Check: $\tfrac{3}{5} \in {}]{-1},1[$, valid.
\end{corrige}

\begin{corrige}[Exercise 11 — The catenary: a hanging chain \hfill (18 marks)]
\textbf{(a) (4 marks)} $y = 10\cosh\!\left(\dfrac{x}{10}\right) \Rightarrow \dfrac{dy}{dx} = 10 \cdot \dfrac{1}{10}\sinh\!\left(\dfrac{x}{10}\right) = \sinh\!\left(\dfrac{x}{10}\right)$. (M1 A1) Then
\[ 1 + \left(\frac{dy}{dx}\right)^2 = 1 + \sinh^2\!\left(\frac{x}{10}\right) = \cosh^2\!\left(\frac{x}{10}\right), \]
by the fundamental identity $\cosh^2 u - \sinh^2 u = 1$. (M1 A1) $\blacksquare$

\textbf{(b) (6 marks)} Since $\cosh\!\left(\tfrac{x}{10}\right) > 0$, $\sqrt{1 + (y')^2} = \cosh\!\left(\tfrac{x}{10}\right)$. (M1) Hence
\[ L = \int_{-10}^{10} \cosh\!\left(\frac{x}{10}\right) dx = \Big[10\sinh\!\left(\tfrac{x}{10}\right)\Big]_{-10}^{10} = 10\sinh 1 - 10\sinh(-1) = 20\sinh 1, \quad\text{(M1 M1 A1)} \]
using $\sinh(-1) = -\sinh 1$. Numerically, $\sinh 1 = \dfrac{e - e^{-1}}{2} \approx 1.17520$, so $L \approx 23.50\ \mathrm{m}$ (2 d.p.). (M1 A1)

\textbf{(c) (3 marks)} The lowest point is at $x = 0$: $y(0) = 10\cosh 0 = 10\ \mathrm{m}$. (B1) At the pylons, $y(\pm 10) = 10\cosh 1$. (B1) The height above the lowest point is $10\cosh 1 - 10 = 10(\cosh 1 - 1) \approx 5.43\ \mathrm{m}$. (B1)

\textbf{(d) (5 marks)} Expected features: (B1 each)
\begin{itemize}
\item curve symmetric about the $y$-axis (since $\cosh$ is even);
\item minimum point $(0, 10)$ clearly marked;
\item curve passing through $(\pm 10, 10\cosh 1) \approx (\pm 10, 15.43)$;
\item $y \to +\infty$ rapidly (exponentially) as $|x|$ increases;
\item smooth U-shape, always above $y = 10$.
\end{itemize}
\begin{popfigure}
\begin{center}
\begin{tikzpicture}
\begin{axis}[width=10cm, height=6.5cm, axis lines=middle, xlabel={$x$ (m)}, ylabel={$y$ (m)}, xmin=-13, xmax=13, ymin=0, ymax=20, xtick={-10,0,10}, ytick={10}, domain=-10:10, samples=80]
\addplot[popblue, very thick] {10*cosh(x/10)};
\addplot[popmuted, dashed] coordinates {(-10,0) (-10,15.43)};
\addplot[popmuted, dashed] coordinates {(10,0) (10,15.43)};
\node[popblue] at (axis cs:6,11.2) {$y=10\cosh(x/10)$};
\node[popdark] at (axis cs:0,8.6) {$(0,10)$};
\end{axis}
\end{tikzpicture}
\end{center}
\legende{The catenary $y = 10\cosh(x/10)$: symmetric about the $y$-axis, minimum at $(0,10)$.}
\end{popfigure}
\end{corrige}

\newpage

\coursec{Summary sheet}

\begin{popfigure}
\begin{tikzpicture}[
  grow cyclic,
  mindmap,
  concept color=popblue,
  every node/.style={concept, font=\small},
  root concept/.append style={font=\bfseries\small, text=white},
  level 1/.append style={level distance=3.4cm, sibling angle=51.4},
  level 2/.append style={level distance=2.1cm, sibling angle=40, font=\scriptsize},
  text width=2.2cm, align=center
]
\node[root concept]{Hyperbolic \& Inverse Hyperbolic Functions}
  child[concept color=poporange]{ node {Definitions $\cosh x,\ \sinh x,\ \tanh x$}
    child { node {Exponential forms} }
    child { node {Numerical values} }
  }
  child[concept color=popgreen]{ node {Graphs}
    child { node {$y=\cosh x$: catenary} }
    child { node {$y=\sinh x$: odd} }
  }
  child[concept color=poppurple]{ node {Identities \& Osborn's rule}
    child { node {$\cosh^2x-\sinh^2x=1$} }
  }
  child[concept color=poppink]{ node {Inverse functions}
    child { node {Log forms e.g.\ $\operatorname{arsinh}x$} }
    child { node {Graphs \& domains} }
  }
  child[concept color=popgold]{ node {Hyperbolic equations}
    child { node {Solve via $e^x$} }
  }
  child[concept color=popblue!70]{ node {Calculus}
    child { node {Derivatives} }
    child { node {Integrals} }
    child { node {Substitutions} }
  }
  child[concept color=popgreen!70!black]{ node {Maclaurin series}
    child { node {$\sinh x$, $\cosh x$} }
  };
\end{tikzpicture}
\legende{Mind map of the chapter: hyperbolic and inverse hyperbolic functions.}
\end{popfigure}

\begin{bilan}
\textbf{Key definitions.}
\begin{itemize}
\item $\cosh x=\dfrac{e^{x}+e^{-x}}{2}$ (even), \quad $\sinh x=\dfrac{e^{x}-e^{-x}}{2}$ (odd), \quad $\tanh x=\dfrac{\sinh x}{\cosh x}=\dfrac{e^{2x}-1}{e^{2x}+1}$.
\item $\operatorname{sech} x=\dfrac{1}{\cosh x}$, \quad $\operatorname{cosech} x=\dfrac{1}{\sinh x}$, \quad $\coth x=\dfrac{1}{\tanh x}$.
\item Special values: $\cosh 0=1$, $\sinh 0=0$, $\tanh 0=0$; $\cosh x\ge 1$ for all $x$; $-1<\tanh x<1$.
\item Inverse functions: $y=\operatorname{arsinh} x \iff x=\sinh y$ (domain $\mathbb{R}$); $y=\operatorname{arcosh} x \iff x=\cosh y$, $y\ge 0$ (domain $x\ge 1$); $y=\operatorname{artanh} x \iff x=\tanh y$ (domain $|x|<1$).
\end{itemize}

\textbf{Formulas to know by heart.}
\begin{itemize}
\item Fundamental identity: $\cosh^{2}x-\sinh^{2}x=1$; hence $1-\tanh^{2}x=\operatorname{sech}^{2}x$ and $\coth^{2}x-1=\operatorname{cosech}^{2}x$.
\item Addition: $\sinh(x\pm y)=\sinh x\cosh y\pm\cosh x\sinh y$; $\cosh(x\pm y)=\cosh x\cosh y\pm\sinh x\sinh y$.
\item Double angle: $\cosh 2x=\cosh^{2}x+\sinh^{2}x=2\cosh^{2}x-1=1+2\sinh^{2}x$; $\sinh 2x=2\sinh x\cosh x$.
\item \cle{Osborn's rule}: take any trigonometric identity, replace $\cos\to\cosh$, $\sin\to\sinh$, and change the sign of every product of two sines (or implied product, e.g.\ $\tan^{2}$, $\sin^{2}$).
\item Logarithmic forms:
$\operatorname{arsinh} x=\ln\!\left(x+\sqrt{x^{2}+1}\right)$, all $x$;
$\operatorname{arcosh} x=\ln\!\left(x+\sqrt{x^{2}-1}\right)$, $x\ge 1$;
$\operatorname{artanh} x=\dfrac{1}{2}\ln\!\left(\dfrac{1+x}{1-x}\right)$, $|x|<1$.
\end{itemize}

\textbf{Calculus results.}
\begin{itemize}
\item Derivatives: $(\sinh x)'=\cosh x$; $(\cosh x)'=\sinh x$ (no minus sign!); $(\tanh x)'=\operatorname{sech}^{2}x$.
\item Inverse derivatives: $(\operatorname{arsinh} x)'=\dfrac{1}{\sqrt{x^{2}+1}}$; $(\operatorname{arcosh} x)'=\dfrac{1}{\sqrt{x^{2}-1}}$ ($x>1$); $(\operatorname{artanh} x)'=\dfrac{1}{1-x^{2}}$ ($|x|<1$).
\item Integrals: $\displaystyle\int\sinh x\,dx=\cosh x+C$; $\displaystyle\int\cosh x\,dx=\sinh x+C$; $\displaystyle\int\operatorname{sech}^{2}x\,dx=\tanh x+C$.
\item Standard integral forms:
$\displaystyle\int\frac{dx}{\sqrt{x^{2}+a^{2}}}=\operatorname{arsinh}\frac{x}{a}+C$;
$\displaystyle\int\frac{dx}{\sqrt{x^{2}-a^{2}}}=\operatorname{arcosh}\frac{x}{a}+C$ ($x>a$);
$\displaystyle\int\frac{dx}{a^{2}-x^{2}}=\frac{1}{a}\operatorname{artanh}\frac{x}{a}+C$ ($|x|<a$).
\item Maclaurin series (valid for all $x$):
$\cosh x=1+\dfrac{x^{2}}{2!}+\dfrac{x^{4}}{4!}+\cdots$; \quad
$\sinh x=x+\dfrac{x^{3}}{3!}+\dfrac{x^{5}}{5!}+\cdots$.
\end{itemize}

\textbf{Methods.}
\begin{itemize}
\item To evaluate numerically: rewrite in terms of $e^{x}$ and use the calculator's $e^{x}$ key.
\item To solve a hyperbolic equation: either use identities to reduce to one function, or set $u=e^{x}$ to obtain a quadratic in $u$; keep only $u>0$, then $x=\ln u$.
\item To prove an identity: start from the exponential definitions, or quote Osborn's rule applied to the corresponding trig identity.
\item For $\int dx/\sqrt{x^{2}+a^{2}}$ substitute $x=a\sinh u$; for $\int dx/\sqrt{x^{2}-a^{2}}$ substitute $x=a\cosh u$; for $\int dx/(a^{2}-x^{2})$ substitute $x=a\tanh u$.
\item For Maclaurin expansions: differentiate repeatedly at $0$, or add/subtract the series of $e^{x}$ and $e^{-x}$.
\end{itemize}

\textbf{Common mistakes.}
\begin{itemize}
\item Writing $\cosh^{2}x+\sinh^{2}x=1$: the correct identity has a \emph{minus} sign.
\item Putting a minus sign in $(\cosh x)'$: it is $+\sinh x$.
\item Forgetting the domains: $\operatorname{arcosh} x$ needs $x\ge 1$; $\operatorname{artanh} x$ needs $|x|<1$.
\item Taking the negative root in $\operatorname{arcosh}$: by definition $\operatorname{arcosh} x\ge 0$.
\item Rejecting valid roots: when solving with $u=e^{x}$, only $u>0$ is acceptable, but every positive $u$ gives a solution.
\item Misapplying Osborn's rule: the sign change applies to \emph{products of two sines}, including hidden ones like $\tan^{2}x=\sin^{2}x/\cos^{2}x$.
\end{itemize}

\textbf{GCE tips.}
\begin{itemize}
\item Quote the exponential definitions first in any proof question; examiners award marks for them.
\item Show the substitution $u=e^{x}$ explicitly when solving equations, and state why negative roots are rejected.
\item In integration, identify the standard form before substituting, and do not forget the constant $C$ and the domain condition.
\item Sketch graphs quickly: $\cosh$ is a U-shape through $(0,1)$, $\sinh$ passes through the origin, $\tanh$ has asymptotes $y=\pm 1$; inverse graphs are reflections in $y=x$.
\end{itemize}
\end{bilan}

\findecours

\end{document}
